% concatenated from 7.22_Hu-Yuan.tex
\documentclass[10pt,reqno]{amsart}
\usepackage{amsfonts}
\usepackage{amsmath}
\usepackage{CJK}
\usepackage{amssymb}
\usepackage{latexsym,bm}
\usepackage[RGB]{xcolor}
\usepackage{mathrsfs}
\usepackage{amsthm}
\usepackage{hyperref}
\usepackage[normalem]{ulem}
%\usepackage{checkref}
\allowdisplaybreaks[4]
%\usepackage{refcheck}
\numberwithin{equation}{section}
\textwidth=15.0cm \textheight=21cm \hoffset=-1.1cm \voffset=-0.3cm
\newtheorem{theorem}{Theorem}[section]
\newtheorem{assumption}[theorem]{Assumption}
\newtheorem{lemma}[theorem]{Lemma}
\newtheorem{remark}{Remark}[section]
\newtheorem{corollary}[theorem]{Corollary}
\newtheorem{proposition}[theorem]{Proposition}
\newcommand{\dif}{\mathrm{d}}
\newcommand{\di}{\mathrm{div}}
\newcommand{\rd}{\rho^d}
\newcommand{\ud}{u^d}
\newcommand{\SUM}[3]{\sum\limits_{{#1}={#2}}^{#3}}
\newcommand{\td}{\theta^d}
\newcommand{\intx}{\int_0^t \int_{\Omega} }
\newcommand{\LIM}[2]{\lim\limits_{{#1}\rightarrow{#2}}}
\newcommand{\blau}[1]{{\textcolor{blue}{#1}}}
\newcommand{\roth}[1]{{\textcolor{red}{#1}}}
\begin{document}
	\title[Spherically symmetric solutions for relaxed CNS]{Global Spherically Symmetric Solutions and Relaxation Limit for the Relaxed Compressible Navier-Stokes Equations}
	\author{Yuxi Hu and Mengran Yuan}
	\thanks{\noindent Yuxi Hu, Department of Mathematics, China University of Mining and Technology, Beijing, 100083 and Key Laboratory of Scientific and Engineering Computing (Ministry of Education), P.R. China, yxhu86@163.com\\
		\indent Mengran Yuan, Department of Mathematics, China University of Mining and Technology, Beijing, 100083, P.R. China,  ymengran2388@126.com
	}
	\thanks{\noindent  }
	\begin{abstract}
		This paper studies an initial boundary value problem for the multidimensional hyperbolized compressible Navier-Stokes equations, in which the classical Newtonian law is replaced by the Maxwell law. We seek spherically symmetric solutions to the studied system in an exterior domain of a ball in $\mathbb R^3$, which are a system possessing a uniform characteristic boundary.
		First, we construct an approximate system featuring a non-characteristic boundary and establish its local well-posedness. Subsequently, by defining a suitable weighted energy functional and carefully handling boundary terms, we derive uniform a priori estimates, enabling the proof of uniform global existence. Leveraging these uniform estimates alongside standard compactness arguments, we establish the global well-posedness of the original system. Additionally, we rigorously justify the global relaxation limit.  \\
		{\bf Keywords}:  Initial boundary value problem; compressible Navier-Stokes equations; Maxwell law; global solution; relaxation limit; \\
		{\bf AMS classification code}: 35A01, 35L50, 35Q30 
	\end{abstract}
	\maketitle
	
	
	\section{Introduction}\label{sec1}
	The compressible isentropic Navier-Stokes equations in the domain \(\mathbb{R}^n \times \mathbb{R}^{+}\) are derived from the fundamental principles of mass conservation and momentum balance. These equations are given by:
	
	\begin{equation} \label{1}
		\begin{cases}
			\partial_t \rho + \operatorname{div}(\rho u) = 0, \\[1mm]
			\partial_t (\rho u) + \operatorname{div}(\rho u \otimes u) + \nabla P = \operatorname{div} S,
		\end{cases}
	\end{equation}
	where \((\rho, u, P, S)\) represent the fluid density, velocity, pressure, and stress tensor, respectively. The pressure \(P\) is assumed to follow the standard \(\gamma\)-law, \(P(\rho) = A \rho^\gamma\) with \(A =1\) without loss of generality. To close the system \eqref{1}, a constitutive equation for the stress tensor \(S\) must be specified. For simple fluids, \(S\) is determined by the Newtonian law:
	
	\begin{equation}\label{h1}
		S = \mu \left( \nabla u + \nabla u^T - \frac{2}{n} \operatorname{div} u I_n \right) + \lambda \operatorname{div} u I_n.
	\end{equation}
	
	The system described by equations \eqref{1} and \eqref{h1} is commonly referred to as the classical compressible isentropic Navier-Stokes equations. However, for complex fluids, such as macromolecular or polymeric fluids, the Newtonian law of viscosity, as given by equation \eqref{h1}, is no longer applicable. In particular, for viscoelastic fluids, Maxwell \cite{maxwell1867iv} introduced a constitutive equation that combines Newton's law of viscosity with Hooke's law of elasticity:
	\begin{equation}\label{h2}
		\tau \dot{S} + S = \mu \left( \nabla u + \nabla u^T - \frac{2}{n} \operatorname{div} u I_n \right) + \lambda \operatorname{div} u I_n,
	\end{equation}
	where \(\dot{S} = \partial_t S + u \cdot \nabla S\) denotes the material derivative. The positive parameter $\tau$ is the relaxation time describing the time lag in the response of the stress tensor to velocity gradient. A fluid obeying equation \eqref{h2} is called Maxwell flow, see \cite{maxwell1867iv}. Replacing the material derivative $\dot S$ by objective derivative, we obtain an objective constitutive equation of compressible Oldroyd-B  type models:
	\begin{align}\label{h3}
		\tau \mathring {S}+S=\mu (\nabla u+\nabla u^T-\frac{2}{n} \mathrm{div} u I_n)+\lambda \operatorname{div} u I_n
	\end{align}
	where $\mathring S:=\dot S+S W(u)-W(u)S-a(S D(u)+D(u)S)$ with $W(u)=\frac{1}{2}(\nabla u-\nabla u^T)$, $D(u)=\frac{1}{2} (\nabla u+\nabla u^T)$ and $-1\le a \le 1$ be a constant. The parameters $a=1, -1, 0$  corresponds to the upper-convective, lower-convective and corotational Maxwell models, see \cite{maxwell1867iv}.
	We should mention that,  even for a simple fluid, the effect of relaxation can not always be neglected, see \cite{pelton2013viscoelastic} with the experiments of high-frequency (20GHZ) vibration of nanoscale mechanical devices immersed in water-glycerol mixtures.  
	
	
	Inspired by the work of Freistühler \cite{freistuhler2020galilei}, which is based on Yong's model \cite{yong2014newtonian}, we investigate the following constitutive equations for the stress tensor \( S = S_1 + S_2 I_n \), where \( S_1 \) and \( S_2 \) satisfy:
	\begin{align}
		\tau_1 \rho \left( \partial_t S_1 + u \cdot \nabla S_1 + S_1 W(u) - W(u) S_1 \right) + S_1 &= \mu \left( \nabla u + \nabla u^T - \frac{2}{n} \mathrm{div} u \, I_n \right), \label{h6} \\
		\tau_2 \rho \left( \partial_t S_2 + u \cdot \nabla S_2 \right) + S_2 &= \lambda \mathrm{div} u. \label{h7}
	\end{align}
	Here, \( \tau_1 \) and \( \tau_2 \) represent the shear and bulk relaxation times, respectively. The objective constitutive equations \eqref{h6}-\eqref{h7} are referred to as the revised Maxwell's law. In its linearized form, this model has been demonstrated to outperform various other models in capturing the dynamic behavior of linear viscoelastic flows, as shown in \cite{chakraborty2015constitutive}. Specifically, for compressible viscoelastic fluids, Chakraborty and Sader \cite{chakraborty2015constitutive} have illustrated that this model provides a comprehensive framework for characterizing the fluid-structure interactions of nanoscale mechanical devices vibrating in simple liquids.
	
	Now, we present some known results for system\eqref{1}, \eqref{h6}-\eqref{h7}. By neglecting the quadratic nonlinear terms  and the factor $\rho$  in \eqref{h6}–\eqref{h7}, Yong \cite{yong2014newtonian} established the local well-posedness of strong solutions and the strong relaxation limit for well-prepared data. These results were subsequently extended to the non-isentropic case by Hu and Racke \cite{hu2017compressible} and to the global relaxation limit with general initial data by Peng \cite{peng2021relaxed}. Additionally, Hu and Wang \cite{HWAML} investigated the finite-time blowup of smooth solutions for certain large initial data. It is important to note that all the aforementioned results focus on the Cauchy problem. In contrast, the initial boundary value problem is more intricate due to its characteristic boundary structure. To the best of our knowledge, the only available results in this context are in the one-dimensional setting, where Hu and Li \cite{HL24} and Hu and Zhao \cite{HZ25} derived global strong solutions for the isentropic and non-isentropic systems, respectively.
	
	In this paper, we are concerned with the initial boundary value problem for the system  \eqref{1}, \eqref{h6}, and \eqref{h7}, within a multidimensional domain that exhibits spherical symmetry. Specifically, we consider the exterior domain of a ball in \(\mathbb{R}^n\) with $n=3$. Denote \(\Omega := \{x \in \mathbb{R}^3 \mid |x| > a\}\) with \(a > 0\). Without loss of generality, we set $a = 1$.
	Our objective is to seek spherically symmetric solutions of the form
	\begin{equation}
		\begin{gathered}
			u(t, x) = v(t, r) \frac{x}{r}, \quad \rho(t, x) = \rho(t, r), \\
			S_1(t, x) = [S^{ij}] = \tilde{S}_1(t, r) \left[\frac{x_i x_j}{r^2} - \frac{1}{3} \delta_{ij}\right], \quad S_2(t, x) = \tilde{S}_2(t, r),
		\end{gathered}
	\end{equation}
	where \(x = (x_1, x_2, x_3) \in \Omega\) and \(r := |x|\). By using this symmetry, we effectively reduce the original system (\ref{1}) to the following set of equations:
	\begin{equation} \label{1.4}
		\begin{cases}
			\partial_t \rho + \partial_r (\rho v) + \dfrac{2}{r} \rho v = 0, \\
			\rho \partial_t v + \rho v \partial_r v + \partial_r P = \dfrac{2}{3} \partial_r \tilde{S}_1 + \dfrac{2}{r} \tilde{S}_1 + \partial_r \tilde{S}_2, \\
			\tau_1 \rho \left(\partial_t \tilde{S}_1 + v \partial_r \tilde{S}_1 \right) + \tilde{S}_1 = 2 \mu \left(\partial_r v - \dfrac{v}{r}\right), \\
			\tau_2 \rho \left(\partial_t \tilde{S}_2 + v \partial_r \tilde{S}_2 \right) + \tilde{S}_2 = \lambda \left(\partial_r v + \dfrac{2}{r} v\right).
		\end{cases}
	\end{equation}
	The system \eqref{1.4} is supplemented with the following initial conditions
	\begin{equation}\label{1.5}
		(\rho(t, r), v(t, r), \tilde{S}_1(t, r), \tilde{S}_2(t, r))|_{t=0} = (\rho_0(r), v_0(r), \tilde{S}_{10}(r), \tilde{S}_{20}(r)), \quad r \in [1, \infty),
	\end{equation}
	and the boundary condition
	\begin{equation}\label{1.6}
		v(t, r)|_{r=1} = 0, \quad t > 0.
	\end{equation} 
	
	Note that when $\tau_1=\tau_2=0$, system (\ref{1.4}) reduces to classical isentropic Navier-Stokes equations with spherical symmetry, for which there are lots of studies on well-posedness theory and large time behavior of both strong and weak solutions,  see \cite{jiang1996global, ding2012global, Cho-Kim}. 
	In particular, for large data that are away from vacuum, Jiang \cite{jiang1996global} proved the global existence and large time behavior of spherically symmetric smooth solutions  in the exterior of a ball in $\mathbb{R}^n$ (\(n = 2\) or 3). When vacuum is taken into account, local and global smooth solutions have also been obtained under certain compatibility conditions, as demonstrated in \cite{ding2012global, Cho-Kim}.  For results concerning weak solutions, we refer to \cite{feireisl2004dynamics,hoff1992spherically,jiang2001spherically,lions1998mathematical,makino1994initial} and references cited therein.
	
	For \(\tau_1 > 0\) and \(\tau_2 > 0\), the fundamental change in the system's structure—from a hyperbolic-parabolic coupled system to a purely hyperbolic system—renders the global existence of classical solutions with large initial data highly unlikely. For instance, Hu and Wang \cite{HWAML} demonstrated that smooth solutions to the relaxed isentropic compressible Navier-Stokes system must blow up in finite time for certain large initial data. Moreover, it can be readily demonstrated that the hyperbolic system \eqref{1.4}–\eqref{1.6} features a characteristic boundary, which introduces significant complexity to the well-posedness theory, even in the local context, see \cite{Chen2007,schochet1986compressible}. 
	
	Our objective is to establish the global well-posedness of strong solutions to system \eqref{1.4}–\eqref{1.6} with small initial data. The strategy we employ is as follows: First, we construct an approximate system with a non-characteristic boundary. By verifying that the boundary is maximally nonnegative, we establish the local well-posedness of strong solutions for this approximated system. Next, we derive uniform a priori estimates and demonstrate the existence of a unique global solution for this approximated system. In this part, the estimates for second-order derivative with respect to $x$ play crucial role. To derive these estimates, we leveraged the structure of the system to construct multiple multipliers. Additionally, boundary terms emerging from integration by parts posed another  challenge. By carefully exploiting the special structure of the system and various multiplier energy methods, the boundary term can be controlled.  Finally, by taking the limit and applying standard compactness arguments, we obtain the desired results.
	
	Now, we introduce some notations. \( W^{m,p} = W^{m,p}(\Omega),\ 0 \leq m \leq \infty,\ 1 \leq p \leq \infty \), denotes the usual Sobolev space with norm \( \|\cdot\|_{W^{m,p}} \). \( H^m \) and \( L^p \) stand for \( W^{m,2} \) resp. \( W^{0,p} \). \( D^\alpha = \partial_t^{\alpha_1} \partial_r^{\alpha_2},\ \alpha = (\alpha_1, \alpha_2),\ |\alpha| = \alpha_1 + \alpha_2 \).
	
	
	
	The following assumptions are needed throughout the paper:
	\begin{assumption}\label{Assump1}
		\begin{enumerate}
			\item[(1)] \(\tau_1\) is the same order as \(\tau_2\), i.e., \(\tau_1 = O(\tau_2)\) as \(\tau_2 \to 0\). Without loss of generality, let \(\tau:= \tau_1 = \tau_2\).
			
			\item[(2)] The initial and boundary data satisfy compatibility condition
			\begin{equation}\label{compatibility}
				\partial_t^k v(0, r)\big|_{r=1} = 0, \quad k = 0, 1. 
			\end{equation}
			\item[(3)] The initial data is well-prepared in the following sense:
			\[
			\|r(\tilde{S}_1(0) - 2\mu(\partial_{r}v_0-\frac{v_0}{r}))\|_{H^1}
			= O(\sqrt{\tau}), 
			\quad \| r(\tilde{S}_2(0) - \lambda(\partial_{r}v_0+\frac{2v_0}{r}))\|_{H^1}
			= O(\sqrt{\tau}), \quad \text{as } \tau \to 0. 
			\]
			
			\item[(4)] \(rV_0^k \in H^{2-k}\) for \(k = 0, 1, 2\) uniformly with respect to $\tau$, where
			\[
			V_0^k:=  \partial_t^k \left(\rho - 1, v, \sqrt{\tau} \tilde{S}_1, \sqrt{\tau} \tilde{S}_2\right) (t = 0, \cdot), \quad k = 0, 1
			\]
			\[
			V_0^2:= \tau \partial_t^2 \left(\rho - 1, v, \sqrt{\tau} \tilde{S}_1, \sqrt{\tau} \tilde{S}_2\right) (t = 0, \cdot)
			\]
			%where \(V_0^1\), \(V_0^2\) are defined recursively by equations (1.10).
		\end{enumerate}  
	\end{assumption}
	
	
	The main results of this paper are as follows.
	\begin{theorem}\label{thm1}
		Let Assumption \eqref{Assump1} hold. Then, there exists a small constant $\epsilon_0>0$ such that if
		\begin{equation}
			E_0:= \sum_{k=0}^2\|rV_0^k\|_{H^{2-k}}^2<\epsilon_0,
		\end{equation}
		there exists a globally defined solution $(\rho,v,\tilde{S}_1,\tilde{S}_2)(r,t)$ to system \eqref{1.4}-\eqref{1.6} satisfying
		$$(\rho-1,v,\tilde{S}_1,\tilde{S}_2)(t,r)\in C([0,\infty),H^{2-\delta_0}(\Omega))\cap C^1([0,\infty),H^{1-\delta_0}(\Omega)) $$
		for any $\delta_0>0$ and 
		\begin{equation}\label{1.8}
			\begin{aligned}
				&\sup _{0 \leq t < \infty}\bigg( \sum_{k=0}^1\left\|r\partial_{t}^k(\rho-1, v, \sqrt{\tau} \tilde{S}_1, \sqrt{\tau} \tilde{S}_2)\right\|_{H^{2-k}}^2
				+\tau^2\left\|r\partial_{tt}(\rho-1, v, \sqrt{\tau} \tilde{S}_1, \sqrt{\tau} \tilde{S}_2)\right\|_{L^{2}}^2\bigg)\\
				&+\int_{0}^{\infty}\bigg(\sum_{\alpha=1}^2\left\|r D^\alpha(\rho, v)\right\|_{L^2}^2
				+\sum_{k=0}^1\left\|r\partial_{t}^k(\tilde{S}_1, \tilde{S}_2)\right\|_{H^{2-k}}^2
				+\tau^2\left\|r(\partial_{tt}\tilde{S}_1,\partial_{tt}\tilde{S}_2)\right\|_{L^{2}}^2\bigg)\mathrm{~d} t
				\leq C E_0
			\end{aligned}
		\end{equation}
		where C is a universal constant independent of $\tau$.
	\end{theorem}
	
	\begin{remark} 
		Observe that the norm 
		\(\|V_0^k\|_{H^{2-k}} = O(1)\) as \(\tau \to 0\), which follows from the assumption of well-prepared initial data. To verify this, we derive the following
		\[
		\left\|r\sqrt{\tau}\partial_{t}\tilde{S}_1(t=0,\cdot)\right\|_{H^1} 
		= \left\|r\left(\frac{2\mu(\partial_{r}v-\frac{v}{r})-\tilde{S}_1}{\sqrt{\tau}}-\sqrt{\tau}v\partial_{r}\tilde{S}_1\right)(t=0,\cdot)\right\|_{H^1} 
		= O(1), \quad \text{as } \tau \to 0,
		\]
		and
		\[
		\left\|r\tau^{\frac{3}{2}}\partial_{tt}\tilde{S}_1(t=0,\cdot)\right\|_{L^2} 
		= \left\|r\left(\sqrt{\tau}2\mu(\partial_{r}v-\frac{v}{r})-\sqrt{\tau}\tilde{S}_1-\tau^{\frac{3}{2}}v\partial_{r}\tilde{S}_1\right)_t(t=0,\cdot)\right\|_{L^2} 
		=O(1), \quad \text{as } \tau \to 0.
		\]
		The other terms can be handled similarly.
	\end{remark}
	
	Based on the uniform estimates of solutions, we have the following convergence theorem.
	\begin{theorem}\label{thm2}
		(Global Weak Convergence).
		Let $(\rho^{\tau},v^{\tau},\tilde{S}_1^{\tau},\tilde{S}_2^{\tau})$ denote the family of global solutions constructed in Theorem \ref{thm1}. Then there exist limit functions
		$$
		(\rho^{0},v^{0}) \in L^{\infty}\left(\mathbb{R}^{+} ; H^2(\Omega)\right), \quad(\tilde{S}_1^{0},\tilde{S}_2^{0}) \in L^2\left(\mathbb{R}^{+} ; H^2(\Omega)\right),
		$$
		such that along a subsequence as $\tau\rightarrow 0$:
		\begin{equation}
			\begin{aligned}
				& \left(\rho^{\tau},v^{\tau}\right) \rightharpoonup^*\left(\rho^{0},v^{0}\right) \quad \text { weakly-* in } L^{\infty}\left(\mathbb{R}^{+} ; H^2(\Omega)\right), \\
				& \left(\tilde{S}_1^{\tau},\tilde{S}_2^{\tau}\right) \rightharpoonup\left(\tilde{S}_1^{0},\tilde{S}_2^{0}\right) \quad \text { weakly in } L^2\left(\mathbb{R}^{+} ; H^2(\Omega)\right) .
			\end{aligned}
		\end{equation}
		The limit $\left(\rho^{0},v^{0}\right)$ constitutes a classical solution to the  spherical symmetric  isentropic compressible Navier-Stokes system  with boundary condition (\ref{1.6}) and initial data $\left(\rho^{0},v^{0}\right)(t=0,x)=(\rho_0,v_0)(x)$. Moreover,
		$$\tilde{S}_1^0=2\mu(\partial_{r}v^0-\frac{v^0}{r}), \quad \tilde{S}_2^0=\lambda(\partial_{r}v^0+\frac{2v^0}{r}).$$
	\end{theorem}
	
	
	
	The paper is organized as follows. In Section 2, we construct an approximate system which possesses a non-characteristic boundary and prove its local well-posedness. In Section 3, we establish  uniform a priori estimates and  obtain a global solution for this approximated system. Based on the uniform estimates derived in Section 3 and usual compactness argument, we justify the limit and prove our main theorem in the final section.
	
	\section{Approximated system and local  existence}\label{sec2}
	In this section, we construct an approximate system with a non-characteristic boundary and establish the local existence of a unique smooth solution to the system. Following the approach of Schochet \cite{schochet1986compressible}, as adopted in \cite{HL24, HZ25}, we approximate system \eqref{1.4} as follows:
	\begin{equation} \label{9}
		\begin{cases}
			\partial_t \rho^{\epsilon}+\partial_r(\rho^{\epsilon} v^{\epsilon})+\dfrac{2}{r} \rho^{\epsilon} v^{\epsilon}=0, \\
			\rho^{\epsilon}\partial_t v^{\epsilon}+\rho^{\epsilon} v^{\epsilon} \partial_r v^{\epsilon}+\partial_r P=\dfrac{2}{3}\partial_r \tilde{S}_1^{\epsilon} +\dfrac{2}{r} \tilde{S}_1^{\epsilon}+\partial_r \tilde{S}_2^{\epsilon},\\
			\tau   \rho^{\epsilon} \left(\partial_t \tilde{S}_1^{\epsilon}+(v^{\epsilon}-\epsilon)\partial_r\tilde{S}_1^{\epsilon} \right) +\tilde{S}_1^{\epsilon}=2\mu(\partial_r v^{\epsilon}- \dfrac{v^{\epsilon}}{r})
			, \\
			\tau  \rho^{\epsilon} \left(\partial_t \tilde{S}_2^{\epsilon}+(v^{\epsilon}-\epsilon)\partial_r\tilde{S}_2^{\epsilon} \right) +\tilde{S}_2^{\epsilon}=\lambda (\partial_r v^{\epsilon}+\dfrac{2}{r}v^{\epsilon}),
		\end{cases}
	\end{equation}
	with initial condition
	\begin{equation}\label{10}
		(\rho^\epsilon,v^\epsilon,\tilde{S}_1^{\epsilon},\tilde{S}_2^{\epsilon})(0,r)=(\rho_0,v_0,\tilde{S}_{10},\tilde{S}_{20})(r),
	\end{equation}
	and boundary condition
	\begin{equation}\label{11}
		v^\epsilon(t,1)=0.
	\end{equation}
	
	We transform the above system into a first-order system for $V^\epsilon:=(\rho^\epsilon,v^\epsilon,\tilde{S}_1^{\epsilon},\tilde{S}_2^{\epsilon})^{T}$,
	\begin{equation}
		A^0(V^\epsilon)\partial_{t}V^\epsilon+A^1(V^\epsilon)\partial_{r}V^\epsilon+B(V^\epsilon)V^\epsilon=0,
	\end{equation} 
	where\\
	$A^0(V^\epsilon)=\left(\begin{array}{ccccc}
		\dfrac{P^{\prime}(\rho^\epsilon)}{\rho^\epsilon} & 0 & 0 & 0 \\
		0 & \rho^\epsilon & 0 & 0 \\
		0 & 0 & \dfrac{\tau  \rho^\epsilon}{3\mu} & 0 \\
		0 & 0 & 0 & \dfrac{\tau \rho^\epsilon}{\lambda}  
	\end{array}\right)$,
	$
	A^1(V^\epsilon)=\left(\begin{array}{ccccc}
		\dfrac{P^{\prime}(\rho^\epsilon)v}{\rho} & P^{\prime}(\rho^\epsilon) & 0 & 0 \\
		P^{\prime}(\rho^\epsilon) & \rho^\epsilon v^\epsilon & -\dfrac{2}{3} & -1 \\
		0 & -\dfrac{2}{3}  & \dfrac{\tau   \rho^\epsilon(v^\epsilon-\epsilon) }{3\mu} & 0 \\
		0 & -1 & 0 &  \dfrac{\tau \rho^\epsilon(v^\epsilon-\epsilon)}{\lambda}
	\end{array}\right)$,\\
	$B(V^\epsilon)=\left(\begin{array}{ccccc}
		0 & \dfrac{2 P^{\prime}(\rho^\epsilon)}{r} & 0 & 0  \\
		0 & 0 & -\dfrac{2}{r} & 0  \\
		0 & \dfrac{2}{3r}  & \dfrac{1}{3\mu} & 0  \\
		0 & -\dfrac{2}{r} & 0 & \dfrac{1}{\lambda}
	\end{array}\right)$,\\
	with initial condition $V^\epsilon(0,r)=V_0:=(\rho_0,v_0,\tilde{S}_{10},\tilde{S}_{20})(r)$ 
	and boundary condition
	$$MV^\epsilon\big|_{\partial\Omega}=0,~~\text{with} 
	~~M=\left(\begin{array}{ccccc}
		0 & 0 & 0 & 0 \\
		0 & 1 & 0 & 0 \\
		0 & 0 & 0 & 0 \\
		0 & 0 & 0 & 0  
	\end{array}\right).$$\\
	
	Note that $det(A^0)^{-1}A^1(V^\epsilon)\big|_{\partial\Omega}=-P^{\prime}(\rho^\epsilon)\epsilon^2\neq0$
	with $P^{\prime}(\rho^\epsilon)>0$ for any $V^\epsilon\in G:=\{(0,\infty)\times\mathbb{R}^3\}$. So, the boundary condition \eqref{11} is a non-characteristic boundary for the approximated system \eqref{9}. 
	
	Next, we need to prove that \eqref{11} satisfies maximally non-negative property. Namely, the matrix $\left.A^1\left(V^\epsilon\right) \cdot \nu\right|_{\partial \Omega}$ is positive semi-definite on the null space $N$ of $M$, yet not on any space encompassing $N$, where $\nu$ denotes the unit outward normal vector on the boundary. Here $\nu=-1$. Let $\xi=\left(\xi_1, 0, \xi_2,\xi_3\right)^T \in \operatorname{ker} M=\operatorname{span}\left\{(1,0,0,0)^T,(0,0,1,0)^T,(0,0,0,1)^T\right\}$, then 
	$$\xi^T A^1 \cdot \nu\big|_{\partial \Omega} \xi=
	\frac{\tau   \rho^{\epsilon}\epsilon}{3\mu}\xi_2^2+\frac{\tau  \rho^{\epsilon}\epsilon}{\lambda}\xi_3^2\geq0.$$
	On the other hand, the only space having $\operatorname{ker} M$ as a proper subspace is $\mathbb{R}^4$. Thus, we take $q=(1,1,0,0)^T\in \mathbb{R}^4$, and  calculate
	$$q^T A^1 \cdot \nu\big|_{\partial \Omega} q=-2P^{\prime}(\rho^{\epsilon})<0.$$
	Thus, the maximally nonnegative property is satisfied. Therefore, classical results implies local well-posedness theory\cite{schochet1986compressible}.
	
	\begin{theorem}
		Suppose $(\rho_0,v_0,\tilde{S}_{10},\tilde{S}_{20})(r)\in H^2$ satisfying the compatibility condition (\ref{compatibility}) and
		\begin{equation*}
			\min\limits_{r \in[1,\infty)} \rho_0(r)>0.
		\end{equation*}
		Then there exists a unique local solution $(\rho^\epsilon,v^\epsilon,\tilde{S}_1^{\epsilon},\tilde{S}_2^{\epsilon})$ to initial boundary value problem (\ref{9})-(\ref{11})
		on some time interval $[0, T]$ with
		$$
		\begin{array}{r}
			\left(\rho^\epsilon,v^\epsilon,\tilde{S}_1^{\epsilon},\tilde{S}_2^{\epsilon}\right) \in C^0\left([0, T], H^2\right) \cap C^1\left([0, T], H^1\right), \\
			\min \limits_{r \in[1,\infty)} \rho(t,r)>0, \quad \forall \quad t>0 . 
		\end{array}
		$$
	\end{theorem}
	\section{Uniform a priori estimates for the approximated system}\label{sec3}
	
	In this section, we derive uniform a priori estimates with respect to \(\tau\) and \(\epsilon\) for the approximated system \eqref{9}-\eqref{11}, yielding uniform global solutions by usual continuation arguments.
	
	For simplicity, we still denote $(\rho^{\epsilon}, v^{\epsilon}, \tilde{S}_1^{\epsilon},\tilde{S}_2^{\epsilon})$ by $(\rho, v,\tilde{S}_1,\tilde{S}_2)$. 
	First, we define the energy 
	
	\begin{equation}
		E(t):=\sup _{0 \leq s \leq t}\bigg( \sum_{k=0}^1\left\|r\partial_{t}^k(\rho-1, v, \sqrt{\tau  } \tilde{S}_1, \sqrt{\tau } \tilde{S}_2)(s,\cdot)\right\|_{H^{2-k}}^2
		+\tau^2\left\|r\partial_{t}^2(\rho-1, v, \sqrt{\tau  } \tilde{S}_1, \sqrt{\tau } \tilde{S}_2)(s,\cdot)\right\|_{L^{2}}^2\bigg)
	\end{equation}
	and  dissipation
	\begin{equation}
		\mathcal{D}(t):=\sum_{|\alpha|=1}^2\left\|r D^\alpha(\rho, v)\right\|_{L^2}^2
		+\sum_{k=0}^1\left\|r\partial_{t}^k(\tilde{S}_1, \tilde{S}_2)\right\|_{H^{2-k}}^2
		+\tau^2\left\|r(\partial_{tt}\tilde{S}_1,\partial_{tt}\tilde{S}_2)\right\|_{L^{2}}^2
	\end{equation}
	where $D^\alpha=\partial_{t}^{\alpha_1}\partial_{r}^{\alpha_2}$ with $|\alpha|=\alpha_1+\alpha_2$. We are aiming to prove the following proposition.
	\begin{proposition}\label{prop1}
		Let Assumption \ref{Assump1} hold and $(\rho, v, \tilde{S}_1, \tilde{S}_2)\in C^0([0,T],H^2)$ be local solutions to system (\ref{9}). Assume that there exists a small $\delta$ such that $E(t)\leq \delta$ for any $0\le t \le T$. Then, for any $0\le t \le T$, we have
		\begin{equation}
			E(t)+\int_0^t \mathcal{D}(s) \mathrm{d}s\leq C \left(E(0)+E^{\frac{3}{2}}(t)+E^{\frac{1}{2}}(t) \int_0^t \mathcal{D}(s) \mathrm{d}s\right)
		\end{equation}
		where $C$ is a constant independent of $\tau$ and $\epsilon$.
	\end{proposition}
	
	In the following lemmas, we suppose \(\delta\) is small enough such that
	\begin{equation}\label{hu3.4}
		\frac{3}{4} \leq \sup_{(t,r)\in[0,T]\times[1,+\infty)} \rho(t,r) \leq \frac{5}{4}.
	\end{equation}
	
	Furthermore, without loss of generality, we assume $\tau \leq 1$ and $\epsilon \ll 1$. $C$ denotes a universal constant which is independent of \(\tau\) and \(\epsilon\).
	
	\subsection{Lower energy estimates}
	With simplified notation, we have equations
	\begin{equation} \label{3.4}
		\begin{cases}
			\partial_t \rho+\partial_r(\rho v)+\dfrac{2}{r} \rho v=0, \\
			\rho\partial_t v+\rho v \partial_r v+\partial_r P=\dfrac{2}{3}\partial_r \tilde{S}_1 +\dfrac{2}{r} \tilde{S}_1+\partial_r \tilde{S}_2,\\
			\tau  \rho \left(\partial_t \tilde{S}_1+(v-\epsilon)\partial_r\tilde{S}_1 \right) +\tilde{S}_1=2\mu\left(\partial_r v- \dfrac{v}{r}\right)
			, \\
			\tau  \rho \left(\partial_t \tilde{S}_2+(v-\epsilon)\partial_r\tilde{S}_2 \right) +\tilde{S}_2=\lambda (\partial_r v+\dfrac{2}{r}v)
			.
		\end{cases}
	\end{equation}
	\begin{lemma}\label{lem1}
		There exists some constant $C$ such that for any $0\le t \le T$
		\begin{align}
			\|r(\rho-1, v, \sqrt{\tau  } \tilde{S}_1, \sqrt{\tau } \tilde{S}_2)\|_{L^2}^2+\int_0^t\|r(\tilde{S}_1,\tilde{S}_2)\|_{L^2}^2 \mathrm{d}t 
			+\int_{0}^{t}\frac{\tau \epsilon}{8 \mu} \tilde{S}_1^2(t,1)\mathrm{d}t \nonumber\\
			\leq 
			C \left(E(0)+E^{\frac{1}{2}}(t) \int_0^t \mathcal{D}(s) \mathrm{d}s\right).
		\end{align}
	\end{lemma}
	
	\begin{proof}
		Multiplying the equation $(\ref{3.4})_2$ by $r^2 v$, we have
		\begin{equation}\nonumber
			(\frac{1}{2}r^2 \rho v^2)_t+\frac{P}{\rho}(r^2 \partial_t\rho+r^2 v \partial_r\rho)=\frac{2}{3}r^2v\partial_{r}\tilde{S}_1+2rv\tilde{S}_1+r^2v\partial_{r}\tilde{S}_2.
		\end{equation}
		Integrating the above equation and using the mass equation, we have
		\begin{equation}\label{0-1}
			\frac{\mathrm{d}}{\mathrm{d} t} \int_{1}^{+\infty}\left(
			\frac{r^2}{\gamma-1}(\rho^{\gamma}-1-\gamma(\rho-1))+\frac{r^2 \rho}{2} v^2\right)\mathrm{d} r=
			\int_{1}^{+\infty}(\frac{2}{3}r^2v\partial_{r}\tilde{S}_1+2rv\tilde{S}_1
			+r^2v\partial_{r}\tilde{S}_2)\mathrm{d} r.
		\end{equation}
		
		Multiplying the equation $(\ref{3.4})_3$ by $\frac{r^2}{3 \mu}\tilde{S}_1$, and multiply the equation $(\ref{3.4})_4$ by $\frac{r^2}{\lambda} \tilde{S}_2$, then one can get
		\begin{equation}
			\begin{aligned}\label{0-2}
				&\frac{\mathrm{d}}{\mathrm{d} t} \int_{1}^{+\infty}\left(\frac{\tau   r^2 \rho}{6 \mu}   \tilde{S}_1^2+\frac{\tau  r^2 \rho}{2\lambda}  \tilde{S}_2^2\right)+\int_{1}^{+\infty} (\frac{r^2}{3 \mu}  \tilde{S}_1^2+\frac{r^2}{\lambda} \tilde{S}_2^2) \mathrm{d} r\\
				&-\epsilon\int_{1}^{+\infty}\frac{\tau   r^2}{3 \mu}  \rho \left(\frac{1}{2}\tilde{S}_1^2\right)_r\mathrm{d} r -\epsilon\int_{1}^{+\infty}\frac{\tau  r^2}{\lambda}  \rho \left(\frac{1}{2}\tilde{S}_2^2\right)_r\mathrm{d} r\\
				&=\int_{1}^{+\infty}(\frac{2}{3}r^2\partial_{r}v\tilde{S}_1-\frac{2r}{3}v\tilde{S}_1+r^2\partial_{r}v\tilde{S}_2+2rv\tilde{S}_2)\mathrm{d} r.
			\end{aligned}
		\end{equation}
		Combining the equations $(\ref{0-1})$ and $(\ref{0-2})$ and using boundary condition (\ref{1.6}), we can get the energy equality
		\begin{equation}
			\begin{aligned}
				&\frac{\mathrm{d}}{\mathrm{d} t} \int_{1}^{+\infty}\left[\frac{r^2}{\gamma-1}(\rho^{\gamma}-1-\gamma(\rho-1))+\frac{r^2 \rho}{2} v^2+\frac{\tau   r^2 \rho}{6 \mu}   \tilde{S}_1^2+\frac{\tau  r^2 \rho}{2\lambda}  \tilde{S}_2^2\right] \mathrm{d} r\\
				&-\epsilon\int_{1}^{+\infty}\frac{\tau   r^2}{3 \mu}  \rho \left(\frac{1}{2}\tilde{S}_1^2\right)_r\mathrm{d} r
				~~~~ -\epsilon\int_{1}^{+\infty}\frac{\tau  r^2}{\lambda}  \rho \left(\frac{1}{2}\tilde{S}_2^2\right)_r\mathrm{d} r +\int_{1}^{+\infty} (\frac{r^2}{3 \mu}  \tilde{S}_1^2+\frac{r^2}{\lambda} 
				\tilde{S}_2^2) \mathrm{d} r=0.
			\end{aligned}
		\end{equation}
		Besides, we have 
		\begin{equation}\nonumber
			\begin{aligned}
				-\frac{\tau \epsilon}{3 \mu} \int_{1}^{+\infty} r^2\rho \left(\frac{1}{2}\tilde{S}_1^2\right)_r\mathrm{d} r
				&=-\frac{\tau   \epsilon}{3 \mu}r^2\rho \frac{1}{2}\tilde{S}_1^2 \bigg|_1^{+\infty} +\frac{\tau \epsilon}{3 \mu} \int_{1}^{+\infty}(2r\rho+r^2\partial_{r}\rho)\frac{1}{2}\tilde{S}_1^2\mathrm{d} r\\
				&= \frac{\tau \epsilon}{6 \mu}\rho(t,1)\tilde{S}_1^2(t,1)
				+\int_{1}^{+\infty} \frac{\tau \epsilon r}{3 \mu} \rho\tilde{S}_1^2
				-C E^{\frac{1}{2}}(t) \mathcal{D}(t)\\
				&\geq \frac{\tau \epsilon}{6 \mu}\rho(t,1)\tilde{S}_1^2(t,1)
				-C E^{\frac{1}{2}}(t) \mathcal{D}(t)
			\end{aligned}
		\end{equation}
		and
		\begin{equation}\nonumber
			\begin{aligned}
				-\frac{\tau  \epsilon}{\lambda} \int_{1}^{+\infty} r^2\rho \left(\frac{1}{2}\tilde{S}_2^2\right)_r\mathrm{d} r
				&=-\frac{\tau  \epsilon}{\lambda} r^2\rho \frac{1}{2}\tilde{S}_2^2 \bigg|_1^{+\infty}
				+\frac{\tau \epsilon}{\lambda} \int_{1}^{+\infty}(2r\rho+r^2\partial_r\rho)\frac{1}{2}\tilde{S}_2^2\mathrm{d} r\\
				&= \frac{\tau \epsilon}{2\lambda}\rho(t,1)\tilde{S}_2^2(t,1) +\int_{1}^{+\infty} \frac{\tau \epsilon r}{\lambda} \rho \tilde{S}_2^2
				-C E^{\frac{1}{2}}(t) \mathcal{D}(t)\\
				&\geq-C E^{\frac{1}{2}}(t) \mathcal{D}(t)
				.
			\end{aligned}
		\end{equation}
		
		
		
		Combining the above result and \eqref{hu3.4}, we obtain
		\begin{equation}
			\begin{aligned}
				\frac{\mathrm{d}}{\mathrm{d} t} \int_{1}^{+\infty}\left[ \frac{r^2}{\gamma-1}\left(\rho^\gamma-1-\gamma(\rho-1)\right) 
				+\frac{r^2 \rho}{2} v^2+\frac{\tau   r^2 \rho}{6 \mu}  \tilde{S}_1^2+\frac{\tau  r^2 \rho}{2\lambda} \tilde{S}_2^2\right] \mathrm{d} r\\
				+\int_{1}^{+\infty} (\frac{r^2} {3\mu} \tilde{S}_1^2+\frac{r^2}{\lambda} \tilde{S}_2^2) \mathrm{d} r
				+\frac{\tau \epsilon}{8 \mu}\tilde{S}_1^2(t,1)\leq
				C E^{\frac{1}{2}}(t) \mathcal{D}(t).
			\end{aligned}
		\end{equation}
		Using Taylor's expansions, we know
		\begin{equation}
			\rho^\gamma-1-\gamma(\rho-1)=\frac{\gamma(\gamma-1)}{2}\xi^{\gamma-2}(\rho-1)^2,
		\end{equation}
		where $\xi\in(1,\rho)$.
		Using \eqref{hu3.4}, we get,
		\begin{equation}
			\|r(\rho-1, v, \sqrt{\tau  } \tilde{S}_1, \sqrt{\tau } \tilde{S}_2)\|_{L^2}^2+\int_0^t\|r(\tilde{S}_1,\tilde{S}_2)\|_{L^2}^2 \mathrm{d}t 
			+\int_{0}^{t}\frac{\tau \epsilon}{8 \mu}\tilde{S}_1^2(t,1)\mathrm{d}t\leq 
			C \left(E(0)+E^{\frac{1}{2}} \int_0^t \mathcal{D}(s) \mathrm{d}s\right).
		\end{equation}
		So, the proof of Lemma \ref{lem1} is finished.
	\end{proof}
	\subsection{First-order estimates}
	
	\begin{lemma}\label{lem2}
		There exists some constant $C$ such that for any $0\le t \le T$
		\begin{equation}\label{hu3.13}
			\begin{aligned}
				&\int_{1}^{+\infty}  r^2\left((\partial_{t}\rho)^2+(\partial_{t}v)^2
				+\tau  (\partial_{t}\tilde{S}_1)^2+\tau (\partial_{t}\tilde{S}_2)^2\right)\mathrm{d}r
				+ \int_0^t \int_{1}^{+\infty}r^2\left(  (\partial_{t}\tilde{S}_1)^2+ (\partial_{t}\tilde{S}_2)^2\right)\mathrm{d}r\mathrm{d}t\\
				&+\int_{0}^{t}\frac{\tau \epsilon}{8 \mu}(\partial_t\tilde{S}_1)^2(t,1)\mathrm{d}t
				\leq C\left(E(0)+ E^{\frac{1}{2}}(t) \int_{0}^{t}\mathcal{D}(s)\mathrm{~d} t\right).	
			\end{aligned}
		\end{equation}
	\end{lemma}
	\begin{proof}
		Taking derivative with respect to $t$ to the equations (\ref{3.4}), one obtain
		
		\begin{equation}
			\label{15}
			\begin{cases}
				\partial_{tt} \rho + \partial_{tr}(\rho v) + \dfrac{2}{r} \partial_t(\rho v) = 0, \\[6pt]
				
				\rho\partial_{tt} v + \partial_t \rho \cdot \partial_t v + \partial_t \rho \cdot v \partial_r v 
				+ \rho \cdot \partial_t( v \partial_r v) + \partial_{tr} P  = \dfrac{2}{3}\partial_{tr} \tilde{S}_1 + \dfrac{2}{r} \partial_t \tilde{S}_1 + \partial_{tr} \tilde{S}_2, \\[6pt]
				
				\tau \rho \left(\partial_{tt} \tilde{S}_1 + \partial_t ((v-\epsilon) \partial_r\tilde{S}_1) \right) 
				+ \tau \partial_t \rho \left(\partial_t \tilde{S}_1 + (v-\epsilon)\partial_r\tilde{S}_1 \right) 
				+ \partial_t \tilde{S}_1  = 2\mu\left(\partial_{tr} v - \dfrac{1}{r}\partial_t v \right), \\[6pt]
				
				\tau \rho \left(\partial_{tt} \tilde{S}_2 + \partial_t ((v-\epsilon)\partial_r\tilde{S}_2) \right) 
				+ \tau \partial_t \rho \left(\partial_t \tilde{S}_2 + (v-\epsilon)\partial_r\tilde{S}_2 \right) 
				+ \partial_t \tilde{S}_2  = \lambda \left(\partial_{tr} v + \dfrac{2}{r}\partial_t v\right).
			\end{cases}
		\end{equation}
	
		Multiplying the equations $(\ref{15})_{1, 2}$ by $r^2 \frac{P^{\prime}(\rho)}{\rho} \partial_t \rho$ and $r^2 \partial_t v$, respectively, 
		and integrating the result,  we get
		\begin{equation}\label{3.17}
			\begin{aligned}
				&\frac{\mathrm{d}}{\mathrm{d} t} \int_{1}^{+\infty}\left(\frac{r^2P^{\prime}(\rho)}{2\rho} \left(\partial_t \rho\right)^2
				+ \frac{r^2\rho}{2}(\partial_tv)^2\right) \mathrm{d} r\\
				&+\int_{1}^{+\infty}
				\left(r^2 \partial_{tr}(\rho v) \frac{\partial_t P} {\rho}
				+2 r \partial_t(\rho v) \frac{\partial_t P}{\rho}
				+r^2 \partial_{tr}P \partial_t v\right)\mathrm{d} r\\
				&
				\leq\int_{1}^{+\infty}\left(\frac{2r^2}{3} \partial_{tr} \tilde{S}_1 \partial_{t}v + 2r \partial_t \tilde{S}_1  \partial_{t}v+r^2\partial_{tr} \tilde{S}_2 \partial_{t}v\right)\mathrm{d} r
				+C E(t)^{\frac{1}{2}}\mathcal{D}(t),	
			\end{aligned}
		\end{equation}
		where
		\begin{equation}\nonumber
			\begin{aligned}
				&\int_{1}^{+\infty}(r^2  \partial_{tr}(\rho v)  \frac{\partial_t P}{\rho}
				+2 r \partial_t(\rho v)  \frac{\partial_t P}{\rho}
				+r^2 \partial_{tr}P \partial_t v)\mathrm{d} r\\
				=&\int_{1}^{+\infty}(r^2 \partial_{tr}v \partial_t P+2r \partial_{t}v\partial_tP+r^2  \partial_{tr}P  \partial_t v)\mathrm{d} r\\
				&+\int_{1}^{+\infty}(r^2  (v\partial_{tr}\rho+\partial_{r}\rho\partial_{t}v+\partial_{t}\rho\partial_{r}v )  \frac{\partial_t P}{\rho}+2 r v \partial_t\rho  \frac{\partial_t P}{\rho})\mathrm{d} r\\
				\geq&\int_{1}^{+\infty}\partial_r(r^2  \partial_{t}v  \partial_t P)\mathrm{d} r-C E(t)^{\frac{1}{2}}\mathcal{D}(t).
			\end{aligned}
		\end{equation}
		Thus, (\ref{3.17}) can be simplified to
		\begin{equation}\label{3.18}
			\begin{aligned}
				\frac{\mathrm{d}}{\mathrm{d} t}& \int_{1}^{+\infty}\left(\frac{r^2P^{\prime}(\rho)}{2\rho} \left(\partial_t \rho\right)^2
				+ \frac{r^2\rho}{2}(\partial_tv)^2\right) \mathrm{d} r
				+\int_{1}^{+\infty}
				\partial_r(r^2  \partial_{t}v  \partial_t P)\mathrm{d} r\\
				&
				\leq\int_{1}^{+\infty}\left(\frac{2r^2}{3} \partial_{tr} \tilde{S}_1 \partial_{t}v + 2r \partial_t \tilde{S}_1  \partial_{t}v+r^2\partial_{tr} \tilde{S}_2 \partial_{t}v\right)\mathrm{d} r
				+C E(t)^{\frac{1}{2}}\mathcal{D}(t).
			\end{aligned}
		\end{equation}
		
		Multiplying the equation $(\ref{15})_3$ by $\dfrac{r^2}{3 \mu} \partial_t \tilde{S}_1$ and integrating over $[1,+\infty)$ with respect to $r$, we have
		\begin{equation}\label{18}
			\begin{aligned}
				\frac{\mathrm{d}}{\mathrm{d} t}& \int_{1}^{+\infty}\left[\frac{\tau}{3 \mu} \cdot \frac{r^2 \rho}{2}\left(\partial_t \tilde{S}_1\right)^2\right]\mathrm{d} r +\int_{1}^{+\infty}(\frac{r^2 \tau}{3 \mu} \partial_t (\rho v) \partial_t \tilde{S}_1 \partial_r \tilde{S}_1
				-\frac{\tau \epsilon r^2}{3 \mu} \partial_{t}\rho\partial_t\tilde{S}_1\partial_r \tilde{S}_1)\mathrm{d} r\\
				&-\int_{1}^{+\infty}\frac{\tau  \epsilon r^2}{3 \mu}\rho(\frac{1}{2}\partial_t\tilde{S}_1^2)_r\mathrm{d} r
				+\int_{1}^{+\infty}\frac{r^2}{3 \mu}(\partial_t \tilde{S}_1)^2\mathrm{d} r
				=\int_{1}^{+\infty}(\frac{2 r^2}{3} \partial_{tr}v \partial_t \tilde{S}_1-\frac{2 r}{3} \partial_{t}v \partial_t \tilde{S}_1)\mathrm{d} r
				,
			\end{aligned}
		\end{equation}
		where
		\begin{equation}\nonumber
			\int_{1}^{+\infty}(\frac{r^2 \tau}{3 \mu} \partial_t (\rho v) \partial_t \tilde{S}_1 \partial_r \tilde{S}_1
			-\frac{\tau \epsilon r^2}{3 \mu} \partial_{t}\rho\partial_t\tilde{S}_1\partial_r \tilde{S}_1)\mathrm{d} r
			\geq -C E^{\frac{1}{2}}(t) \mathcal{D}(t),
		\end{equation}
		and
		\begin{equation}
			\begin{aligned}
				\nonumber
				-\int_{1}^{+\infty}\frac{\tau  \epsilon r^2}{3 \mu}\rho(\frac{1}{2}\partial_t\tilde{S}_1^2)_r\mathrm{d}r
				&=\frac{\tau \epsilon}{6 \mu}\rho(\partial_t\tilde{S}_1)^2(t,1)
				+\int_{1}^{+\infty}\frac{\tau  \epsilon }{3 \mu}(r^2\rho)_r\frac{1}{2}\partial_t\tilde{S}_1^2\mathrm{d}r\\
				&\geq\frac{\tau \epsilon}{6 \mu}\rho(t,1)(\partial_t\tilde{S}_1)^2(t,1) -C E^{\frac{1}{2}}(t) \mathcal{D}(t).
			\end{aligned}
		\end{equation}
		
		Multiplying the equation $(\ref{15})_4$ by $\dfrac{r^2}{\lambda} \partial_t \tilde{S}_2$ and integrating over $[1,+\infty)$ with respect to $r$, we have
		\begin{equation}
			\begin{aligned}\label{19}
				\frac{\mathrm{d}}{\mathrm{d} t}& \int_{1}^{+\infty}\left[\frac{\tau }{\lambda} \cdot \frac{r^2 \rho}{2}\left(\partial_t \tilde{S}_2\right)^2\right]\mathrm{d}r
				+\int_{1}^{+\infty}(\frac{r^2 \tau }{\lambda}   \partial_t (\rho v) \partial_t \tilde{S}_2 \partial_r \tilde{S}_2
				-\frac{\tau \epsilon r^2}{\lambda}\partial_{t}\rho\partial_t\tilde{S}_2\partial_r \tilde{S}_2)\mathrm{d}r\\
				&-\int_{1}^{+\infty}\frac{\tau \epsilon r^2}{\lambda}\rho(\frac{1}{2}\partial_t\tilde{S}_2^2)_r\mathrm{d}r
				+\int_{1}^{+\infty}\frac{r^2}{\lambda}(\partial_t \tilde{S}_2)^2\mathrm{d}r
				=\int_{1}^{+\infty}(r^2 \partial_{tr}v \partial_t \tilde{S}_2+2r \partial_{t}v \partial_t \tilde{S}_2)\mathrm{d}r
				,
			\end{aligned}
		\end{equation}
		where 
		\begin{equation}\nonumber
			\int_{1}^{+\infty}(\frac{r^2 \tau }{\lambda}   \partial_t (\rho v) \partial_t \tilde{S}_2 \partial_r \tilde{S}_2
			-\frac{\tau \epsilon r^2}{\lambda}\partial_{t}\rho\partial_t\tilde{S}_2\partial_r \tilde{S}_1)\mathrm{d}r
			\geq -C E^{\frac{1}{2}}(t) \mathcal{D}(t),
		\end{equation}
		and
		\begin{equation}
			\begin{aligned}
				\nonumber
				-\int_{1}^{+\infty}\frac{\tau \epsilon r^2}{\lambda}\rho(\frac{1}{2}\partial_t\tilde{S}_2^2)_r\mathrm{d}r
				&=\frac{\tau \epsilon}{2\lambda}\rho(t,1)(\partial_t\tilde{S}_2)^2(t,1)
				+\int_{1}^{+\infty}\frac{\tau \epsilon }{\lambda}(r^2\rho)_r\frac{1}{2}\partial_t\tilde{S}_2^2\mathrm{d}r
				\geq-C E^{\frac{1}{2}}(t) \mathcal{D}(t).
			\end{aligned}
		\end{equation}
		
		Combining the equations (\ref{3.18})--(\ref{19}) and the above estimates, we derive
		\begin{equation}\label{order-1-t}
			\begin{aligned}
				&\frac{\mathrm{d}}{\mathrm{d} t} \int_{1}^{+\infty}\left(\frac{r^2 P^{\prime}(\rho)}{2\rho} (\partial_{t}\rho)^2+\frac{r^2\rho}{2}(\partial_{t}v)^2
				+\frac{\tau  r^2 \rho}{6 \mu} (\partial_{t}\tilde{S}_1)^2 +\frac{\tau r^2 \rho}{2\lambda} (\partial_{t}\tilde{S}_2)^2\right) \mathrm{d}r\\
				&+\int_{1}^{+\infty}\left(\frac{r^2}{3 \mu}(\partial_t \tilde{S}_1)^2+\frac{r^2}{\lambda}(\partial_t \tilde{S}_2)^2\right)\mathrm{d}r
				+\frac{\tau \epsilon}{8 \mu}(\partial_t\tilde{S}_1)^2(t,1)\\
				&\leq \int_{1}^{+\infty}\left(\frac{2}{3} \partial_r(r^2 \partial_t v \partial_t \tilde{S}_1) 
				+\partial_r(r^2 \partial_t v \partial_t \tilde{S}_2)
				-\partial_r(r^2  \partial_{t}v  \partial_t P)\right)\mathrm{d}r
				+C E^{\frac{1}{2}}(t) \mathcal{D}(t)
				\leq C E^{\frac{1}{2}}(t) \mathcal{D}(t)
			\end{aligned}
		\end{equation}
		where we used the compatibility condition (\ref{compatibility}) and \eqref{hu3.4}. 
		
		By integrating the above inequality over $(0, t)$, and noticing that
		\begin{align*}
			\int_{1}^{+\infty}\left(\frac{r^2 P^{\prime}(\rho)}{2\rho} (\partial_{t}\rho)^2+\frac{r^2\rho}{2}(\partial_{t}v)^2
			+\frac{\tau  r^2 \rho}{6 \mu} (\partial_{t}\tilde{S}_1)^2 +\frac{\tau r^2 \rho}{2\lambda} (\partial_{t}\tilde{S}_2)^2\right) (t=0,r)\mathrm{d}r\\
			\leq C \|V_0^1\|_{L^2}^2\leq C E(0),
		\end{align*}
		\eqref{hu3.13} follows immediately and this finish the proof of  Lemma $\ref{lem2}$.
	\end{proof}
	
	\begin{lemma}\label{lem2-1}
		There exists some constant C such that for any $0\leq t \leq T$
		\begin{equation}
			\begin{aligned}
				&\int_{1}^{+\infty}  r^2\left((\partial_{r}\rho)^2+(\partial_{r}v)^2
				+\tau  (\partial_{r}\tilde{S}_1)^2+\tau (\partial_{r}\tilde{S}_2)^2\right)\mathrm{d}r
				+ \int_0^t \int_{1}^{+\infty}r^2\left(  (\partial_{r}\tilde{S}_1)^2+ (\partial_{r}\tilde{S}_2)^2\right)\mathrm{d}r\mathrm{d}t\\
				&+\int_{0}^{t}\frac{\tau\epsilon}{16\mu}(\partial_{r}\tilde{S}_1)^2(t,1)\mathrm{d}t
				\leq C\left(E(0)+ E^{\frac{1}{2}}(t) \int_{0}^{t}\mathcal{D}(s)\mathrm{~d} t\right)
				.	
			\end{aligned}
		\end{equation}
	\end{lemma}
	\begin{proof}
		Taking derivatives with respect to $r$ to the equations (\ref{3.4}), we get
		\begin{equation}\label{22}
			\begin{cases}
				\partial_{tr} \rho+\partial_{rr}(\rho v)+\dfrac{2}{r} \partial_r(\rho v)-\dfrac{2}{r^2} (\rho v)=0, \\
				\rho\partial_{tr} v+\rho v \partial_{rr} v+\partial_{rr} P-\dfrac{2}{3}\partial_{rr}  \tilde{S}_1-\dfrac{2}{r} \partial_r \tilde{S}_1+\dfrac{2}{r^2}\tilde{S}_1-\partial_{rr} \tilde{S}_2=f_2,\\
				\tau   \rho \partial_{tr} \tilde{S}_1+\tau   \rho v \partial_{rr}\tilde{S}_1 
				-2\mu\left(\partial_{rr} v- \dfrac{1}{r}\partial_r v
				+\dfrac{1}{r^2} v \right)+\partial_r \tilde{S}_1-\tau   \epsilon\rho\partial_{rr}\tilde{S}_1= f_3, \\
				\tau  \rho \partial_{tr} \tilde{S}_2+\tau  \rho v \partial_{rr}\tilde{S}_2
				-\lambda \left(\partial_{rr} v+\dfrac{2}{r}\partial_r v-\dfrac{2}{r^2} v\right)+\partial_r \tilde{S}_2-\tau  \epsilon\rho\partial_{rr}\tilde{S}_2=f_4,
			\end{cases}
		\end{equation}
		where
		\begin{equation}
			\begin{aligned}\nonumber
				&f_2:=-\partial_r \rho \partial_t v-\partial_r \rho v \partial_r v-\rho (\partial_{r}v)^2,\\
				&f_3:=-\tau   \rho\partial_r v\partial_r \tilde{S}_1
				-\tau   \partial_r \rho \left(\partial_t \tilde{S}_1+v\partial_r\tilde{S}_1 \right)+\tau   \epsilon\partial_{r}\rho\partial_r\tilde{S}_1,\\
				&f_4:=-\tau  \rho\partial_r v\partial_r \tilde{S}_2
				-\tau  \partial_r \rho \left(\partial_t \tilde{S}_2+v\partial_r\tilde{S}_2 \right)+\tau  \epsilon\partial_{r}\rho\partial_r\tilde{S}_2.
			\end{aligned}
		\end{equation}
		
		Multiplying the equation $(\ref{22})_{1,2}$ by $r^2 \frac{P^{\prime}(\rho)}{\rho} \partial_r \rho$ and $r^2 \partial_{r}v$, respectively,  and integrating the results, we can obtain
		\begin{equation}\label{27}
			\begin{aligned}
				&\frac{\mathrm{d}}{\mathrm{d} t}\int_{1}^{+\infty}(\dfrac{P^{\prime}(\rho)r^2}{2\rho}\left(\partial_r \rho\right)^2+\dfrac{r^2\rho}{2} (\partial_{r}v)^2) \mathrm{d}r
				+\int_{1}^{+\infty}\big[
				r^2\partial_{rr}\left(\rho v\right)\frac{ \partial_r P}{\rho}
				+2r\partial_{r}\left(\rho v\right)\frac{ \partial_r P}{\rho}-2v\partial_{r}P\big]\mathrm{d}r\\
				&+\int_{1}^{+\infty}\left(r^2 \partial_{rr}P \partial_{r}v-\dfrac{2 r^2}{3} \partial_{rr}\tilde{S}_1 \partial_{r}v-2r \partial_{r}\tilde{S}_1 \partial_{r}v+2\tilde{S}_1\partial_{r}v-r^2 \partial_{rr}\tilde{S}_2\partial_{r}v\right)\mathrm{d}r\\
				&\le\int_{1}^{+\infty}r^2f_2 \partial_{r}v\mathrm{d}r+ C E(t)^{\frac{1}{2}}\mathcal{D}(t),
			\end{aligned}
		\end{equation}
		for which we have 
		\begin{equation}\nonumber
			\begin{aligned}
				&\int_{1}^{+\infty}(r^2  \partial_{rr}(\rho v)  \frac{\partial_r P}{\rho}
				+2 r \partial_r(\rho v)  \frac{\partial_r P}{\rho}
				+r^2 \partial_{rr}P \partial_r v)\mathrm{d}r\\
				=&\int_{1}^{+\infty} (r^2\partial_{rr} v\partial_r P+2r\partial_{r} v\partial_r P+r^2\partial_{rr} P\partial_r v)	\mathrm{d}r
				+\int_{1}^{+\infty}(r^2(2\partial_r\rho\partial_rv
				+v\partial_{rr}\rho)\frac{\partial_r P}{\rho}+2r v \partial_r \rho\partial_r P)\mathrm{d}r\\
				\geq&\int_{1}^{+\infty} \partial_{r}(r^2\partial_{r} v  \partial_r P)\mathrm{d}r
				-C E(t)^{\frac{1}{2}}\mathcal{D}(t),
			\end{aligned}
		\end{equation}
		and
		\begin{equation}\nonumber
			\int_{1}^{+\infty}r^2f_2 \partial_{r}v\mathrm{d}r\leq C E(t)^{\frac{1}{2}}\mathcal{D}(t).
		\end{equation}
		Now, in order to handle the term $-2v\partial_{r}P$ in (\ref{27}), we multiply $(\ref{3.4})_2$ by $2v$ and integrate over $[1,+\infty)$ to get
		\begin{equation}\label{3.29}
			\int_{1}^{+\infty}\left(\rho (v^2)_t+2\rho v^2\partial_{r}v+2v\partial_{r}P-\dfrac{4}{3}v\partial_r \tilde{S}_1-\dfrac{4 v }{r}\tilde{S}_1-2v\partial_{r}\tilde{S}_2\right)\mathrm{d}r=0.
		\end{equation}
		
		Thus, adding equation (\ref{3.29}) to (\ref{27}),
		one can get
		\begin{equation}\label{3.30}
			\begin{aligned}
				&\frac{\mathrm{d}}{\mathrm{d} t}\int_{1}^{+\infty}\left[\dfrac{P^{\prime}(\rho)r^2}{2\rho}\left(\partial_r \rho\right)^2+\dfrac{r^2\rho}{2} (\partial_{r}v)^2+\rho v^2\right]\mathrm{d}r
				+ \int_{1}^{+\infty}\partial_{r}(r^2\partial_{r} v  \partial_r P)	\mathrm{d}r\\
				&
				+\int_{1}^{+\infty}\left(-\dfrac{2 r^2}{3} \partial_{rr}\tilde{S}_1 \partial_{r}v-2r \partial_{r}\tilde{S}_1 \partial_{r}v+2\tilde{S}_1\partial_{r}v-r^2 \partial_{rr}\tilde{S}_2\partial_{r}v\right)\mathrm{d}r\\
				&+\int_{1}^{+\infty}\left(-\dfrac{4}{3}v\partial_r \tilde{S}_1-\dfrac{4 v }{r}\tilde{S}_1- 2v\partial_{r}\tilde{S}_2\right)\mathrm{d}r
				\leq C E(t)^{\frac{1}{2}}\mathcal{D}(t).
			\end{aligned}
		\end{equation}
		
		Multiplying the equation $(\ref{22})_3$ by $\dfrac{r^2}{3 \mu} \partial_r \tilde{S}_1$ and integrating over $[1,+\infty)$,  we have
		\begin{equation}\label{29}
			\begin{aligned}
				&\frac{\mathrm{d}}{\mathrm{d} t}\int_{1}^{+\infty}\left[\frac{\tau  }{3 \mu} \cdot \frac{r^2 \rho}{2}(\partial_r \tilde{S}_1)^2\right]\mathrm{d}r
				+\int_{1}^{+\infty}\frac{r^2}{3 \mu}(\partial_r \tilde{S}_1)^2\mathrm{d}r
				-\int_{1}^{+\infty}\frac{\tau  \epsilon }{3 \mu}r^2\rho(\frac{1}{2}(\partial_r \tilde{S}_1)^2)_r\mathrm{d}r\\
				&-\int_{1}^{+\infty}\bigg(\frac{2r^2}{3}\partial_{rr}v\partial_r \tilde{S}_1 -\frac{2r}{3}\partial_{r}v\partial_r \tilde{S}_1
				+\frac{2}{3}v\partial_r \tilde{S}_1\bigg)\mathrm{d}r
				=\int_{1}^{+\infty}\frac{r^2}{3\mu}f_3\partial_r \tilde{S}_1\mathrm{d}r ,
			\end{aligned}
		\end{equation}
		where 
		\begin{equation}
			\begin{aligned}
				\nonumber
				-\int_{1}^{+\infty}\frac{\tau \epsilon }{3\mu}r^2\rho(\frac{1}{2}(\partial_r\tilde{S}_1)^2)_r\mathrm{d}r
				&=\frac{\tau \epsilon }{6\mu}\rho (\partial_r\tilde{S}_1)^2(t,1)
				+\int_{1}^{+\infty}\frac{\tau \epsilon }{3\mu}(2r\rho+r^2\partial_{r}\rho)\frac{1}{2}(\partial_r\tilde{S}_1)^2\mathrm{d}r\\
				%&= \frac{\tau \epsilon }{6\mu}\rho(t,1)(\partial_r\tilde{S}_1)^2(t,1)+ \int_{1}^{+\infty}\frac{\tau\epsilon r}{3\mu} \rho (\partial_r\tilde{S}_1)^2\mathrm{d}r
				%+\int_{1}^{+\infty}\frac{\tau\epsilon r^2} {6\mu} \partial_{r}\rho(\partial_r\tilde{S}_1)^2\mathrm{d}r\\
				%&\geq\frac{\tau \epsilon }{6\mu}\rho(t,1)(\partial_r\tilde{S}_1)^2(t,1)
				%+\int_{1}^{+\infty}\frac{\tau\epsilon r^2} {6\mu} \partial_{r}\rho(\partial_r\tilde{S}_1)^2\mathrm{d}r\\
				&
				\geq \frac{\tau \epsilon }{6\mu}\rho(t,1)(\partial_r\tilde{S}_1)^2(t,1)-C E^{\frac{1}{2}}(t) \mathcal{D}(t),
			\end{aligned}
		\end{equation}
		and
		\begin{equation}\nonumber
			\int_{1}^{+\infty}\frac{r^2}{3\mu}f_3\partial_r \tilde{S}_1\mathrm{d}r\leq C E^{\frac{1}{2}}(t) \mathcal{D}(t).
		\end{equation}
		Similarly,  multiplying the equation $(\ref{22})_4$ by $\dfrac{r^2}{\lambda} \partial_r \tilde{S}_2$,  we get
		\begin{equation}\label{30}
			\begin{aligned}
				&\frac{\mathrm{d}}{\mathrm{d} t}\int_{1}^{+\infty}\left[\frac{\tau }{\lambda} \cdot \frac{r^2 \rho}{2}(\partial_r \tilde{S}_2)^2\right]\mathrm{d}r
				+\int_{1}^{+\infty}\frac{r^2}{\lambda}(\partial_r \tilde{S}_2)^2\mathrm{d}r
				-\int_{1}^{+\infty}\frac{\tau \epsilon }{\lambda}r^2\rho(\frac{1}{2}(\partial_r \tilde{S}_2)^2)_r\mathrm{d}r\\
				&-\int_{1}^{+\infty}\bigg(r^2\partial_{rr}v\partial_r \tilde{S}_2+2r\partial_{r}v\partial_r \tilde{S}_2
				-2v\partial_r \tilde{S}_2\bigg)\mathrm{d}r
				=\int_{1}^{+\infty}\frac{r^2}{\lambda}f_4\partial_r \tilde{S}_2\mathrm{d}r,
			\end{aligned}
		\end{equation}
		where
		\begin{equation}
			\begin{aligned}
				\nonumber
				-\int_{1}^{+\infty}\frac{\tau \epsilon }{\lambda}r^2\rho(\frac{1}{2}(\partial_r\tilde{S}_2)^2)_r\mathrm{d}r
				&=\frac{\tau \epsilon }{2\lambda}\rho (\partial_r\tilde{S}_2)^2(t,1)
				+\int_{1}^{+\infty}\frac{\tau \epsilon }{\lambda}(2r\rho+r^2\partial_{r}\rho)\frac{1}{2}(\partial_r\tilde{S}_2)^2\mathrm{d}r\\
				&\geq-C E^{\frac{1}{2}}(t) \mathcal{D}(t),
			\end{aligned}
		\end{equation}
		and
		\begin{equation}\nonumber
			\int_{1}^{+\infty}\frac{r^2}{\lambda}f_4\partial_r \tilde{S}_2\mathrm{d}r\leq C E^{\frac{1}{2}}(t) \mathcal{D}(t).
		\end{equation}
		
		Combining the equations (\ref{3.30})--(\ref{30}) and the above estimates, we derive that
		\begin{equation}\label{3.33}
			\begin{aligned}
				&\frac{\mathrm{d}}{\mathrm{d} t}\int_{1}^{+\infty}(\dfrac{P^{\prime}(\rho)r^2}{2\rho}\left(\partial_r \rho\right)^2+\dfrac{r^2\rho}{2} (\partial_{r}v)^2+\rho v^2
				+ \frac{\tau r^2\rho}{6\mu}(\partial_r \tilde{S}_1)^2
				+ \frac{\tau r^2\rho}{2\lambda}(\partial_r \tilde{S}_2)^2) \mathrm{d}r\\
				&
				+ \int_{1}^{+\infty}(\partial_{r}(r^2\partial_{r} v  \partial_r P)	-\dfrac{2}{3}\partial_{r}(r^2\partial_{r}v\partial_{r}\tilde{S}_1)-\partial_{r}(r^2\partial_{r}v\partial_{r}\tilde{S}_2))\mathrm{d}r\\
				&+\int_{1}^{+\infty}(\frac{r^2}{3\mu}(\partial_r \tilde{S}_1)^2+\frac{r^2}{\lambda}(\partial_r \tilde{S}_2)^2)\mathrm{d}r
				+\frac{\tau\epsilon}{6\mu}\rho(t,1)(\partial_{r}\tilde{S}_1)^2(t,1)
				\\
				&+\int_{1}^{+\infty}\left(2\tilde{S}_1\partial_{r}v-2v\partial_r \tilde{S}_1-\dfrac{4 v }{r}\tilde{S}_1\right)\mathrm{d}r\\
				&\leq C E(t)^{\frac{1}{2}}\mathcal{D}(t).
			\end{aligned}
		\end{equation}
		To address the last term on the left-hand side of equation (\ref{3.33}), we
		multiply $(\ref{3.4})_3$ by $\frac{2}{\mu}\tilde{S}_1$ and add the result to equation (\ref{3.33}) to get
		\begin{equation}\label{3.34}
			\begin{aligned}
				&\frac{\mathrm{d}}{\mathrm{d} t}\int_{1}^{+\infty}(\dfrac{P^{\prime}(\rho)r^2}{2\rho}\left(\partial_r \rho\right)^2+\dfrac{r^2\rho}{2} (\partial_{r}v)^2+\rho v^2
				+ \frac{\tau r^2\rho}{6\mu}(\partial_r \tilde{S}_1)^2
				+ \frac{\tau r^2\rho}{2\lambda}(\partial_r \tilde{S}_2)^2) \mathrm{d}r\\
				&+\int_{1}^{+\infty}(\frac{r^2}{3\mu}(\partial_r \tilde{S}_1)^2+\frac{r^2}{\lambda}(\partial_r \tilde{S}_2)^2)\mathrm{d}r
				+\frac{\tau\epsilon}{6\mu}\rho(t,1)(\partial_{r}\tilde{S}_1)^2 (t,1)
				\\
				&
				+ \int_{1}^{+\infty}\left(\partial_{r}(r^2\partial_{r} v  \partial_r P)	-\dfrac{2}{3}\partial_{r}(r^2\partial_{r}v\partial_{r}\tilde{S}_1)-\partial_{r}(r^2\partial_{r}v\partial_{r}\tilde{S}_2)\right)\mathrm{d}r\\
				&+\int_{1}^{+\infty}\left(2\tilde{S}_1\partial_{r}v-2v\partial_r \tilde{S}_1-\dfrac{4 v }{r}\tilde{S}_1\right)\mathrm{d}r\\
				&+\int_{1}^{+\infty}\left[\left(\tau   \rho \left(\partial_t \tilde{S}_1+(v-\epsilon)\partial_r\tilde{S}_1 \right) +\tilde{S}_1-2\mu\left(\partial_r v- \dfrac{v}{r}\right)\right) \frac{2}{\mu}\tilde{S}_1\right] \mathrm{d}r
				\leq C E(t)^{\frac{1}{2}}\mathcal{D}(t),
			\end{aligned}
		\end{equation}
		where
		\begin{equation}\nonumber
			-\int_{1}^{+\infty}\frac{2\tau\epsilon}{\mu}\rho\tilde{S}_1
			\partial_r\tilde{S}_1\mathrm{d}r
			=\frac{\tau\epsilon}{\mu}\rho(t,1)\tilde{S}_1^2(t,1)+\int_{1}^{+\infty}\frac{\tau\epsilon}{\mu}\partial_r\rho\tilde{S}_1^2
			\mathrm{d}r
			\geq %\frac{\tau\epsilon}{\mu}\rho(t,1)\tilde{S}_1^2(t,1)
			-C E(t)^{\frac{1}{2}}\mathcal{D}(t).
		\end{equation}
		Integrating the equation (\ref{3.34}) over $(0,t)$ and using \eqref{hu3.4}, we have
		\begin{equation}\label{34}
			\begin{aligned}
				&\int_{1}^{+\infty}\bigg(\dfrac{r^2 P^{\prime}(\rho)}{2\rho} \left(\partial_r \rho\right)^2 +\dfrac{1}{2}\rho\left(r^2 (\partial_rv)^2+2 v^2\right)+\dfrac{\tau  }{4 \mu} \rho \left[\dfrac{2}{3} r^2(\partial_r\tilde{S}_1)^2+4 \tilde{S}_1^2\right]
				+\dfrac{\tau }{2\lambda}\rho r^2 (\partial_r \tilde{S}_2)^2\bigg)\mathrm{d}r\\
				&+\int_{0}^{t}\int_{1}^{+\infty}\bigg(\dfrac{r^2}{3\mu} \left(\partial_r\tilde{S}_1\right)^2+\dfrac{2}{\mu}\tilde{S}_1^2
				+\dfrac{r^2}{\lambda}\left(\partial_r\tilde{S}_2\right)^2
				\bigg)\mathrm{d}r\mathrm{d}t
				+\int_{0}^{t}\frac{3\tau\epsilon}{8\mu}(\partial_{r}\tilde{S}_1)^2 (t,1)\mathrm{d}t\\
				&\leq \int_{0}^{t} r^2\partial_{r}v\left(\partial_{r}P-\dfrac{2}{3}\partial_r\tilde{S}_1-\partial_r\tilde{S}_2\right)(t,1)\mathrm{d}t
				-\int_{0}^{t}2v\tilde{S}_1(t,1)\mathrm{d}t
				+C\left(E(0)+ E^{\frac{1}{2}}(t)\int_{0}^{t} \mathcal{D}(s)\mathrm{d}s\right)
				.
			\end{aligned}
		\end{equation}
		Using the boundary condition (\ref{1.6}), we know $-\int_{0}^{t}2v\tilde{S}_1(t,1)\mathrm{d}t=0$. For the first term on the right hand side of equation \eqref{34}, we have
		\begin{equation}
			\begin{aligned}
				\left[r^2\partial_r v \left( \partial_r P - \dfrac{2}{3}\partial_r \tilde{S}_1 - \partial_r \tilde{S}_2 \right)\right](t,1)
				&= r^2\partial_r v (\partial_r P - \dfrac{2}{3}\partial_r \tilde{S}_1 - \dfrac{2}{r}\tilde{S}_1 - \partial_r \tilde{S}_2 )(t,1) + (2r\tilde{S}_1 \partial_r v) (t,1) \\
				&= -[r^2\partial_r v\left( \rho\partial_t v + \rho v\partial_r v \right)](t,1)+ (2r\tilde{S}_1 \partial_r v) (t,1) \\
				&=  2(\tilde{S}_1 \partial_r v) (t,1),
			\end{aligned}
		\end{equation}
		where we used the momentum equation $(\ref{3.4})_2$ and boundary condition (\ref{1.6}).
		By using the equation $(\ref{3.4})_3$ and $\epsilon-Young$ inequality, we have 
		\begin{equation}\nonumber
			\begin{aligned}
				\int_{0}^{t}2(\tilde{S}_1 \partial_r v )(t,1) \mathrm{d}t
				=&\int_{0}^{t}\left[2\tilde{S}_1\left(\dfrac{\tau  \rho}{2\mu}(\partial_t\tilde{S}_1+(v-\epsilon)\partial_r\tilde{S}_1)+\dfrac{1}{2\mu}\tilde{S}_1+\dfrac{v}{r}\right)\right](t,1)\mathrm{d}t\\
				=&\int_{0}^{t}\left[\dfrac{\tau \rho}{\mu}(\dfrac{1}{2}\tilde{S}_1^2)_t(t,1)
				-\dfrac{\epsilon\tau }{\mu} \rho \tilde{S}_1\partial_r\tilde{S}_1(t,1)
				+\dfrac{1}{\mu}\tilde{S}_1^2(t,1)\right]\mathrm{d}t\\
				\leq& \dfrac{\tau}{2\mu} \rho\tilde{S}_1^2(t,1)\big|_0^t
				-\int_{0}^{t}\frac{\tau}{2\mu}\partial_{t}\rho\tilde{S}_1^2(t,1)\mathrm{d}t
				+\int_{0}^{t}\dfrac{1}{\mu}\tilde{S}_1^2(t,1)\mathrm{d}t\\
				&+\int_{0}^{t}\dfrac{\epsilon\tau }{\mu}\rho(t,1)\left(\eta(\partial_{r}\tilde{S}_1)^2(t,1)+C(\eta)\tilde{S}_1^2(t,1)\right)\mathrm{d}t
			\end{aligned}
		\end{equation}
		where
		$$
		\begin{aligned}
			-\int_{0}^{t}\frac{\tau}{2\mu}\partial_{t}\rho\tilde{S}_1^2(t,1)\mathrm{d}t
			&\leq
			\int_{0}^{t}\frac{\tau}{2\mu}\|\partial_{t}\rho\tilde{S}_1^2\|_{L^{\infty}}\mathrm{d}t
			\leq \int_{0}^{t}\frac{\tau}{2\mu}\|\partial_{t}\rho\|_{H^{1}}\|\tilde{S}_1\|_{H^{1}}^2\mathrm{d}t\\
			&\leq C\left(E(0)+ E^{\frac{1}{2}}(t)\int_{0}^{t} \mathcal{D}(s)\mathrm{d}s\right),
		\end{aligned}
		$$
		and
		\begin{equation}
			\begin{aligned}\nonumber
				&\int_{0}^{t}\dfrac{1}{\mu} \tilde{S}_1^2(t,1)\mathrm{d}t
				\leq \int_{0}^{t}(\dfrac{1}{\mu}\|\tilde{S}_1\|_{L^\infty}^2)\mathrm{d}t
				\leq \int_{0}^{t}\dfrac{1}{\mu}\|\partial_{r}\tilde{S}_1\|_{L^2}\|\tilde{S}_1\|_{L^2}\mathrm{d}t\\
				&\leq \int_{0}^{t}\dfrac{1}{\mu} (\eta\|\partial_{r}\tilde{S}_1\|_{L^2}^2+C(\eta)\|\tilde{S}_1\|_{L^2}^2)\mathrm{d}t
				\leq \int_{0}^{t}\dfrac{\eta}{\mu}\|\partial_{r}\tilde{S}_1\|_{L^2}^2\mathrm{d}t+C\left(E(0)+E^{\frac{1}{2}}(t) \int_{0}^t \mathcal{D}(s)\mathrm{d}s\right).
			\end{aligned}
		\end{equation}
		Here and after, $\eta$ is any positive constant to be chosen later.
		
		Combining the above results, Lemma \ref{lem1} and \eqref{hu3.4}, we derive
		\begin{equation*}
			\begin{aligned}
				&\int_{1}^{+\infty}\bigg(\dfrac{r^2 P^{\prime}(\rho)}{2\rho} \left(\partial_r \rho\right)^2 +\dfrac{\rho}{2}\left[r^2 (\partial_rv)^2+2 v^2\right]+\dfrac{\tau\rho r^2}{6\mu}   \left(\partial_r\tilde{S}_1\right)^2
				+\dfrac{\tau\rho}{\mu}  \tilde{S}_1^2
				+\dfrac{\tau \rho r^2}{2\lambda}  \left(\partial_r \tilde{S}_2\right)^2
				\bigg)\mathrm{d}r\\
				&+\int_0^t\int_{1}^{+\infty}\left(\dfrac{r^2}{4\mu} \left(\partial_r\tilde{S}_1\right)^2+\dfrac{2}{\mu}\tilde{S}_1^2
				+\dfrac{r^2}{\lambda}\left(\partial_r\tilde{S}_2\right)^2 \right)\mathrm{d}r\mathrm{d}t
				+\int_{0}^{t}\frac{\tau\epsilon}{16\mu}(\partial_{r}\tilde{S}_1)^2 (t,1)\mathrm{d}t\\
				&\leq \dfrac{\tau}{2\mu} \rho(t,1)\tilde{S}_1^2(t,1)
				+C\left(E(0)+ E^{\frac{1}{2}}(t)\int_{0}^{t} \mathcal{D}(s)\mathrm{d}s\right)
				,
			\end{aligned}
		\end{equation*}
		and, by applying the {\it Gagliardo-Nirenberg} inequality and {\it Young} inequality, we have 
		$$\begin{aligned}
			\dfrac{\tau}{2\mu} \rho(t,1)\tilde{S}_1^2(t,1)
			&\leq\dfrac{\tau}{\mu} r^2 \|\tilde{S}_1\|_{L^\infty}^2
			\leq \dfrac{\tau}{\mu}r^2\|\partial_{r}\tilde{S}_1\|_{L^2}\|\tilde{S}_1\|_{L^2}
			\leq \dfrac{\tau}{\mu}r^2 (\eta\|\partial_{r}\tilde{S}_1\|_{L^2}^2+C(\eta)\|\tilde{S}_1\|_{L^2}^2)\\
			&
			\leq \dfrac{\tau\eta}{\mu}r^2\|\partial_{r}\tilde{S}_1\|_{L^2}^2+
			C\left(E(0)+ E^{\frac{1}{2}}(t)\int_{0}^{t} \mathcal{D}(s)\mathrm{d}s\right).
		\end{aligned}$$
		By choosing $\eta<\frac{1}{12}$, we get the Lemma $\ref{lem2-1}$ immediately.
		%The estimates $(\ref{order-1-t})$ and $(\ref{order-1-r})$ prove Lemma $\ref{lem2}$.
	\end{proof}
	The next lemma show the dissipative estimates of $D(\rho, v)$.
	\begin{lemma}\label{lem3.4}
		There exists a constant $C$ such that for any $0\le t \le T$ 
		\begin{equation}
			\begin{aligned}
				\int_0^t\int_{1}^{+\infty}(v^2+r^2|D(\rho, v)|^2)\mathrm{d}r\mathrm{d}t
				\leq C\left(E(0)+E^{\frac{1}{2}}(t)\int_{0}^{t} \mathcal{D}(s)\mathrm{d}s\right).
			\end{aligned}
		\end{equation}
	\end{lemma}
	\begin{proof}
		By manipulating $\dfrac{\lambda}{\mu}(\ref{3.4})_{3}+(\ref{3.4})_{4}$
		and applying Lemma \ref{lem2}, we derive
		
		\begin{equation}\label{dis1}
			\int_0^t\int_{1}^{+\infty}r^2(\partial_{r}v)^2
			\mathrm{d}r\mathrm{d}t
			\leq C(E(0)+E^{\frac{1}{2}}(t)\int_{0}^{t} \mathcal{D}(s)\mathrm{d}s).
		\end{equation}
		Similarly, by computing $-\dfrac{\lambda}{2\mu}(\ref{3.4})_{3}+(\ref{3.4})_{4}$, one can get
		\begin{equation}
			\int_0^t\int_{1}^{+\infty}v^2
			\mathrm{d}r\mathrm{d}t
			\leq C(E(0)+E^{\frac{1}{2}}(t)\int_{0}^{t} \mathcal{D}(s)\mathrm{d}s).
		\end{equation}
		Using mass  equation $(\ref{3.4})_{1}$, we get immediately
		\begin{equation}\label{dis2}
			\int_0^t\int_{1}^{+\infty}r^2(\partial_{t}\rho)^2
			\mathrm{d}r\mathrm{d}t
			\leq C(E(0)+E^{\frac{1}{2}}(t)\int_{0}^{t} \mathcal{D}(s)\mathrm{d}s).
		\end{equation}
		On the other hand, multiplying the equation
		$(\ref{3.4})_{2}$ by $r^2\partial_{t}v$ yields
		\begin{equation}\nonumber
			\begin{aligned}
				&\int_0^t\int_{1}^{+\infty}r^2(\partial_{t}v)^2\mathrm{d}r\mathrm{d}t\\
				=&\int_0^t\int_{1}^{+\infty}(-r^2\partial_{r}P\partial_{t}v+\frac{2r^2}{3}\partial_r\tilde{S}_1\partial_{t}v+2r\tilde{S}_1\partial_{t}v
				+r^2\partial_r\tilde{S}_2\partial_{t}v
				-r^2\rho v \partial_{r}v\partial_{t}v)\mathrm{d}r\mathrm{d}t\\
				=&-\int_0^t\frac{\mathrm{d}}{\mathrm{d}t}\int_{1}^{+\infty}r^2\partial_{r}Pv\mathrm{d}r\mathrm{d}t
				-\int_0^t\int_{1}^{+\infty}r^2\partial_{t}P\partial_{r}v\mathrm{d}r\mathrm{d}t
				-\int_0^t\int_{1}^{+\infty}2r\partial_{t}Pv\mathrm{d}r\mathrm{d}t
				\\
				&+\int_0^t\int_{1}^{+\infty}\frac{2r^2}{3}\partial_r\tilde{S}_1\partial_{t}v\mathrm{d}r\mathrm{d}t
				+\int_0^t\int_{1}^{+\infty}2r\tilde{S}_1\partial_{t}v\mathrm{d}r\mathrm{d}t\\
				&+\int_0^t\int_{1}^{+\infty}r^2\partial_r\tilde{S}_2\partial_{t}v\mathrm{d}r\mathrm{d}t
				-\int_0^t\int_{1}^{+\infty}r^2\rho v \partial_{r}v\partial_{t}v\mathrm{d}r\mathrm{d}t\\
				\leq&\frac{1}{2}\int_0^t\int_{1}^{+\infty}r^2(\partial_{t}v)^2\mathrm{d}r\mathrm{d}t
				+C\int_{1}^{+\infty}r^2((\partial_{r}\rho)^2+v^2)\mathrm{d}r
				+C(E(0)+E^{\frac{1}{2}}(t)\int_{0}^{t} \mathcal{D}(s)\mathrm{d}s)\\
				&+C\int_0^t\int_{1}^{+\infty}\left(r^2((\partial_{t}\rho)^2
				+(\partial_{r}v)^2+(\partial_r\tilde{S}_1)^2 +\tilde{S}_1^2+(\partial_r\tilde{S}_2)^2)
				+v^2\right)
				\mathrm{d}r\mathrm{d}t
			\end{aligned}
		\end{equation}
		which implies
		\begin{equation}\label{dis3}
			\int_0^t\int_{1}^{+\infty}r^2(\partial_{t}v)^2
			\mathrm{d}r\mathrm{d}t
			\leq C(E(0)+E^{\frac{1}{2}}(t)\int_{0}^{t} \mathcal{D}(s)\mathrm{d}s).
		\end{equation}
		
		By equation $(\ref{3.4})_2$, we also get
		\begin{equation}\label{dis4}
			\int_0^t\int_{1}^{+\infty}r^2(\partial_{r}\rho)^2
			\mathrm{d}r\mathrm{d}t
			\leq C(E(0)+E^{\frac{1}{2}}(t)\int_{0}^{t} \mathcal{D}(s)\mathrm{d}s).
		\end{equation}
		
		Thus, from the equations (\ref{dis1})-(\ref{dis4}), we get the desired result.
	\end{proof}
	
	Combining Lemmas \ref{lem1}-\ref{lem3.4},  we get the following lemma.
	\begin{lemma}\label{lem3.5}
		There exists some constant C such that for any $0\le t \le T$
		\begin{equation}
			\begin{aligned}
				&\sum_{|\alpha|=0}^1\left\|rD^\alpha(\rho-1, v, \sqrt{\tau} \tilde{S}_1, \sqrt{\tau} \tilde{S}_2)\right\|_{L^{2}}^2
				+\int_{0}^{t}\bigg(\left\|r D(\rho, v)\right\|_{L^2}^2
				+\sum_{|\alpha|=0}^1\left\|r D^\alpha(\tilde{S}_1, \tilde{S}_2)\right\|_{L^{2}}^2+v^2\bigg)\mathrm{~d} t
				\\&
				\leq C \left(E(0)+E^{\frac{1}{2}}(t) \int_{0}^{t}\mathcal{D}(s)\mathrm{d}s\right).
			\end{aligned}
		\end{equation}
	\end{lemma}
	\subsection{Second-order estimates}
	\begin{lemma}\label{lem3}
		There exists some constant C such that for any $0\le t \le T$
		\begin{equation}
			\begin{aligned}
				&\int_{1}^{+\infty} \tau^2r^2 \left((\partial_{tt}\rho)^2+(\partial_{tt}v)^2+\tau(\partial_{tt}\tilde{S}_1)^2+\tau(\partial_{tt}\tilde{S}_2)^2\right) \mathrm{d} r \\
				&+\int_0^t \int_{1}^{+\infty}\tau^2r^2\left((\partial_{tt}\tilde{S}_1)^2+(\partial_{tt}\tilde{S}_2)^2\right) \mathrm{d} r \mathrm{d} t 
				\leq C\left(E(0)+ E^{\frac{1}{2}}(t) \int_{0}^t \mathcal{D}(s)\mathrm{d}s\right).
			\end{aligned}
		\end{equation}
	\end{lemma}
	\begin{proof}
		Taking derivative with respect to $t$ twice to equation $(\ref{3.4})$ , one get
		
		\begin{equation} \label{38}
			\begin{cases}
				\partial_{ttt} \rho+\partial_{ttr}(\rho v)+\dfrac{2}{r}\partial_{tt}(\rho v)=0, \\
				\rho\partial_{ttt} v+2\partial_{t}\rho \partial_{tt}v+ \partial_{tt}\rho \partial_{t}v+\partial_{tt}\rho (v \partial_r v)+2\partial_{t}\rho\partial_{t}(v\partial_{r}v)+\rho\partial_{tt}(v\partial_{r}v)\\
				\qquad\qquad\qquad\qquad\qquad\qquad\qquad\qquad\qquad\qquad+\partial_{ttr} P=\dfrac{2}{3}\partial_{ttr}  \tilde{S}_1+\dfrac{2}{r} \partial_{tt}\tilde{S}_1+\partial_{ttr} \tilde{S}_2,\\
				\tau \rho \left(\partial_{ttt} \tilde{S}_1+\partial_{tt} ((v-\epsilon)\partial_r\tilde{S}_1) \right)
				+2\tau \partial_t \rho \left(\partial_{tt} \tilde{S}_1
				+\partial_{t} ((v-\epsilon)\partial_r\tilde{S}_1) \right) \\
				\qquad\qquad\qquad\qquad\qquad+\tau \partial_{tt} \rho \left(\partial_t \tilde{S}_1+(v-\epsilon)\partial_r\tilde{S}_1 \right) +\partial_{tt} \tilde{S}_1=2\mu\left(\partial_{ttr} v- \dfrac{1}{r}\partial_{tt} v \right), \\
				\tau  \rho \left(\partial_{ttt} \tilde{S}_2+\partial_{tt} ((v-\epsilon)\partial_r\tilde{S}_2) \right)
				+2\tau  \partial_t \rho \left(\partial_{tt} \tilde{S}_2+\partial_{t} ((v-\epsilon)\partial_r\tilde{S}_2) \right) \\
				\qquad\qquad\qquad\qquad\qquad+\tau  \partial_{tt} \rho \left(\partial_t \tilde{S}_2+(v-\epsilon)\partial_r\tilde{S}_2 \right) +\partial_{tt} \tilde{S}_2=\lambda \left(\partial_{ttr} v+\dfrac{2}{r}\partial_{tt} v\right).
			\end{cases}
		\end{equation}
		Multiplying equations
		$(\ref{38})_1$ and $(\ref{38})_2$ 
		by $r^2\frac{{P}^{\prime}(\rho)}{\rho}\partial_{tt}\rho$ and $r^2\partial_{tt}v$, respectively, and  integrating the result, we obtain
			\begin{align}\label{3.47}
				\nonumber
				&\frac{\mathrm{d}}{\mathrm{d} t}\int_{1}^{+\infty}\bigg(\frac{r^2 P^{\prime}(\rho)}{2\rho} (\partial_{tt}\rho)^2+\frac{r^2\rho}{2}(\partial_{tt}v)^2
				\bigg)\mathrm{d}r\\ \nonumber
				&+\int_{1}^{+\infty}\bigg(-[\frac{ P^{\prime}(\rho)}{\rho}]_t\frac{r^2}{2}(\partial_{tt}\rho)^2
				+2r^2\partial_{t}\rho (\partial_{tt}v)^2
				+r^2\partial_{tt}\rho [\frac{1}{2}(\partial_{t}v)^2]_t
				+r^2\partial_{tt}\rho \partial_{tt}v (\frac{v^2}{2})_r\\ \nonumber
				&\qquad\qquad\quad+r^2\partial_{t}\rho \partial_{r}v \left[(\partial_{t}v)^2\right]_t
				+2r^2\partial_{t}\rho v \partial_{tt}v \partial_{tr}v
				+r^2\rho \partial_{r}v (\partial_{tt}v)^2+r^2\rho \left[(\partial_{t}v)^2\right]_r \partial_{tt}v\bigg)\mathrm{~~d}r\\ \nonumber
				&+\int_{1}^{+\infty}\left(r^2\partial_{ttr}(\rho v)\frac{ P^{\prime}(\rho)}{\rho}\partial_{tt}\rho
				+2r\partial_{tt}(\rho v)\frac{ P^{\prime}(\rho)}{\rho}\partial_{tt}\rho+r^2\partial_{ttr}P \partial_{tt}v\right) \mathrm{d}r\\ 
				=&\int_{1}^{+\infty}\left(\frac{2}{3}r^2\partial_{tt}v\partial_{ttr} \tilde{S}_1 +2r\partial_{tt}v\partial_{tt} \tilde{S}_1+r^2\partial_{tt}v\partial_{ttr} \tilde{S}_2\right)\mathrm{d}r.
			\end{align}

		It is straightforward to see that the second term in the above equation can be estimated by $CE(t)^{\frac{1}{2}} \mathcal{D}(t)$.
		We now focus on the last term on the left-hand side of the above equation, which is given by
		\begin{equation}\nonumber
			\begin{aligned}
				&\int_{1}^{+\infty}\left(r^2\partial_{ttr}(\rho v)\frac{ P^{\prime}(\rho)}{\rho}\partial_{tt}\rho
				+2r\partial_{tt}(\rho v)\frac{ P^{\prime}(\rho)}{\rho}\partial_{tt}\rho+r^2\partial_{ttr}P \partial_{tt}v\right) \mathrm{d}r\\
				=&\int_{1}^{+\infty} \left(r^2\partial_{ttr}vP^{\prime}(\rho)\partial_{tt}\rho
				+2r\partial_{tt}vP^{\prime}(\rho)\partial_{tt}\rho+r^2\partial_{ttr}P \partial_{tt}v\right)\mathrm{d}r \\
				&+\int_{1}^{+\infty}\bigg(\frac{r^2 P^{\prime}(\rho)}{\rho}\partial_{tt}\rho
				\left(\partial_{r}\rho \partial_{tt}v+2\partial_{t}\rho \partial_{tr}v +2\partial_{tr}\rho \partial_{t}v+\partial_{tt}\rho \partial_{r}v\right)\\
				&-\left(r^2\frac{ P^{\prime}(\rho)}{\rho}v\right)_r \frac{1}{2}(\partial_{tt}\rho)^2
				+2r\frac{ P^{\prime}(\rho)}{\rho} \partial_{tt}\rho\left(2\partial_{t}\rho \partial_{t}v+\partial_{tt}\rho v\right)\bigg) \mathrm{d}r\\
				\geq&\int_{1}^{+\infty} \left(r^2\partial_{ttr}vP^{\prime}(\rho)\partial_{tt}\rho
				+2r\partial_{tt}vP^{\prime}(\rho)\partial_{tt}\rho+r^2\partial_{ttr}P \partial_{tt}v\right) \mathrm{d}r
				-C E(t)^{\frac{1}{2}} \mathcal{D}(t).
			\end{aligned}
		\end{equation}
		Furthermore,
		\begin{equation}\nonumber
			\begin{aligned}
				&\int_{1}^{+\infty} \left(r^2\partial_{ttr}vP^{\prime}(\rho)\partial_{tt}\rho
				+2r\partial_{tt}vP^{\prime}(\rho)\partial_{tt}\rho+r^2\partial_{ttr}P \partial_{tt}v\right) \mathrm{d}r\\
				=&\int_{1}^{+\infty}\left(\partial_{r}\left(r^2\partial_{tt}P\partial_{tt}v\right)-r^2P^{\prime\prime}(\rho)(\partial_{t}\rho)^2\partial_{ttr}v
				-2rP^{\prime\prime}(\rho)(\partial_{t}\rho)^2\partial_{tt}v\right)\mathrm{d}r\\
				=&\int_{1}^{+\infty}\left(r^2 \partial_{r}\left(P^{\prime\prime}(\rho)(\partial_{t}\rho)^2\right)\partial_{tt}v\right)\mathrm{d}r \ge-CE(t)^{\frac{1}{2}} \mathcal{D}(t).
			\end{aligned}
		\end{equation}
		Therefore, by combining the above estimation results, equation (\ref{3.47}) reduces to
		\begin{equation}\label{39}
			\begin{aligned}
				&\frac{\mathrm{d}}{\mathrm{d} t}\int_{1}^{+\infty}\bigg(\frac{r^2 P^{\prime}(\rho)}{2\rho} (\partial_{tt}\rho)^2+\frac{r^2\rho}{2}(\partial_{tt}v)^2
				\bigg)\mathrm{d}r\\
				&\leq\int_{1}^{+\infty}\left(\frac{2}{3}r^2\partial_{tt}v\partial_{ttr} \tilde{S}_1 +2r\partial_{tt}v\partial_{tt} \tilde{S}_1+r^2\partial_{tt}v\partial_{ttr} \tilde{S}_2\right)\mathrm{d}r+CE(t)^{\frac{1}{2}} \mathcal{D}(t).
			\end{aligned}
		\end{equation}
		
		Multiplying the equation $(\ref{38})_3$ 
		by  $\dfrac{r^2}{3\mu}\partial_{tt} \tilde{S}_1$ and integrating over $[1,+\infty)$ with respect to $r$ yields
		\begin{equation}\label{42}
			\begin{aligned}
				&\frac{\mathrm{d}}{\mathrm{d} t}\int_{1}^{+\infty}\left[\frac{\tau}{3 \mu} \cdot \frac{r^2 \rho}{2}(\partial_{tt} \tilde{S}_1)^2\right]\mathrm{d}r
				+\int_{1}^{+\infty}\frac{r^2}{3 \mu}(\partial_{tt}\tilde{S}_1)^2\mathrm{d}r\\
				&+\int_{1}^{+\infty}\bigg(\frac{2\tau}{3 \mu}r^2\partial_{t}\rho(\partial_{tt} \tilde{S}_1)^2+\frac{2\tau}{3 \mu}r^2\rho\partial_{t}v\partial_{tr} \tilde{S}_1\partial_{tt} \tilde{S}_1+\frac{\tau}{3 \mu}r^2\rho\partial_{tt}v\partial_{r} \tilde{S}_1\partial_{tt} \tilde{S}_1
				+\frac{2\tau}{3 \mu}r^2\partial_{t}\rho\partial_{t}v\partial_{r} \tilde{S}_1\partial_{tt} \tilde{S}_1\\
				&\qquad\qquad\qquad
				+\frac{2\tau}{3 \mu}r^2\partial_{t}\rho v \partial_{tr} \tilde{S}_1\partial_{tt} \tilde{S}_1
				+\frac{\tau}{3 \mu}r^2\partial_{tt}\rho[\frac{1}{2}(\partial_{t} \tilde{S}_1)^2]_t
				+\frac{\tau}{3 \mu}r^2\partial_{tt}\rho v\partial_{r} \tilde{S}_1\partial_{tt} \tilde{S}_1\bigg)\mathrm{d}r\\
				&+\int_{1}^{+\infty}
				-\frac{\tau\epsilon}{3 \mu}r^2\rho[\frac{1}{2}(\partial_{tt}\tilde{S}_1)^2]_r\mathrm{d}r
				+\int_{1}^{+\infty}\left(-\frac{\tau\epsilon}{3 \mu}r^2\partial_{tt}\rho \partial_{r}\tilde{S}_1\partial_{tt} \tilde{S}_1
				-\frac{2\tau\epsilon}{3 \mu}r^2 \partial_{t}\rho \partial_{tr}\tilde{S}_1\partial_{tt}\tilde{S}_1\right)\mathrm{d}r\\
				&=\int_{1}^{+\infty}\left(\frac{2}{3}r^2\partial_{tt} \tilde{S}_1\partial_{ttr}v
				-\frac{2}{3}r\partial_{tt} \tilde{S}_1\partial_{tt}v\right) \mathrm{d}r
				,
			\end{aligned}
		\end{equation}
		where
		\begin{equation}\nonumber
			\begin{aligned}
				&\int_{1}^{+\infty}\left(\frac{2\tau}{3 \mu}r^2\partial_{t}\rho(\partial_{tt} \tilde{S}_1)^2
				+\frac{\tau}{3 \mu}r^2\partial_{tt}\rho [\frac{1}{2}(\partial_{t} \tilde{S}_1)^2]_t\right)\mathrm{d}r
				\geq-\frac{C}{\tau}E(t)^{\frac{1}{2}} \mathcal{D}(t),\\
				&\int_{1}^{+\infty}\left(-\frac{\tau\epsilon}{3 \mu}r^2\partial_{tt}\rho \partial_{r}\tilde{S}_1\partial_{tt} \tilde{S}_1
				-\frac{2\tau\epsilon}{3 \mu}r^2 \partial_{t}\rho \partial_{tr}\tilde{S}_1\partial_{tt}\tilde{S}_1\right)\mathrm{d}r
				\geq -CE(t)^{\frac{1}{2}} \mathcal{D}(t),\\
				&\int_{1}^{+\infty}\bigg(\frac{2\tau}{3 \mu}r^2\rho\partial_{t}v\partial_{tr} \tilde{S}_1\partial_{tt} \tilde{S}_1+\frac{\tau}{3 \mu}r^2\rho\partial_{tt}v\partial_{r} \tilde{S}_1\partial_{tt} \tilde{S}_1
				+\frac{2\tau}{3 \mu}r^2\partial_{t}\rho\partial_{t}v\partial_{r} \tilde{S}_1\partial_{tt} \tilde{S}_1
				+\frac{2\tau}{3 \mu}r^2\partial_{t}\rho v \partial_{tr} \tilde{S}_1\partial_{tt} \tilde{S}_1\\
				&~~~~
				+\frac{\tau}{3 \mu}r^2\partial_{tt}\rho v\partial_{r} \tilde{S}_1\partial_{tt} \tilde{S}_1-\frac{\tau\epsilon}{3 \mu}r^2\partial_{tt}\rho \partial_{r}\tilde{S}_1\partial_{tt} \tilde{S}_1
				-\frac{2\tau\epsilon}{3 \mu}r^2\partial_{t}\rho\partial_{tr}\tilde{S}_1\partial_{tt}\tilde{S}_1\bigg)\mathrm{d}r
				\geq-CE(t)^{\frac{1}{2}} \mathcal{D}(t),
			\end{aligned}
		\end{equation}
		and
		\begin{equation}\nonumber
			\begin{aligned}
				\int_{1}^{+\infty}-\frac{\tau\epsilon}{3 \mu}r^2\rho[\frac{1}{2}(\partial_{tt}\tilde{S}_1)^2]_r \mathrm{d}r
				&=-\frac{\tau\epsilon}{3 \mu}r^2\rho\frac{1}{2}(\partial_{tt}\tilde{S}_1)^2\bigg|_{1}^{+\infty}
				+\int_{1}^{+\infty}\frac{\tau\epsilon}{3 \mu}(r^2\rho)_r\frac{1}{2}(\partial_{tt}\tilde{S}_1)^2\mathrm{d}r\\
				&=\frac{\tau\epsilon}{6\mu}\rho(t,1)(\partial_{tt}\tilde{S}_1)^2(t,1) +\int_{1}^{+\infty}\frac{\tau\epsilon}{6 \mu}(r^2\partial_{r}\rho+2r\rho)(\partial_{tt}\tilde{S}_1)^2\mathrm{d}r\\
				&\geq  \int_{1}^{+\infty}\frac{\tau\epsilon r^2}{6 \mu}\partial_{r}\rho(\partial_{tt}\tilde{S}_1)^2\mathrm{d}r
				\geq-\frac{C}{\tau}E(t)^{\frac{1}{2}} \mathcal{D}(t).
			\end{aligned}
		\end{equation}
		
		Therefore, we derive
		\begin{equation}
		\begin{aligned}\label{42aa}
			&\frac{\mathrm{d}}{\mathrm{d} t}\int_{1}^{+\infty}\frac{\tau r^2 \rho}{6\mu}(\partial_{tt} \tilde{S}_1)^2\mathrm{d}r
			+\int_{1}^{+\infty}\frac{r^2}{3 \mu}(\partial_{tt}\tilde{S}_1)^2\mathrm{d}r\\
			&\qquad
			\leq \int_{1}^{+\infty}\left(\frac{2}{3}r^2\partial_{tt} \tilde{S}_1\partial_{ttr}v
			-\frac{2}{3}r\partial_{tt} \tilde{S}_1\partial_{tt}v\right) \mathrm{d}r
			+\frac{C}{\tau}E(t)^{\frac{1}{2}} \mathcal{D}(t).
		\end{aligned}
		\end{equation}
		Multiplying the equation $(\ref{38})_4$ 
		by $\dfrac{r^2}{\lambda}\partial_{tt} \tilde{S}_2$ and integrating over $[1,+\infty)$ with respect to $r$ yields
		
			\begin{align}\label{43}
				&\frac{\mathrm{d}}{\mathrm{d} t}\int_{1}^{+\infty}\left[\frac{\tau }{\lambda} \cdot \frac{r^2 \rho}{2}(\partial_{tt} \tilde{S}_2)^2\right]\mathrm{d}r
				+\int_{1}^{+\infty}\frac{r^2}{\lambda}(\partial_{tt}{S_2})^2\mathrm{d}r \nonumber \\
				&+\int_{1}^{+\infty}\bigg(\frac{2\tau }{\lambda}r^2\partial_{t}\rho(\partial_{tt} \tilde{S}_2)^2
				+\frac{2\tau }{\lambda}r^2\rho\partial_{t}v\partial_{tr} \tilde{S}_2\partial_{tt} \tilde{S}_2
				+\frac{\tau }{\lambda}r^2\rho\partial_{tt}v\partial_{r} \tilde{S}_2\partial_{tt} \tilde{S}_2
				+\frac{2\tau }{\lambda}r^2\partial_{t}\rho\partial_{t}v\partial_{r} \tilde{S}_2\partial_{tt} \tilde{S}_2\nonumber \\
				&\qquad\qquad\qquad
				+\frac{2\tau}{\lambda}r^2\partial_{t}\rho v \partial_{tr} \tilde{S}_2\partial_{tt} \tilde{S}_2
				+\frac{\tau}{\lambda}r^2\partial_{tt}\rho[\frac{1}{2}(\partial_{t} \tilde{S}_2)^2]_t
				+\frac{\tau}{\lambda}r^2\partial_{tt}\rho v\partial_{r} \tilde{S}_2\partial_{tt} \tilde{S}_2\bigg)\mathrm{d}r\nonumber \\
				&+\int_{1}^{+\infty}
				-\frac{\tau \epsilon}{\lambda}r^2\rho [\frac{1}{2}(\partial_{tt}\tilde{S}_2)^2]_r\mathrm{d}r
				+\int_{1}^{+\infty}\left(-\frac{\tau \epsilon}{\lambda}r^2\partial_{tt}\rho \partial_{r}\tilde{S}_2\partial_{tt} \tilde{S}_2
				-\frac{2\tau\epsilon}{\lambda}r^2\partial_{t}\rho\partial_{tr}\tilde{S}_2\partial_{tt}\tilde{S}_2\right)\mathrm{d}r \nonumber \\
				\leq&\int_{1}^{+\infty}\left(r^2\partial_{tt} \tilde{S}_2\partial_{ttr}v
				+2r\partial_{tt} \tilde{S}_2\partial_{tt}v\right)\mathrm{d}r.
			\end{align}
	
		Similar to the method used for handling equation (\ref{42}), we derive
		\begin{equation}\label{43a}
			\begin{aligned}
			\frac{\mathrm{d}}{\mathrm{d} t}\int_{1}^{+\infty}\frac{\tau r^2 \rho}{2\lambda}(\partial_{tt} \tilde{S}_2)^2\mathrm{d}r
			+\int_{1}^{+\infty}\frac{r^2}{\lambda}(\partial_{tt}\tilde{S}_2)^2\mathrm{d}r
			\leq \int_{1}^{+\infty}\bigg(r^2\partial_{tt} \tilde{S}_2\partial_{ttr}v
			+2r\partial_{tt} \tilde{S}_2\partial_{tt}v\bigg)\mathrm{d}r\\
			+\frac{C}{\tau }E(t)^{\frac{1}{2}} \mathcal{D}(t).
			\end{aligned}
		\end{equation}
		
		By adding the equations $(\ref{39})$, $(\ref{42aa})$ and $(\ref{43a})$ and using $\partial_{tt}v\big|_{1}^{+\infty}=0$, we obtain
		
		\begin{equation}\label{42a}
			\begin{aligned}
				&\frac{\mathrm{d}}{\mathrm{d} t}\int_{1}^{+\infty}
				\left(\frac{r^2 P^{\prime}(\rho)}{2\rho} (\partial_{tt}\rho)^2
				+\frac{r^2\rho}{2}(\partial_{tt}v)^2
				+\frac{\tau r^2 \rho}{6\mu} (\partial_{tt}\tilde{S}_1)^2 +\frac{\tau r^2 \rho}{2\lambda}  (\partial_{tt}\tilde{S}_1)^2 \right)\mathrm{d}r\\
				&+\int_{1}^{+\infty} \left(\frac{r^2}{3 \mu}(\partial_{tt}\tilde{S}_1)^2
				+\frac{r^2}{\lambda}(\partial_{tt}\tilde{S}_2)^2\right)\mathrm{d}r \le \frac{C}{\tau} E^{\frac{1}{2}}(t) \mathcal{D}(t)
				.
			\end{aligned}
		\end{equation}
		
		
		%\begin{equation}\nonumber
		%	\begin{aligned}
			%		&r^2\partial_{tt}v\left[-\partial_{tt}P+\dfrac{2}{3}\partial_{tt}\tilde{S}_1+\partial_{tt}\tilde{S}_2\right]\bigg|_{1}^{+\infty}=0\\
			%		&=r^2v_{tt}\left[-\partial_{tt}P+\dfrac{2}{3}\partial_{tt}\tilde{S}_1+\partial_{tt}\tilde{S}_2+\dfrac{2}{r}\partial_{t}\tilde{S}_1\right]\bigg|_{1}^{+\infty}-
			%		\dfrac{2}{r}\partial_{t}\tilde{S}_1\cdot r^2 v_{tt} \bigg|_{1}^{+\infty}\\
			%		&=0-2r\partial_{t}\tilde{S}_1\cdot v_{tt} \bigg|_{1}^{+\infty}=0.\\
			%	\end{aligned}
		%\end{equation}
		%where used the  equation $(\ref{15})_2$. 
		%Therefore, by integrating equation $(\ref{34})$ over $(0,t)$, we derive
		
		Integrating the above inequality over $(0, t)$, and noticing that
		\begin{align*}\nonumber
			\int_{1}^{+\infty}\tau^2
			\left(\frac{r^2 P^{\prime}(\rho)}{2\rho} (\partial_{tt}\rho)^2
			+\frac{r^2\rho}{2}(\partial_{tt}v)^2
			+\frac{\tau r^2 \rho}{6\mu} (\partial_{tt}\tilde{S}_1)^2  
			+\frac{\tau r^2 \rho}{2\lambda}  (\partial_{tt}\tilde{S}_1)^2 \right)(t=0,r)\mathrm{d}r\\
			\leq C \|V_0^2\|_{L^2}^2\leq C E(0).
		\end{align*}
		Thus, we have
		\begin{equation}\label{order-2-tt}
			\begin{aligned}
				&\int_{1}^{+\infty} \tau^2
				\left(\frac{r^2 P^{\prime}(\rho)}{2\rho} (\partial_{tt}\rho)^2
				+\frac{r^2\rho}{2}(\partial_{tt}v)^2
				+\frac{\tau r^2 \rho}{6\mu} (\partial_{tt}\tilde{S}_1)^2 +\frac{\tau r^2 \rho}{2\lambda}  (\partial_{tt}\tilde{S}_1)^2 \right)\mathrm{d}r\\
				&+\int_{0}^{t}\int_{1}^{+\infty} \tau^2\left(\frac{r^2}{3 \mu}(\partial_{tt}\tilde{S}_1)^2
				+\frac{r^2}{\lambda}(\partial_{tt}\tilde{S}_2)^2\right)
				\mathrm{d}r\mathrm{d}t
				\leq C\left(E(0)+ E^{\frac{1}{2}}(t) \int_{0}^t \mathcal{D}(s)\mathrm{d}s\right).
			\end{aligned}
		\end{equation}
		This finish the proof of the Lemma $\ref{lem3}$.
	\end{proof}
	
	\begin{lemma}\label{lem4}
		There exists some constant $C$ such that for any $0\leq t \leq T$
		\begin{equation}
			\begin{aligned}
				&\int_{1}^{+\infty}r^2\bigg(\left(\partial_{tr} \rho\right)^2 + (\partial_{tr}v)^2+ (\partial_{t}v)^2
				+\tau (\partial_{tr}\tilde{S}_1)^2
				+\tau(\partial_{t}\tilde{S}_1)^2+\tau(\partial_{tr} \tilde{S}_2)^2\bigg)\mathrm{d}r\\
				&+\int_{0}^{t}\int_{1}^{+\infty}r^2\bigg((\partial_{tr}\tilde{S}_1)^2+ (\partial_{t}\tilde{S}_1)^2
				+(\partial_{tr}\tilde{S}_2)^2\bigg)\mathrm{d}r \mathrm{d} t\\
				&+\int_{0}^{t}\left(\frac{\tau\epsilon}{16\mu} (\partial_{tr}\tilde{S}_1)^2(t,1)
				+\frac{3\tau\epsilon}{8\lambda}(\partial_{tr}\tilde{S}_2)^2(t,1)\right)\mathrm{d}t\\
				&\leq C\left(E(0)+E^{\frac{3}{2}}(t)+ E^{\frac{1}{2}}(t)\int_{0}^{t} \mathcal{D}(s)\mathrm{d}s\right).
			\end{aligned}
		\end{equation}
	\end{lemma}
	\begin{proof}
		Taking the derivative of equation \eqref{3.4} with respect to $t$ and $r$ once, respectively, we get
		\begin{equation}\label{47}
			\begin{cases}
				\partial_{ttr} \rho+\partial_{trr}(\rho v)+\dfrac{2}{r} \partial_{tr}(\rho v)-\dfrac{2}{r^2} \partial_{t}(\rho v)=0, \\
				\rho\partial_{ttr} v+\rho v \partial_{trr} v+\partial_{trr} P-\dfrac{2}{3}\partial_{trr}  \tilde{S}_1-\dfrac{2}{r} \partial_{tr} \tilde{S}_1+\dfrac{2}{r^2}\partial_{t}\tilde{S}_1-\partial_{trr} \tilde{S}_2=g_2,\\
				\tau \rho \partial_{ttr} \tilde{S}_1+\tau \rho v \partial_{trr}\tilde{S}_1 
				-2\mu\left(\partial_{trr} v- \dfrac{1}{r}\partial_{tr} v
				+\dfrac{1}{r^2} \partial_{t}v \right)+\partial_{tr} \tilde{S}_1
				-\tau \epsilon\rho\partial_{trr}\tilde{S}_1
				= g_3, \\
				\tau  \rho \partial_{ttr} \tilde{S}_2+\tau  \rho v \partial_{trr}\tilde{S}_2
				-\lambda \left(\partial_{trr} v+\dfrac{2}{r}\partial_{tr} v-\dfrac{2}{r^2} \partial_{t}v\right)+\partial_{tr} \tilde{S}_2
				-\tau  \epsilon\rho\partial_{trr}\tilde{S}_2
				=g_4,
			\end{cases}
		\end{equation}
		where
		\begin{equation}
			\begin{aligned}\nonumber
				&g_2:=\partial_{t}f_2-\partial_{t}(\rho v)\partial_{rr} v,\\
				&g_3:=\partial_{t}f_3
				-\tau \partial_{t}\rho \partial_{tr} \tilde{S}_1
				-\tau \partial_{t}(\rho v) \partial_{rr}\tilde{S}_1
				+\tau \epsilon\partial_{t}\rho\partial_{rr}\tilde{S}_1,\\
				&g_4:=\partial_{t}f_4-\tau  \partial_{t}\rho \partial_{tr} \tilde{S}_2-\tau  \partial_{t}(\rho v) \partial_{rr}\tilde{S}_2
				+\tau  \epsilon\partial_{t}\rho\partial_{rr}\tilde{S}_2.
			\end{aligned}
		\end{equation}
		
		Multiplying the equation $(\ref{47})_1$ and $(\ref{47})_2$ 
		by  $r^2\frac{\partial_{tr}P}{\rho}$ and $r^2\partial_{tr}v$, respectively, and  integrating the result, we have
		\begin{equation}\label{48}
			\begin{aligned}
				&\frac{\mathrm{d}}{\mathrm{d} t}\int_{1}^{+\infty}\left(\dfrac{r^2 P^{\prime}(\rho)}{2\rho} \left(\partial_{tr} \rho\right)^2+\dfrac{r^2\rho}{2} (\partial_{tr}v)^2\right)\mathrm{d}r
				-\int_{1}^{+\infty}\frac{r^2}{2}[\dfrac{ P^{\prime}(\rho)}{\rho}]_t\left(\partial_{tr} \rho\right)^2\mathrm{d}r\\
				&+\frac{\mathrm{d}}{\mathrm{d} t}\int_{1}^{+\infty} \dfrac{r^2 P^{\prime\prime}(\rho)}{\rho} \partial_{t}\rho\partial_{r}\rho\partial_{tr}\rho\mathrm{~d}r
				-\int_{1}^{+\infty} [\dfrac{r^2 P^{\prime\prime}(\rho)}{\rho} \partial_{t}\rho\partial_{r}\rho]_t\partial_{tr}\rho\mathrm{~d}r\\
				&+\int_{1}^{\infty}\bigg(r^2\partial_{trr}\left(\rho v\right) \frac{\partial_{tr}P}{\rho}
				+2r \partial_{tr} (\rho v) \frac{\partial_{tr}P}{\rho}
				-2\partial_{t} (\rho v) \frac{\partial_{tr}P}{\rho}\bigg)\mathrm{d}r\\
				&+\int_{1}^{+\infty}\bigg(
				r^2 \partial_{trr}P \partial_{tr}v-\dfrac{2 r^2}{3} \partial_{trr} \tilde{S}_1\partial_{tr}v-2r \partial_{tr}\tilde{S}_1 \partial_{tr}v+2\partial_{t}\tilde{S}_1\partial_{tr}v
				-r^2\partial_{trr}\tilde{S}_2\partial_{tr}v\bigg)\mathrm{d}r\\
				&=\int_{1}^{+\infty}r^2g_2 \partial_{tr}v\mathrm{d}r.
			\end{aligned}
		\end{equation} 
		
		Firstly, we have
		\begin{equation}\nonumber
			\begin{aligned}
				-&\int_{1}^{+\infty}\frac{r^2}{2}[\dfrac{ P^{\prime}(\rho)}{\rho}]_t\left(\partial_{tr} \rho\right)^2\mathrm{d}r
				-\int_{1}^{+\infty} [\dfrac{r^2 P^{\prime\prime}(\rho)}{\rho} \partial_{t}\rho\partial_{r}\rho]_t\partial_{tr}\rho\mathrm{~d}r
				\geq -C E^{\frac{1}{2}}(t) \mathcal{D}(t),\\
				&\int_{1}^{+\infty}-2\partial_{t} (\rho v) \frac{\partial_{tr}P}{\rho}\mathrm{d}r
				\geq \int_{1}^{+\infty}-2\partial_{t}v\partial_{tr}P\mathrm{d}r
				-C E^{\frac{1}{2}}(t) \mathcal{D}(t).
			\end{aligned}
		\end{equation}
		Moreover, we have
		\begin{equation}\nonumber
			\begin{aligned}
				&\int_{1}^{\infty}\bigg(r^2\partial_{trr}\left(\rho v\right) \frac{\partial_{tr}P}{\rho}
				+2r \partial_{tr} (\rho v) \frac{\partial_{tr}P}{\rho}
				+r^2 \partial_{trr}P \partial_{tr}v\bigg)\mathrm{d}r\\
				=&\int_{1}^{\infty}\bigg(r^2\partial_{trr}v \partial_{tr}P+2r \partial_{tr} v \partial_{tr}P
				+r^2 \partial_{trr}P \partial_{tr}v\bigg)\mathrm{d}r\\
				&+\int_{1}^{\infty}\bigg(\frac{r^2\partial_{tr}P}{\rho}(2\partial_{r}\rho \partial_{tr} v+2\partial_{tr}\rho \partial_{r} v+\partial_{t}\rho \partial_{rr} v+\partial_{rr}\rho \partial_{t} v)\bigg)\mathrm{d}r
				-\int_{1}^{\infty}[\frac{r^2vP^{\prime}(\rho)}{2\rho}]_r(\partial_{tr}\rho)^2\mathrm{d}r\\
				&
				-\int_{1}^{\infty}r^2P^{\prime\prime}(\rho)[\frac{\partial_{t}\rho\partial_{r}\rho v}{\rho}]_t \partial_{rr}\rho\mathrm{d}r
				+\int_{1}^{\infty}2r \frac{\partial_{tr}P}{\rho}(\partial_{t}\rho \partial_{r} v+\partial_{r}\rho \partial_{t} v+v \partial_{tr}\rho)\mathrm{d} r\\
				\geq& \int_{1}^{\infty}\partial_{r}\left(r^2\partial_{tr}P\partial_{tr}v\right)\mathrm{d}r -C E^{\frac{1}{2}}(t) \mathcal{D}(t),
			\end{aligned}
		\end{equation}
		and
		\begin{equation}\nonumber
			\int_{1}^{+\infty}r^2g_2 \partial_{tr}v\mathrm{d}r\leq C E^{\frac{1}{2}}(t) \mathcal{D}(t).
		\end{equation}
		
		Thus, combining the above estimates, equation $(\ref{48})$ can be simplified to
		\begin{align}\label{49}
				\frac{\mathrm{d}}{\mathrm{d} t}&\int_{1}^{+\infty}\left(\dfrac{r^2 P^{\prime}(\rho)}{2\rho} \left(\partial_{tr} \rho\right)^2+\dfrac{r^2\rho}{2} (\partial_{tr}v)^2\right)\mathrm{d}r
				+\int_{1}^{+\infty}\partial_{r}\left(r^2\partial_{tr}P\partial_{tr}v\right)\mathrm{d}r
				-\int_{1}^{+\infty}2\partial_{t}v\partial_{tr}P\mathrm{d}r\nonumber \\
				+&\int_{1}^{+\infty}\bigg(-\dfrac{2 r^2}{3}\partial_{trr}\tilde{S}_1 \partial_{tr}v-2r \partial_{tr}\tilde{S}_1 \partial_{tr}v+2\partial_{t}\tilde{S}_1 \partial_{tr}v
				-r^2 \partial_{trr}\tilde{S}_2 \partial_{tr}v\bigg)\mathrm{d}r
				\nonumber \\
				\leq& C E^{\frac{1}{2}}(t) \mathcal{D}(t)
				-\frac{\mathrm{d}}{\mathrm{d} t}\int_{1}^{+\infty} \dfrac{r^2 P^{\prime\prime}(\rho)}{\rho} \partial_{t}\rho\partial_{r}\rho\partial_{tr}\rho\mathrm{~d}r.
		\end{align}
		
		Multiplying the equation $(\ref{47})_3$  by  $\dfrac{r^2}{3\mu}\partial_{tr} \tilde{S}_1$, and  integrating the result, we obtain
		\begin{equation}\label{50}
			\begin{aligned}
				&\frac{\mathrm{d}}{\mathrm{d} t}\int_{1}^{+\infty}\left[\frac{\tau}{3 \mu} \frac{r^2 \rho}{2}(\partial_{tr}\tilde{S}_1)^2\right]\mathrm{d}r
				+\int_{1}^{+\infty}\frac{r^2}{3 \mu}(\partial_{tr}\tilde{S}_1)^2 \mathrm{d}r
				-\int_{1}^{+\infty}
				\frac{\tau\epsilon}{3 \mu}r^2\rho[\frac{1}{2}(\partial_{tr}\tilde{S}_1)^2]_r\mathrm{d}r\\
				&-\int_{1}^{+\infty}\bigg(\frac{2r^2}{3}\partial_{trr}v\partial_{tr}\tilde{S}_1-\frac{2r}{3}\partial_{tr}v\partial_{tr}\tilde{S}_1
				+\frac{2}{3}\partial_{t}v\partial_{tr}\tilde{S}_1\bigg)\mathrm{d}r
				=\int_{1}^{+\infty}\frac{r^2}{3\mu}g_3\partial_{tr} \tilde{S}_1\mathrm{d}r,\\
			\end{aligned}
		\end{equation}
		where
		\begin{equation}\nonumber
			\begin{aligned}
				-\int_{1}^{+\infty}
				\frac{\tau\epsilon}{3 \mu}r^2\rho[\frac{1}{2}(\partial_{tr}\tilde{S}_1)^2]_r\mathrm{d}r
				&=\frac{\tau\epsilon}{6\mu}\rho(t,1)(\partial_{tr}\tilde{S}_1)^2(t,1)
				+\int_{1}^{+\infty}
				\frac{\tau\epsilon}{6\mu}(r^2\rho)_r(\partial_{tr}\tilde{S}_1)^2\mathrm{d}r\\
				&\geq\frac{\tau\epsilon}{6\mu}\rho(t,1)(\partial_{tr}\tilde{S}_1)^2(t,1) -C E^{\frac{1}{2}}(t) \mathcal{D}(t),\\
				\int_{1}^{+\infty}\frac{r^2}{3\mu}g_3\partial_{tr} \tilde{S}_1\mathrm{d}r
				&\leq C E^{\frac{1}{2}}(t) \mathcal{D}(t).\\
			\end{aligned}
		\end{equation}
		Similarly, multiplying the equation $(\ref{47})_4$ by  $\dfrac{r^2}{\lambda}\partial_{tr} \tilde{S}_2$, and then integrating the result, one get 
		\begin{equation}
			\begin{aligned}\label{51}
				&\frac{\mathrm{d}}{\mathrm{d} t}\int_{1}^{+\infty}\left[\frac{\tau }{\lambda} \frac{r^2 \rho}{2}(\partial_{tr}{S_2})^2\right]\mathrm{d}r
				+\int_{1}^{+\infty}\frac{r^2}{\lambda}(\partial_{tr}{S_2})^2
				\mathrm{d}r
				-\int_{1}^{+\infty}
				\frac{\tau \epsilon}{\lambda}r^2\rho[\frac{1}{2}(\partial_{tr}\tilde{S}_2)^2]_r\mathrm{d}r\\
				&-\int_{1}^{+\infty}\bigg(r^2\partial_{trr}v\partial_{tr}{S_2}
				-2r\partial_{tr}v\partial_{tr}{S_2}
				+2\partial_tv\partial_{tr}{S_2}\bigg)\mathrm{d}r
				=\int_{1}^{+\infty}\frac{r^2}{\lambda}g_4\partial_{tr} \tilde{S}_2\mathrm{d}r
				,
			\end{aligned}
		\end{equation}
		where
		\begin{equation}\nonumber
			\begin{aligned}
				-\int_{1}^{+\infty}
				\frac{\tau\epsilon}{\lambda}r^2\rho[\frac{1}{2}(\partial_{tr}\tilde{S}_2)^2]_r\mathrm{d}r
				&=\frac{\tau\epsilon}{2\lambda}\rho(t,1)(\partial_{tr}\tilde{S}_2)^2(t,1)
				+\int_{1}^{+\infty}
				\frac{\tau\epsilon}{2\lambda}(r^2\rho)_r(\partial_{tr}\tilde{S}_2)^2\mathrm{d}r\\
				&\geq\frac{\tau\epsilon}{2\lambda}\rho(t,1)(\partial_{tr}\tilde{S}_2)^2(t,1) -C E^{\frac{1}{2}}(t) \mathcal{D}(t)
				,\\
				\int_{1}^{+\infty}\frac{r^2}{\lambda}g_4\partial_{tr} \tilde{S}_2\mathrm{d}r
				&\leq C E^{\frac{1}{2}}(t) \mathcal{D}(t).\\
			\end{aligned}
		\end{equation}
		Adding equations $(\ref{49})-(\ref{51})$ together and combining the above estimates with \eqref{hu3.4},  we can obtain
		\begin{equation}\label{52}
			\begin{aligned}
				&\frac{\mathrm{d}}{\mathrm{d} t}\int_{1}^{+\infty}\bigg(\dfrac{r^2 P^{\prime}(\rho)}{2\rho} \left(\partial_{tr} \rho\right)^2 +\dfrac{r^2\rho}{2} (\partial_{tr}v)^2+\frac{\tau}{3 \mu} \frac{r^2 \rho}{2} (\partial_{tr}\tilde{S}_1)^2
				+\dfrac{\tau }{\lambda}\frac{r^2 \rho}{2} (\partial_{tr} \tilde{S}_2)^2\bigg)\mathrm{d}r\\
				&+\int_{1}^{+\infty}\bigg(\dfrac{r^2}{3\mu} (\partial_{tr}\tilde{S}_1)^2
				+\dfrac{r^2}{\lambda}(\partial_{tr}\tilde{S}_2)^2\bigg)\mathrm{d}r
				+\frac{\tau\epsilon}{8\mu} (\partial_{tr}\tilde{S}_1)^2(t,1)
				+\frac{3\tau\epsilon}{8\lambda}(\partial_{tr}\tilde{S}_2)^2(t,1)\\
				&+\int_{1}^{+\infty}\bigg(\partial_{r}\left(r^2\partial_{tr}P\partial_{tr}v\right)-\dfrac{2}{3}\partial_{r}(r^2\partial_{tr}\tilde{S}_1\partial_{tr}v) 
				-\partial_{r}(r^2\partial_{tr}\tilde{S}_2\partial_{tr}v)
				\bigg)\mathrm{d}r\\
				&+\int_{1}^{+\infty}\bigg(-2\partial_{t}v\partial_{tr}P+2\partial_{t}\tilde{S}_1\partial_{tr}v -\frac{2}{3}\partial_{t}v\partial_{tr}\tilde{S}_1+2\partial_{t}v\partial_{tr}\tilde{S}_2\bigg)\mathrm{d}r\\
				&\leq 
				C E(t)^{\frac{1}{2}}\mathcal{D}(t)
				-\frac{\mathrm{d}}{\mathrm{d} t}\int_{1}^{+\infty} \dfrac{r^2 P^{\prime\prime}(\rho)}{\rho} \partial_{t}\rho\partial_{r}\rho\partial_{tr}\rho\mathrm{~d}r.  
			\end{aligned}
		\end{equation}
		To handle the last term on the left-hand side of the above equation, by multiplying equation $(\ref{15})_2$ by $2\partial_{t}v$ and $(\ref{15})_3$ by $\dfrac{2}{\mu}\partial_{t}\tilde{S}_1$, respectively,  integrating the results, and adding the results to equation $(\ref{52})$, we get
		\begin{equation}
			\begin{aligned}\label{3.64}
				&\frac{\mathrm{d}}{\mathrm{d} t}\int_{1}^{+\infty}\bigg(\dfrac{r^2 P^{\prime}(\rho)}{2\rho} \left(\partial_{tr} \rho\right)^2 +\dfrac{r^2\rho}{2} (\partial_{tr}v)^2+\frac{\tau}{3 \mu} \frac{r^2 \rho}{2} (\partial_{tr}\tilde{S}_1)^2
				+\dfrac{\tau}{\lambda}\frac{r^2 \rho}{2} (\partial_{tr} \tilde{S}_2)^2\bigg)\mathrm{d}r\\
				&+\int_{1}^{+\infty}\bigg(\dfrac{r^2}{3\mu} (\partial_{tr}\tilde{S}_1)^2
				+\dfrac{r^2}{\lambda}(\partial_{tr}\tilde{S}_2)^2\bigg)\mathrm{d}r
				+\frac{\tau\epsilon}{8\mu} (\partial_{tr}\tilde{S}_1)^2(t,1)
				+\frac{3\tau\epsilon}{8\lambda}(\partial_{tr}\tilde{S}_2)^2(t,1)\\
				&+\int_{1}^{+\infty}\bigg(\partial_{r}\left(r^2\partial_{tr}P\partial_{tr}v\right)-\dfrac{2}{3}\partial_{r}(r^2\partial_{tr}\tilde{S}_1\partial_{tr}v) 
				-\partial_{r}(r^2\partial_{tr}\tilde{S}_2\partial_{tr}v)
				\bigg)\mathrm{d}r\\
				&+\int_{1}^{+\infty}\bigg(-2\partial_{t}v\partial_{tr}P+2\partial_{t}\tilde{S}_1\partial_{tr}v -\frac{2}{3}\partial_{t}v\partial_{tr}\tilde{S}_1+2\partial_{t}v\partial_{tr}\tilde{S}_2\bigg)\mathrm{d}r\\
				&+\int_{1}^{+\infty}\left(\rho\partial_{tt} v+\partial_{tr} P
				-(\dfrac{2}{3}\partial_{tr} \tilde{S}_1+\dfrac{2}{r} \partial_t\tilde{S}_1+\partial_{tr} \tilde{S}_2)\right) 2\partial_{t}v\mathrm{d}r
				\\
				&+\int_{1}^{+\infty}\left(\tau \rho \left(\partial_{tt} \tilde{S}_1+\partial_t ((v-\epsilon) \partial_r\tilde{S}_1) \right) +\partial_t \tilde{S}_1
				-2\mu(\partial_{tr} v- \dfrac{1}{r}\partial_t v )\right)\dfrac{2}{\mu}\partial_{t}\tilde{S}_1\mathrm{d}r
				\\
				&\leq 
				C E(t)^{\frac{1}{2}}\mathcal{D}(t)
				-\frac{\mathrm{d}}{\mathrm{d} t}\int_{1}^{+\infty} \dfrac{r^2 P^{\prime\prime}(\rho)}{\rho} \partial_{t}\rho\partial_{r}\rho\partial_{tr}\rho\mathrm{~d}r.
			\end{aligned}
		\end{equation}
		Notice that
		\begin{equation}\nonumber
			\begin{aligned}
				-\int_{1}^{+\infty}\frac{2\tau\epsilon}{\mu}\rho\partial_{t}\tilde{S}_1
				\partial_{tr}\tilde{S}_1\mathrm{d}r
				&=\frac{\tau\epsilon}{\mu}\rho(t,1)(\partial_{t}\tilde{S}_1)^2(t,1)
				+\int_{1}^{+\infty}\frac{\tau\epsilon}{\mu}\partial_r\rho(\partial_{t}\tilde{S}_1)^2
				\mathrm{d}r
				\geq %\frac{\tau\epsilon}{\mu}\rho(t,1)(\partial_{t}\tilde{S}_1)^2(t,1)
				-C E(t)^{\frac{1}{2}}\mathcal{D}(t),
			\end{aligned}
		\end{equation}
		and we integrate equation (\ref{3.64}) over $(0,t)$ to obtain
		\begin{equation}\nonumber
			\begin{aligned}
				&\int_{1}^{+\infty}\bigg(\dfrac{r^2 P^{\prime}(\rho)}{2\rho} \left(\partial_{tr} \rho\right)^2 +\dfrac{1}{2}\rho\left[r^2 (\partial_{tr}v)^2+2 (\partial_{t}v)^2\right]+ \left[\dfrac{\tau\rho}{6\mu} r^2(\partial_{tr}\tilde{S}_1)^2+\dfrac{\tau\rho}{\mu} \partial_{t}\tilde{S}_1^2\right]
				+\dfrac{\tau r^2\rho}{2\lambda} (\partial_{tr} \tilde{S}_2)^2\bigg)\mathrm{d}r\\
				&\quad+\int_{0}^{t}\int_{1}^{+\infty}\bigg(\left[\dfrac{r^2}{3\mu}(\partial_{tr}\tilde{S}_1)^2+\dfrac{2}{\mu} \partial_{t}\tilde{S}_1^2\right]
				+\dfrac{r^2}{\lambda}(\partial_{tr}\tilde{S}_2)^2\bigg)\mathrm{d}r\mathrm{d}t\\
				&\quad+\int_{0}^{t}\left(\frac{\tau\epsilon}{8\mu}(\partial_{tr}\tilde{S}_1)^2(t,1)
				+\frac{3\tau\epsilon}{8\lambda}(\partial_{tr}\tilde{S}_2)^2(t,1)\right)\mathrm{d}t\\
				&\leq 
				\int_{0}^{t}-\left(r^2\partial_{tr}v \partial_{tr}P
				-\dfrac{2}{3} r^2  \partial_{tr}v \partial_{tr}\tilde{S}_1
				-r^2  \partial_{tr}v \partial_{tr}\tilde{S}_2
				-2\partial_{t}v\partial_{t}\tilde{S}_1\right)\bigg|_{1}^{+\infty} \mathrm{d}t\\
				&\quad+C\left(E(0)+E(t)^{\frac{3}{2}}+ E^{\frac{1}{2}}(t) \int_{0}^t \mathcal{D}(s)\mathrm{d}s\right)
				.
			\end{aligned}
		\end{equation}
		
		Noting that $2\partial_{t}v\partial_{t}\tilde{S}_1\big|_{1}^{+\infty}=0$, and using momentum equation \eqref{22}, one has
		
		\begin{equation}
			\begin{aligned}
				&r^2  \partial_{tr}v\left[-\partial_{tr}P+\frac{2}{3}\partial_{tr}\tilde{S}_1
				+\partial_{tr}\tilde{S}_2\right]\bigg|_{1}^{+\infty}\\
				&=r^2  \partial_{tr}v\left[-\partial_{tr}P+\frac{2}{3}\partial_{tr}\tilde{S}_1+\frac{2}{r}\partial_{t}\tilde{S}_1
				+\partial_{tr}\tilde{S}_2\right]\bigg|_{1}^{+\infty}-2r\partial_{t}\tilde{S}_1\partial_{tr}v\bigg|_{1}^{+\infty}\\
				&=r^2  \partial_{tr}v[\rho\partial_{tt} v+\partial_t \rho  \partial_t v+\partial_t \rho \cdot v \partial_r v+\rho \partial_t( v \partial_r v)]\bigg|_{1}^{+\infty}
				-2r\partial_{t}\tilde{S}_1\partial_{tr}v\bigg|_{1}^{+\infty}\\
				&=2(\partial_{t}\tilde{S}_1\partial_{tr}v)(t,1).
			\end{aligned}
		\end{equation}
		Furthermore, using equation $(\ref{15})_3$, we have
		\begin{equation}
			\begin{aligned}\nonumber
				&\int_{0}^{t}2(\partial_{t}\tilde{S}_1\partial_{tr}v)(t,1)\mathrm{d}t\\
				=&\int_{0}^{t} 2\partial_{t}\tilde{S}_1(t,1)
				\bigg[\frac{\tau\rho}{2\mu} \left(\partial_{tt} \tilde{S}_1+\partial_t ((v-\epsilon)\partial_r\tilde{S}_1) \right) +\frac{\tau\partial_{t}\rho}{2\mu} \left(\partial_t \tilde{S}_1+(v-\epsilon)\partial_r\tilde{S}_1 \right) \\
				&\quad+\frac{1}{2\mu}\partial_t \tilde{S}_1+\dfrac{1}{r}\partial_t v\bigg](t,1)
				\mathrm{d}t\\
				=&\frac{\tau}{2\mu}\rho(\partial_{t}\tilde{S}_1)^2(t,1)\big|_0^t
				+\int_{0}^{t}\frac{\tau}{2\mu}\partial_{t}\rho(\partial_{t}\tilde{S}_1)^2(t,1)\mathrm{d}t
				-\int_{0}^{t}\frac{\epsilon\tau}{\mu}\partial_{t}\rho\partial_{t}\tilde{S}_1\partial_{r}\tilde{S}_1(t,1)\mathrm{d}t\\
				&\quad-\int_{0}^{t}\frac{\epsilon\tau}{\mu}\rho\partial_{t}\tilde{S}_1\partial_{tr}\tilde{S}_1(t,1)\mathrm{d}t
				+\int_{0}^{t}\frac{1}{\mu}(\partial_{t}\tilde{S}_1)^2(t,1)\mathrm{d}t.
			\end{aligned}
		\end{equation}
		Now, we deal with each term on the right hand side of the above equation. Firstly, we have 
		\begin{equation}\nonumber
			\begin{aligned}
				&\int_{0}^{t}\frac{\tau}{2\mu}\partial_{t}\rho(\partial_{t}\tilde{S}_1)^2(t,1)\mathrm{d}t
				-\int_{0}^{t}\frac{\epsilon\tau}{\mu}\partial_{t}\rho\partial_{t}\tilde{S}_1\partial_{r}\tilde{S}_1(t,1)\mathrm{d}t
				\\
				\leq&\int_{0}^{t}\frac{\tau}{2\mu}\|\partial_{t}\rho(\partial_{t}\tilde{S}_1)^2\|_{L^\infty}\mathrm{d}t
				+\int_{0}^{t}\frac{\epsilon\tau}{\mu}\|\partial_{t}\rho\partial_{t}\tilde{S}_1\partial_{r}\tilde{S}_1\|_{L^\infty}\mathrm{d}t\\
				\leq&\int_{0}^{t}\frac{\tau}{2\mu}\|\partial_{t}\rho\|_{H^1}\|\partial_{t}\tilde{S}_1\|_{H^1}^2\mathrm{d}t
				+\int_{0}^{t}\frac{\epsilon\tau}{\mu}\|\partial_{t}\rho\|_{H^1}\|\partial_{t}\tilde{S}_1\|_{H^1}\|\partial_{r}\tilde{S}_1\|_{H^1}\mathrm{d}t\\
				\leq&C\left(E(0)+E^{\frac{3}{2}}(t)+E^{\frac{1}{2}}(t) \int_{0}^t \mathcal{D}(s)\mathrm{d}s\right).
			\end{aligned}
		\end{equation}
		Secondly, we know
		\begin{equation}\nonumber
			\begin{aligned}
				&-\int_{0}^{t}\frac{\epsilon\tau}{\mu}\rho\partial_{t}\tilde{S}_1\partial_{tr}\tilde{S}_1(t,1)\mathrm{d}t\\
				\leq&\frac{\epsilon\tau}{\mu}\int_{0}^{t}\rho(t,1) \left(C(\eta)(\partial_{t}\tilde{S}_1)^2(t,1)+\eta(\partial_{tr}\tilde{S}_1)^2(t,1)\right)\mathrm{d}t\\
				\leq&\frac{\epsilon\tau}{\mu}\eta\int_{0}^{t}\rho(t,1)(\partial_{tr}\tilde{S}_1)^2(t,1)\mathrm{d}t
				+C\left(E(0)+E^{\frac{1}{2}}(t) \int_{0}^t \mathcal{D}(s)\mathrm{d}s\right),
			\end{aligned}
		\end{equation}
		and
		\begin{equation}
			\begin{aligned}\nonumber
				&\int_{0}^{t}\dfrac{1}{\mu} (\partial_{t}\tilde{S}_1)^2(t,1)\mathrm{d}t
				\leq \int_{0}^{t}(\dfrac{1}{\mu}\|\partial_{t}\tilde{S}_1\|_{L^\infty}^2)\mathrm{d}t
				\leq \int_{0}^{t}\dfrac{1}{\mu}\|\partial_{tr}\tilde{S}_1\|_{L^2}\|\partial_{t}\tilde{S}_1\|_{L^2}\mathrm{d}t\\
				&\leq \int_{0}^{t}\dfrac{1}{\mu} (\eta\|\partial_{tr}\tilde{S}_1\|_{L^2}^2+C(\eta)\|\partial_{t}\tilde{S}_1\|_{L^2}^2)\mathrm{d}t
				\leq \int_{0}^{t}\dfrac{\eta}{\mu}\|\partial_{tr}\tilde{S}_1\|_{L^2}^2\mathrm{d}t+C\left(E(0)+E^{\frac{1}{2}}(t) \int_{0}^t \mathcal{D}(s)\mathrm{d}s\right).
			\end{aligned}
		\end{equation}
		Here and after, $\eta$ is any positive constant to be chosen later.

		Combining the above results, one can derive
		\begin{equation}\label{57}
			\begin{aligned}
				&\int_{1}^{+\infty}\bigg(\dfrac{r^2 P^{\prime}(\rho)}{2\rho} \left(\partial_{tr} \rho\right)^2 +\dfrac{1}{2}\rho\left[r^2 (\partial_{tr}v)^2+2 (\partial_{t}v)^2\right]
				+\dfrac{\tau\rho r^2}{6\mu}  (\partial_{tr}\tilde{S}_1)^2
				+\dfrac{\tau\rho }{\mu} (\partial_{t}\tilde{S}_1)^2
				+\dfrac{\tau\rho r^2}{2\lambda} (\partial_{tr} \tilde{S}_2)^2\bigg)\mathrm{d}r\\
				&+\int_{0}^{t}\left(\frac{\tau\epsilon}{16\mu} (\partial_{tr}\tilde{S}_1)^2(t,1)
				+\frac{3\tau\epsilon}{8\lambda}\rho(t,1)(\partial_{tr}\tilde{S}_2)^2(t,1)\right)\mathrm{d}t\\
				&+\int_{0}^{t}\int_{1}^{+\infty}\bigg(\dfrac{r^2}{4\mu}(\partial_{tr}\tilde{S}_1)^2+\frac{2}{\mu} (\partial_{t}\tilde{S}_1)^2
				+\dfrac{r^2}{\lambda}(\partial_{tr}\tilde{S}_2)^2\bigg)\mathrm{d}r\\
				&\leq \frac{\tau}{2\mu}\rho(\partial_{t}\tilde{S}_1)^2(t,1)+C\left(E(0)+E^{\frac{3}{2}}(t)+ E^{\frac{1}{2}}(t)\int_{0}^{t} \mathcal{D}(s)\mathrm{d}s\right).
			\end{aligned}
		\end{equation}
		By applying the {\it Gagliardo-Nirenberg} inequality and {\it Young} inequality, we know
		\begin{equation}
			\begin{aligned}\nonumber
				&\frac{\tau}{2\mu}\rho(\partial_{t}\tilde{S}_1)^2(t,1)
				\leq \frac{\tau}{\mu}r^2(\partial_{t}\tilde{S}_1)^2(t,1)
				\leq \frac{\tau}{\mu}r^2\|\partial_{t}\tilde{S}_1\|_{L^{\infty}}^2
				\leq\frac{\tau}{\mu}r^2\|\partial_{tr}\tilde{S}_1\|_{L^{2}}\|\partial_{t}\tilde{S}_1\|_{L^{2}}\\
				&\leq \frac{\tau}{\mu}r^2(\eta\|\partial_{tr}\tilde{S}_1\|_{L^{2}}^2+C(\eta)\|\partial_{t}\tilde{S}_1\|_{L^{2}}^2)
				\leq \frac{\tau\eta}{\mu} r^2\|\partial_{tr}\tilde{S}_1\|_{L^{2}}^2
				+C\left(E(0)+E^{\frac{1}{2}}(t) \int_{0}^t \mathcal{D}(s)\mathrm{d}s\right).
			\end{aligned}
		\end{equation}
		Choosing $\eta<\frac{1}{12}$, 
		we get the Lemma \ref{lem4} immediately.
	\end{proof}
	
	Next, before estimating the second-order partial derivative with respect to $r$, we first derive the second-order dissipation estimates for $(\rho, v)$.
	
	\begin{lemma}\label{lem5}
		There exists a constant $C$ such that
		\begin{equation}
			\begin{aligned}
				&\int_0^t\int_{1}^{+\infty}r^2\left((\partial_{tt}v)^2+(\partial_{tr}v)^2+(\partial_{tt}\rho)^2+(\partial_{tr}\rho)^2\right)\mathrm{d}r\mathrm{d}t
				\leq C\left(E(0)+E^{\frac{3}{2}}(t)+E^{\frac{1}{2}}(t)\int_{0}^{t} \mathcal{D}(s)\mathrm{d}s\right)
			\end{aligned}
		\end{equation}
		and
		\begin{equation}
			\begin{aligned}
				&\int_0^t\int_{1}^{+\infty}r^2(\partial_{rr}v)^2\mathrm{d}r\mathrm{d}t\\
				&\leq
				C\left(\epsilon\int_{0}^{t}\int_{1}^{+\infty}r^2((\partial_{rr}\tilde{S}_1)^2+(\partial_{rr}\tilde{S}_2)^2)\mathrm{d}r\mathrm{d}t+E(0)+E^{\frac{3}{2}}(t)+E^{\frac{1}{2}}(t)\int_{0}^{t} \mathcal{D}(s)\mathrm{d}s\right).
			\end{aligned}
		\end{equation}
	\end{lemma}
	\begin{proof}
		We manipulate $\dfrac{\lambda}{\mu}(\ref{15})_{3}+(\ref{15})_{4}$
		and by applying Lemmas \ref{lem1}-\ref{lem3}, we can obtain
		\begin{equation}\label{2dis1}
			\begin{aligned}
				\int_0^t\int_{1}^{+\infty}r^2(\partial_{tr}v)^2
				\mathrm{d}r\mathrm{d}t
				\leq&\int_0^t\int_{1}^{+\infty}\tau^2r^2((\partial_{tt}\tilde{S}_1)^2+\epsilon(\partial_{tr}\tilde{S}_1)^2+(\partial_{t}\tilde{S}_1)^2+\epsilon(\partial_{r}\tilde{S}_1)^2)\mathrm{d}r\mathrm{d}t\\
				&+\int_0^t\int_{1}^{+\infty}\tau^2r^2((\partial_{tt}\tilde{S}_2)^2+\epsilon(\partial_{tr}\tilde{S}_2)^2+(\partial_{t}\tilde{S}_2)^2+\epsilon(\partial_{r}\tilde{S}_2)^2)\mathrm{d}r\mathrm{d}t\\
				\leq& C(E(0)+E^{\frac{3}{2}}(t)+E^{\frac{1}{2}}(t)\int_{0}^{t} \mathcal{D}(s)\mathrm{d}s).
			\end{aligned}
		\end{equation}
		%	Similarly, by computing $-\dfrac{\lambda}{2\mu}(\ref{15})_{3}+(\ref{15})_{4}$, we can get
		%	\begin{equation}
			%		\int_0^t\int_{1}^{+\infty}r^2(\partial_{t}v)^2
			%		\mathrm{d}r\mathrm{d}t
			%		\leq C(E(0)+E^{\frac{1}{2}}(t)\int_{0}^{t} \mathcal{D}(s)\mathrm{d}s).
			%	\end{equation}
		
		From equation $(\ref{15})_{1}$, we get immediately
		\begin{equation}\label{2dis2}
			\int_0^t\int_{1}^{+\infty}r^2(\partial_{tt}\rho)^2
			\mathrm{d}r\mathrm{d}t
			\leq C(E(0)+E^{\frac{3}{2}}(t)+E^{\frac{1}{2}}(t)\int_{0}^{t} \mathcal{D}(s)\mathrm{d}s).
		\end{equation}
		On the other hand, by multiplying the equation
		$(\ref{15})_{2}$ by $r^2\partial_{tt}v$, it yields
		\begin{equation}\nonumber
			\begin{aligned}
				&\int_0^t\int_{1}^{+\infty}r^2\rho(\partial_{tt}v)^2\mathrm{d}r\mathrm{d}t\\
				\leq&-\int_0^t\frac{\mathrm{d}}{\mathrm{d}t}\int_{1}^{+\infty}r^2\partial_{tr}P\partial_{t}v\mathrm{d}r\mathrm{d}t
				-\int_0^t\int_{1}^{+\infty}r^2\partial_{tt}P\partial_{tr}v\mathrm{d}r\mathrm{d}t
				-\int_0^t\int_{1}^{+\infty}2r\partial_{tt}P\partial_{t}v\mathrm{d}r\mathrm{d}t\\
				&+\int_0^t\int_{1}^{+\infty}\frac{2r^2}{3}\partial_{tr}\tilde{S}_1\partial_{tt}v\mathrm{d}r\mathrm{d}t
				+\int_0^t\int_{1}^{+\infty}2r\partial_{t}\tilde{S}_1\partial_{tt}v\mathrm{d}r\mathrm{d}t
				+\int_0^t\int_{1}^{+\infty}r^2\partial_{tr}\tilde{S}_2\partial_{tt}v\mathrm{d}r\mathrm{d}t\\
				&+C(E(0)+E^{\frac{1}{2}}(t)\int_{0}^{t} \mathcal{D}(s)\mathrm{d}s)\\
				\leq&\frac{1}{2}\int_0^t\int_{1}^{+\infty}r^2(\partial_{tt}v)^2\mathrm{d}r\mathrm{d}t
				+C\int_{1}^{+\infty}r^2((\partial_{tr}\rho)^2+(\partial_{t}v)^2)\mathrm{d}r\\
				&+C\int_0^t\int_{1}^{+\infty}r^2((\partial_{tt}\rho)^2
				+(\partial_{tr}v)^2+(\partial_{t}v)^2
				+(\partial_{tr}\tilde{S}_1)^2+(\partial_{t}\tilde{S}_1)^2+(\partial_{tr}\tilde{S}_2)^2)
				\mathrm{d}r\mathrm{d}t\\
				&+C(E(0)+E^{\frac{1}{2}}(t)\int_{0}^{t} \mathcal{D}(s)\mathrm{d}s)
			\end{aligned}
		\end{equation}
		which implies
		\begin{equation}\label{2dis3}
			\int_0^t\int_{1}^{+\infty}r^2(\partial_{tt}v)^2
			\mathrm{d}r\mathrm{d}t
			\leq C(E(0)+E^{\frac{3}{2}}(t)+E^{\frac{1}{2}}(t)\int_{0}^{t} \mathcal{D}(s)\mathrm{d}s).
		\end{equation}
		By equation $(\ref{15})_2$, we also get
		\begin{equation}\label{2dis4}
			\int_0^t\int_{1}^{+\infty}r^2(\partial_{tr}\rho)^2
			\mathrm{d}r\mathrm{d}t
			\leq C(E(0)+E^{\frac{3}{2}}(t)+E^{\frac{1}{2}}(t)\int_{0}^{t} \mathcal{D}(s)\mathrm{d}s).
		\end{equation}
		
		Furthermore, by manipulating $\dfrac{\lambda}{\mu}(\ref{22})_{3}+(\ref{22})_{4}$, we can obtain
		
		\begin{equation*}%\label{2dis5}
			\begin{aligned}
				\int_0^t\int_{1}^{+\infty}r^2(\partial_{rr}v)^2
				\mathrm{d}r\mathrm{d}t
				\leq&\int_0^t\int_{1}^{+\infty}\tau^2r^2((\partial_{tr}\tilde{S}_1)^2+\epsilon(\partial_{rr}\tilde{S}_1)^2+(\partial_{t}\tilde{S}_1)^2+\epsilon(\partial_{r}\tilde{S}_1)^2)\mathrm{d}r\mathrm{d}t\\
				&+\int_0^t\int_{1}^{+\infty}\tau^2r^2((\partial_{tr}\tilde{S}_2)^2+\epsilon(\partial_{rr}\tilde{S}_2)^2+(\partial_{t}\tilde{S}_2)^2+\epsilon(\partial_{r}\tilde{S}_2)^2)\mathrm{d}r\mathrm{d}t\\
				\leq& C\epsilon\int_{0}^{t}\int_{1}^{+\infty}r^2((\partial_{rr}\tilde{S}_1)^2+(\partial_{rr}\tilde{S}_2)^2)\mathrm{d}r\mathrm{d}t
				+C\bigg(E(0)+E^{\frac{3}{2}}(t)\\
				&\qquad\qquad\qquad\qquad\qquad\qquad\qquad\qquad\qquad\qquad+E^{\frac{1}{2}}(t)\int_{0}^{t} \mathcal{D}(s)\mathrm{d}s\bigg).
			\end{aligned}
		\end{equation*}
		
		Thus,  we get the desired result.
	\end{proof}
	
	\begin{lemma}\label{lem8}
		There exists some constant C such that for any $0\leq t\leq T$
		\begin{equation}\label{sed-rr}
			\begin{aligned}
				&\int_{1}^{+\infty}\bigg(\left(r\partial_{rr} \rho+r\partial_{r} \rho\right)^2
				+r^2 (\partial_{rr}v)^2+(\partial_{r}v)^2
				+\tau\left(r^2(\partial_{rr}\tilde{S}_1)^2+(\partial_{r}\tilde{S}_1)^2+\frac{1}{ r^2}(\tilde{S}_1)^2\right)\\
				&+\tau\left(r\partial_{rr} \tilde{S}_2+2\partial_{r} \tilde{S}_2\right)^2\bigg)\mathrm{d}r\\
				&+\int_0^t\int_{1}^{+\infty}\bigg(\left(r^2(\partial_{rr}\tilde{S}_1)^2+(\partial_{r}\tilde{S}_1)^2+\frac{1}{ r^2}(\tilde{S}_1)^2\right)
				+\left(r\partial_{rr} \tilde{S}_2+2\partial_{r} \tilde{S}_2\right)^2\bigg)\mathrm{d}r\mathrm{d}t\\
				&
				+\int_{0}^{t}\left(\frac{\tau\epsilon}{4\mu}(\partial_{rr}\tilde{S}_1)^2(t,1)
				+\frac{3\tau\epsilon}{8\lambda}(\partial_{rr}\tilde{S}_2)^2(t,1)\right)\mathrm{d}t\\
				&\leq C \left(E(0)+E^{\frac{3}{2}}(t)+E^{\frac{1}{2}}(t) \int_{0}^{t}\mathcal{D}(s)\mathrm{d}s
				+\int_{0}^{t}((\partial_{rr}v)^2(t,1)+(\partial_{tr}v)^2(t,1))\mathrm{d}t\right).
				\\ 
			\end{aligned}
		\end{equation}
	\end{lemma}
	\begin{proof}
		Taking derivative to the equations $(\ref{3.4})$ with respect to $r$ twice, we get
		\begin{equation}\label{60}
			\begin{cases}
				\partial_{trr} \rho+\partial_{rrr}(\rho v)+\dfrac{2}{r} \partial_{rr}(\rho v)-\dfrac{4}{r^2} \partial_{r}(\rho v)+\dfrac{4}{r^3}\rho v=0, \\
				\rho\partial_{trr} v+\rho v \partial_{rrr} v+\partial_{rrr} P-\dfrac{2}{3}\partial_{rrr}  \tilde{S}_1-\dfrac{2}{r} \partial_{rr} \tilde{S}_1+\dfrac{4}{r^2}\partial_{r}\tilde{S}_1-\dfrac{4}{r^3}\tilde{S}_1-\partial_{rrr} \tilde{S}_2=l_2,\\
				\tau \rho \partial_{trr} \tilde{S}_1+\tau \rho v \partial_{rrr}\tilde{S}_1 
				-2\mu\left(\partial_{rrr} v- \dfrac{1}{r}\partial_{rr} v
				+\dfrac{2}{r^2} \partial_{r}v-\dfrac{2}{r^3}v \right)+\partial_{rr} \tilde{S}_1-\tau \epsilon\rho\partial_{rrr}\tilde{S}_1
				=l_3, \\
				\tau  \rho \partial_{trr} \tilde{S}_2+\tau  \rho v \partial_{rrr}\tilde{S}_2
				-\lambda \left(\partial_{rrr} v+\dfrac{2}{r}\partial_{rr} v-\dfrac{4}{r^2} \partial_{r}v+\dfrac{4}{r^3}v\right)+\partial_{rr} \tilde{S}_2 -\tau  \epsilon\rho\partial_{rrr}\tilde{S}_2
				=l_4,
			\end{cases}
		\end{equation}
		where
		\begin{equation}
			\begin{aligned}\nonumber
				&l_2:=\partial_{r}f_2-\partial_{r}(\rho v)\partial_{rr} v,\\
				&l_3:=\partial_{r}f_3
				-\tau \partial_{r}\rho \partial_{tr} \tilde{S}_1
				-\tau \partial_{r}(\rho v) \partial_{rr}\tilde{S}_1
				+\tau \epsilon\partial_{r}\rho\partial_{rr}\tilde{S}_1,\\
				&l_4:=\partial_{r}f_4-\tau  \partial_{r}\rho \partial_{tr} \tilde{S}_2-\tau  \partial_{r}(\rho v) \partial_{rr}\tilde{S}_2
				+\tau  \epsilon\partial_{r}\rho\partial_{rr}\tilde{S}_2.
			\end{aligned}
		\end{equation}
		
		Multiplying equations $(\ref{60})_1$ and $(\ref{60})_2$ 
		by  $r^2\dfrac{\partial_{rr}P}{\rho}$ and $r^2\partial_{rr}v$, respectively, and  integrating the result, we get
		\begin{align}\label{3.66}
				&\frac{\mathrm{d}}{\mathrm{d} t}\int_{1}^{+\infty}\left(\frac{r^2 P^{\prime}(\rho)}{2\rho} \left(\partial_{rr} \rho\right)^2+\frac{r^2\rho}{2} (\partial_{rr}v)^2\right)\mathrm{d}r
				-\int_{1}^{+\infty}\frac{r^2}{2}[\frac{P^{\prime}(\rho)}{\rho}]_t\left(\partial_{rr} \rho\right)^2\mathrm{d}r \nonumber\\
				&+\frac{\mathrm{d}}{\mathrm{d} t}\int_{1}^{+\infty}\frac{r^2 P^{\prime\prime}(\rho)}{\rho} \partial_{rr} \rho(\partial_{r}\rho)^2\mathrm{d}r -\int_{1}^{+\infty}r^2[\frac{P^{\prime\prime}(\rho)}{\rho}(\partial_r\rho)^2]_t\partial_{rr}\rho\mathrm{d}r \nonumber\\
				&+\int_{1}^{+\infty}\bigg(r^2\partial_{rrr}\left(\rho v\right)  \frac{\partial_{rr}P}{\rho}+2r \partial_{rr} (\rho v) \frac{\partial_{rr}P}{\rho} 
				-4\partial_{r} (\rho v) \frac{\partial_{rr}P}{\rho}+\frac{4 v\partial_{rr}P}{r}\bigg)\mathrm{d}r \nonumber\\
				&+\int_{1}^{+\infty}\bigg(
				r^2 \partial_{rrr}P \partial_{rr}v-\frac{2 r^2}{3} \partial_{rrr}\tilde{S}_1 \partial_{rr}v-2r\partial_{rr}\tilde{S}_1\partial_{rr}v+4\partial_{r}\tilde{S}_1  \partial_{rr}v-\frac{4}{r}\tilde{S}_1\partial_{rr}v
				-r^2 \partial_{rrr}\tilde{S}_2\partial_{rr}v\bigg)\mathrm{d}r \nonumber\\
				&=\int_{1}^{+\infty}r^2l_2 \partial_{rr}v\mathrm{d}r.
		\end{align}
		Firstly, we have
		\begin{equation}\nonumber
			\begin{aligned}
				&-\int_{1}^{+\infty}\frac{r^2}{2}[\frac{P^{\prime}(\rho)}{\rho}]_t\left(\partial_{rr} \rho\right)^2\mathrm{d}r -\int_{1}^{+\infty}r^2[\frac{P^{\prime\prime}(\rho)}{\rho}(\partial_r\rho)^2]_t\partial_{rr}\rho\mathrm{d}r
				\geq -C E(t)^{\frac{1}{2}}\mathcal{D}(t),
			\end{aligned}
		\end{equation}
		and
		\begin{equation}\nonumber
			\begin{aligned}
				&\int_{1}^{+\infty}\bigg(r^2\partial_{rrr}\left(\rho v\right)  \frac{\partial_{rr}P}{\rho}+2r \partial_{rr} (\rho v) \frac{\partial_{rr}P}{\rho}+r^2 \partial_{rrr}P \partial_{rr}v 
				\bigg)\mathrm{d}r\\
				&=\int_{1}^{+\infty}(r^2 \partial_{rrr}v\partial_{rr}P+2r \partial_{rr}v\partial_{rr}P+r^2 \partial_{rrr}P \partial_{rr}v)\mathrm{d}r\\
				&~~~~+\int_{1}^{+\infty}r^2(3\partial_{r}\rho\partial_{rr}v+3\partial_{rr}\rho\partial_{r}v+v\partial_{rrr}\rho)\frac{\partial_{rr}P}{\rho} \mathrm{d}r+\int_{1}^{+\infty}2r(v\partial_{rr}\rho+2\partial_{r}\rho\partial_{r}v)\frac{\partial_{rr}P}{\rho} \mathrm{d}r\\
				&\geq \int_{1}^{+\infty}\partial_{r}(r^2 \partial_{rr}v\partial_{rr}P)\mathrm{d}r
				-C E(t)^{\frac{1}{2}}\mathcal{D}(t)
			\end{aligned}
		\end{equation}
		where we used
		\begin{equation}\nonumber
			\begin{aligned}
				\int_{1}^{+\infty}r^2v\partial_{rrr}\rho\frac{\partial_{rr}P}{\rho} \mathrm{d}r&=
				\int_{1}^{+\infty}\frac{P^\prime(\rho)}{\rho}r^2v((\partial_{rr}\rho)^2)_r
				+\int_{1}^{+\infty}\frac{P^{\prime\prime}(\rho)(\partial_{r}\rho)^2}{\rho}r^2v\partial_{rrr}\rho\\
				&=-\int_{1}^{+\infty}\frac{P^\prime(\rho)}{\rho}(r^2v)_r(\partial_{rr}\rho)^2
				-\int_{1}^{+\infty}(\frac{P^{\prime\prime}(\rho)(\partial_{r}\rho)^2}{\rho}r^2v)_r\partial_{rr}\rho\\
				&\geq-C E(t)^{\frac{1}{2}}\mathcal{D}(t).
			\end{aligned}
		\end{equation}
		Moreover, we know
		\begin{equation}\nonumber
			\begin{aligned}
				-\int_{1}^{+\infty}4\partial_{r} (\rho v) \frac{\partial_{rr}P}{\rho}\mathrm{d}r
				\geq -\int_{1}^{+\infty}4\partial_{r}v\partial_{rr}P\mathrm{d}r
				-C E(t)^{\frac{1}{2}}\mathcal{D}(t),
				~~~~\int_{1}^{+\infty}r^2l_2 \partial_{rr}v\mathrm{d}r
				\leq C E(t)^{\frac{1}{2}}\mathcal{D}(t).
			\end{aligned}
		\end{equation}
		Thus, we can sort out the equation (\ref{3.66}) to obtain
		\begin{equation}\label{72}
			\begin{aligned}
				&\frac{\mathrm{d}}{\mathrm{d} t}\int_{1}^{+\infty}\left(\frac{r^2 P^{\prime}(\rho)}{2\rho} \left(\partial_{rr} \rho\right)^2+\frac{r^2\rho}{2} (\partial_{rr}v)^2\right)\mathrm{d}r
				+\frac{\mathrm{d}}{\mathrm{d} t}\int_{1}^{+\infty}\frac{r^2 P^{\prime\prime}(\rho)}{\rho} \partial_{rr} \rho(\partial_{r}\rho)^2\mathrm{d}r\\
				&+\int_{1}^{+\infty}(-\frac{2 r^2}{3} \partial_{rrr}\tilde{S}_1 \partial_{rr}v-2r\partial_{rr}\tilde{S}_1\partial_{rr}v+4\partial_{r}\tilde{S}_1  \partial_{rr}v-\frac{4}{r}\tilde{S}_1\partial_{rr}v
				-r^2 \partial_{rrr}\tilde{S}_2 \partial_{rr}v)\mathrm{d}r\\
				&+\int_{1}^{+\infty}\partial_{r}(r^2 \partial_{rr}v\partial_{rr}P)\mathrm{d}r
				+\int_{1}^{+\infty}(-4\partial_{r}v\partial_{rr}P+\frac{4 v}{r}\partial_{rr}P)\mathrm{d}r
				\leq C E(t)^{\frac{1}{2}}\mathcal{D}(t).\\
			\end{aligned}
		\end{equation}
		
		
		The term \((-4\partial_{r}v\partial_{rr}P+\dfrac{4 v}{r}\partial_{rr}P)\) in the above equation cannot be directly estimated. To address this, we exploit the structure of system (\ref{3.4}) by performing the following multiplications:
		multiply equation $(\ref{22})_1$ by $(\dfrac{4\partial_r P}{\rho} + \dfrac{2r\partial_{rr}P}{\rho})$, $(\ref{60})_1$ by $\dfrac{2r\partial_r P}{\rho}$,  $(\ref{3.4})_2$ by $(\dfrac{4}{r^2}v - 2\partial_{rr}v - \dfrac{4}{r}\partial_r v)$,  $(\ref{22})_2$ by $(4\partial_r v + 2r\partial_{rr}v - \dfrac{4}{r}v)$, and $(\ref{60})_2$ by $(2r\partial_r v - 2v)$, respectively. 
		
		Summing these manipulated equations and combining the result with equation (\ref{72}), we obtain
			\begin{align}\label{3.81}
				&\frac{\mathrm{d}}{\mathrm{d} t}\int_{1}^{+\infty}\bigg(\frac{r^2 P^{\prime}(\rho)}{2\rho} \left(\partial_{rr} \rho\right)^2
				+\frac{r^2\rho}{2} (\partial_{rr}v)^2\bigg)\mathrm{d}r
				+\frac{\mathrm{d}}{\mathrm{d} t}\int_{1}^{+\infty}\frac{r^2 P^{\prime\prime}(\rho)}{\rho} \partial_{rr} \rho(\partial_{r}\rho)^2\mathrm{d}r
				 \nonumber \\
				&+\int_{1}^{+\infty}\partial_{r}(r^2 \partial_{rr}v\partial_{rr}P)\mathrm{d}r
				+\int_{1}^{+\infty}(-4\partial_{r}v\partial_{rr}P+\frac{4 v}{r}\partial_{rr}P)\mathrm{d}r \nonumber\\
				&+\int_{1}^{+\infty}\bigg[\left(\partial_{tr}\rho+\partial_{rr}(\rho v)+\dfrac{2}{r} \partial_r(\rho v)-\dfrac{2}{r^2} \rho v\right)(\dfrac{4\partial_{r}{P}}{\rho}+\dfrac{2r\partial_{rr}{P}}{\rho}) \nonumber\\
				&+\left(\partial_{trr}\rho+\partial_{rrr}(\rho v)+\dfrac{2}{r} \partial_{rr}(\rho v)-\dfrac{4}{r^2} \partial_{r}(\rho v)+\dfrac{4}{r^3}\rho v\right) \dfrac{2r\partial_{r}{P}}{\rho} \nonumber\\
				&+\left(\rho\partial_{t}v+\rho v \partial_r v+\partial_r P-\dfrac{2}{3}\partial_r \tilde{S}_1-\dfrac{2}{r} \tilde{S}_1-\partial_r \tilde{S}_2\right)(\dfrac{4}{r^2}v-2\partial_{rr}v-\dfrac{4}{r}\partial_{r}v) \nonumber\\
				&+\left(\rho\partial_{tr}v+\rho v \partial_{rr} v+\partial_{rr} P-\dfrac{2}{3}\partial_{rr}  \tilde{S}_1-\dfrac{2}{r} \partial_r \tilde{S}_1+\dfrac{2}{r^2}\tilde{S}_1-\partial_{rr} \tilde{S}_2-f_2\right)(4\partial_{r}v+2r\partial_{rr}v-\dfrac{4}{r}v) \nonumber\\
				&+\left(\rho\partial_{trr}v+\rho v \partial_{rrr} v+\partial_{rrr} P-\dfrac{2}{3}\partial_{rrr}  \tilde{S}_1-\dfrac{2}{r} \partial_{rr} \tilde{S}_1+\dfrac{4}{r^2}\partial_{r}\tilde{S}_1-\dfrac{4}{r^3}\tilde{S}_1-\partial_{rrr} \tilde{S}_2-l_2\right)
				(2r\partial_{r}v-2v)\bigg]\mathrm{d}r \nonumber\\
				&+\int_{1}^{+\infty}\bigg(-\frac{2 r^2}{3} \partial_{rrr}\tilde{S}_1 \partial_{rr}v-2r\partial_{rr}\tilde{S}_1\partial_{rr}v+4\partial_{r}\tilde{S}_1  \partial_{rr}v-\frac{4}{r}\tilde{S}_1\partial_{rr}v
				-r^2 \partial_{rrr}\tilde{S}_2 \partial_{rr}v\bigg)\mathrm{d}r \nonumber\\
				&\leq C E(t)^{\frac{1}{2}}\mathcal{D}(t).
			\end{align}
		
		Similar to before, we will temporarily not consider the terms involving $\tilde{S}_1$ and $\tilde{S}_2$. 
		We introduce $G_1$ and $G_2$, defined as
		\begin{equation}\nonumber
			\begin{aligned}
				G_1
				&=\partial_{tr} \rho(\dfrac{4\partial_{r}{P}}{\rho}+\dfrac{2r\partial_{rr}{P}}{\rho})
				+\partial_{trr} \rho\dfrac{2r\partial_{r}{P}}{\rho}+(\rho\partial_{t} v+\rho v \partial_{r} v)(\dfrac{4}{r^2}v-2\partial_{rr}v-\dfrac{4}{r}\partial_{r}v)\\
				&\quad+(\rho\partial_{tr} v+\rho v \partial_{rr} v)(4\partial_{r}v+2r\partial_{rr}v-\dfrac{4}{r}v)
				+(\rho\partial_{trr} v+\rho v \partial_{rrr} v)(2r\partial_{r}v-2v),
			\end{aligned}
		\end{equation}
		and
		\begin{equation}\nonumber
			\begin{aligned}
				G_2&=
				(\partial_{rr}(\rho v)+\dfrac{2}{r} \partial_r(\rho v)-\dfrac{2}{r^2} \rho v) (\dfrac{4\partial_{r}{P}}{\rho}+\dfrac{2r\partial_{rr}{P}}{\rho})\\
				&\quad+(\partial_{rrr}(\rho v)+\dfrac{2}{r} \partial_{rr}(\rho v)-\dfrac{4}{r^2} \partial_{r}(\rho v)+\dfrac{4}{r^3}\rho v) \dfrac{2r\partial_{r}{P}}{\rho}\\
				&\quad+\partial_{r} P (\dfrac{4}{r^2}v-2\partial_{rr}v-\dfrac{4}{r}\partial_{r}v)
				+\partial_{rr} P (4\partial_{r}v+2r\partial_{rr}v-\dfrac{4}{r}v)
				+\partial_{rrr} P(2r\partial_{r}v-2v).
			\end{aligned}
		\end{equation}
		Specifically, we have
		%\begin{equation}\nonumber
			\begin{align*}
				&\int_{1}^{+\infty}G_1\mathrm{d}r\\
				\geq&\int_{1}^{+\infty}\bigg(\frac{2 P^{\prime}(\rho)}{\rho}((\partial_{r}\rho)^2)_t
				+\frac{2r P^{\prime\prime}(\rho)}{\rho}\partial_{tr}\rho(\partial_{r}\rho)^2
				+\left[\frac{2rP^{\prime}(\rho)}{\rho}\partial_{r}\rho\partial_{rr}\rho\right]_t 
				-\left[\frac{P^{\prime}(\rho)}{\rho}\right]_t2r\partial_{r}\rho\partial_{rr}\rho\bigg)\mathrm{d}r\\
				&+\int_{1}^{+\infty}\bigg(
				\frac{2\rho}{r^2}(v^2)_t -2\rho\partial_{t} v\partial_{rr} v-\frac{4\rho}{r}\partial_{t} v\partial_{r} v
				+2\rho((\partial_{r}v)^2)_t+2r\rho\partial_{tr} v\partial_{rr} v-\frac{4\rho}{r}v\partial_{tr} v\bigg)\mathrm{d}r\\
				&+\int_{1}^{+\infty}\bigg(
				2r\rho\partial_{trr}v\partial_{r} v-2\rho v\partial_{trr} v+2r\rho v \partial_{rrr}v\partial_{r} v-2\rho v^2 \partial_{rrr}v
				\bigg)\mathrm{d}r-C E(t)^{\frac{1}{2}}\mathcal{D}(t) \displaybreak\\
				\geq&\int_{1}^{+\infty}\bigg(\frac{2 P^{\prime}(\rho)}{\rho}((\partial_{r}\rho)^2)_t
				+\left[\frac{2rP^{\prime}(\rho)}{\rho}\partial_{r}\rho\partial_{rr}\rho\right]_t\bigg)\mathrm{d}r
				+\int_{1}^{+\infty}\bigg((\frac{2\rho}{r^2}v^2)_t -(2\rho v\partial_{rr} v)_t -(\frac{4\rho}{r} v\partial_{r} v)_t\\
				&\qquad\qquad\quad+(2\rho(\partial_{r} v)^2)_t
				+(2r\rho\partial_{rr}v\partial_{r} v)_t
				-(2r\rho v\partial_{r} v)_r\partial_{rr} v
				+(2\rho v^2)_r\partial_{rr} v
				\bigg)\mathrm{d}r
				-C E(t)^{\frac{1}{2}}\mathcal{D}(t)\\
				\geq& \frac{\mathrm{d}}{\mathrm{d} t}\int_{1}^{+\infty}
				\bigg(\frac{2 P^{\prime}(\rho)}{\rho}(\partial_{r}\rho)^2+\frac{P^{\prime}(\rho)}{\rho}2r\partial_{r}\rho\partial_{rr}\rho
				+2\rho(\partial_{r}v)^2+\frac{2\rho}{r^2}v^2+2r\rho\partial_{rr}v\partial_{r}v-2\rho v\partial_{rr}v\\
				&\qquad\qquad\quad-\frac{4\rho v}{r}\partial_{r}v\bigg)\mathrm{d}r
				-C E(t)^{\frac{1}{2}}\mathcal{D}(t),
			\end{align*}
		%\end{equation}
		and
		\begin{equation}\nonumber
			\begin{aligned}
				&\int_{1}^{+\infty}G_2\mathrm{d}r\\
				\geq&
				\int_{1}^{+\infty}\bigg((4\partial_{r}P\partial_{rr}v+\frac{8}{r}\partial_{r}P\partial_{r}v-\frac{8}{r^2}v\partial_{r}P)+(2r\partial_{rr}P\partial_{rr}v
				+4\partial_{rr}P\partial_{r}v-\frac{4}{r}v\partial_{rr}P)\\
				&\qquad+(2r\partial_{rr}P\partial_{rrr}v+4\partial_{r}P\partial_{rr}v
				-\frac{8}{r}\partial_{r}P\partial_{r}v+\frac{8}{r^2}v\partial_{r}P)+(-2\partial_{r}P\partial_{rr}v-\frac{4}{r}\partial_{r}P\partial_{r}v+\frac{4}{r^2}v\partial_{r}P)\\
				&\qquad+(2r\partial_{rr}P\partial_{rr}v+4\partial_{rr}P\partial_{r}v-
				\frac{4}{r}v\partial_{rr}P)+(2r\partial_{rrr}P\partial_{rr}v-2v\partial_{rrr}P)\bigg)\mathrm{d}r
				-C E(t)^{\frac{1}{2}}\mathcal{D}(t)\\
				=&\int_{1}^{+\infty}\bigg(2\partial_{r}P\partial_{rr}v-\frac{4}{r}\partial_{r}P\partial_{r}v+\frac{4}{r^2}\partial_{r}Pv+2r\partial_{r}P\partial_{rrr}v
				+4\partial_{r}P\partial_{rr}v
				+4r\partial_{rr}P\partial_{rr}v+8\partial_{rr}P\partial_{r}v\\
				&\qquad-\frac{8}{r}v\partial_{rr}P+2r\partial_{rrr}P\partial_{r}v-2v\partial_{rrr}P\bigg)\mathrm{d}r
				-C E(t)^{\frac{1}{2}}\mathcal{D}(t)\\
				=&\int_{1}^{+\infty}\bigg(\partial_{r}\left(2r\partial_{r}P\partial_{rr}v\right)+\partial_{r}\left(2r\partial_{rr}P\partial_{r}v\right)+
				\partial_{r}\left(4\partial_{r}v\partial_{r}P\right)
				-2\partial_{r}\left(v\partial_{rr}P\right)
				-\partial_{r}(\frac{4}{r}v\partial_{r}P)\\
				&\qquad+(4\partial_{r}v\partial_{rr}P-\dfrac{4 v}{r}\partial_{rr}P)\bigg)\mathrm{d}r
				-C E(t)^{\frac{1}{2}}\mathcal{D}(t).
			\end{aligned}
		\end{equation}
		To simplify notation, the terms involving $\tilde{S}_1$ and $\tilde{S}_2$ in equation (\ref{3.81}) are denoted as \(G_3\), 
		\begin{align*}\nonumber
				G_3&=
				\left(-\dfrac{2}{3}\partial_r \tilde{S}_1-\dfrac{2}{r} \tilde{S}_1-\partial_r \tilde{S}_2\right)(\dfrac{4}{r^2}v-2\partial_{rr}v-\dfrac{4}{r}\partial_{r}v)\\
				&\quad+\left(-\dfrac{2}{3}\partial_{rr}  \tilde{S}_1-\dfrac{2}{r} \partial_r \tilde{S}_1+\dfrac{2}{r^2}\tilde{S}_1-\partial_{rr} \tilde{S}_2\right)(4\partial_{r}v+2r\partial_{rr}v-\dfrac{4}{r}v)\\
				&\quad+\left(-\dfrac{2}{3}\partial_{rrr}  \tilde{S}_1-\dfrac{2}{r} \partial_{rr} \tilde{S}_1+\dfrac{4}{r^2}\partial_{r}\tilde{S}_1-\dfrac{4}{r^3}\tilde{S}_1-\partial_{rrr} \tilde{S}_2\right)
				(2r\partial_{r}v-2v)\\
				&\quad+\bigg(-\frac{2 r^2}{3} \partial_{rrr}\tilde{S}_1 \partial_{rr}v-2r\partial_{rr}\tilde{S}_1\partial_{rr}v+4\partial_{r}\tilde{S}_1  \partial_{rr}v-\frac{4}{r}\tilde{S}_1\partial_{rr}v
				-r^2 \partial_{rrr}\tilde{S}_2 \partial_{rr}v\bigg)
		\end{align*}
		which will be addressed later by combining equations $(\ref{60})_3$ and $(\ref{60})_4$.
		
		Therefore, equation (\ref{3.81}) can be simplified to
		
		\begin{align}\label{61}
				\frac{\mathrm{d}}{\mathrm{d} t} &\int_{1}^{+\infty}\bigg(\frac{r^2 P^{\prime}(\rho)}{2\rho} \left(\partial_{rr} \rho\right)^2+\frac{ P^{\prime}(\rho)}{\rho}2(\partial_{r}\rho)^2+\frac{P^{\prime}(\rho)}{\rho}2r\partial_{r}\rho\partial_{rr}\rho
				+\frac{r^2\rho}{2} (\partial_{rr}v)^2+2\rho(\partial_{r}v)^2+\frac{2\rho}{r^2}v^2 \nonumber\\
				&\qquad\quad+2r\rho\partial_{rr}v\partial_{r}v-2\rho\partial_{rr}v v-\frac{4\rho v}{r}\partial_{r}v\bigg)\mathrm{d}r
				+\frac{\mathrm{d}}{\mathrm{d} t}\int_{1}^{+\infty}\frac{r^2 P^{\prime\prime}(\rho)}{\rho} \partial_{rr} \rho(\partial_{r}\rho)^2\mathrm{d}r
				+\int_{1}^{+\infty}G_3\mathrm{d}r \nonumber\\
				\leq&
				\int_{1}^{+\infty}I_1\mathrm{d}r+ C E(t)^{\frac{1}{2}}\mathcal{D}(t),
		\end{align}
		where
		
	\begin{align*}				I_1=-\bigg(\partial_{r}\left(r^2\partial_{rr}P\partial_{rr}v\right) +\partial_{r}\left(2r\partial_{r}P\partial_{rr}v\right)+\partial_{r}\left(2r\partial_{rr}P\partial_{r}v\right)+
	\partial_{r}\left(4\partial_{r}v\partial_{r}P\right)-2\partial_{r}\left(v\partial_{rr}P\right)\\
	-\partial_{r}(\frac{4}{r}v\partial_{r}P)\bigg).
	\end{align*}
		
		Multiplying the equation $(\ref{60})_3$ 
		by  $\dfrac{r^2}{3\mu}\partial_{rr} \tilde{S}_1$  and  integrating over \(r \in [1, +\infty)\), we obtain
		
		\begin{equation}\label{62}
			\begin{aligned}
				&\frac{\mathrm{d}}{\mathrm{d} t}\int_{1}^{+\infty}\left[\frac{\tau}{3 \mu} \frac{r^2 \rho}{2}(\partial_{rr} \tilde{S}_1)^2\right]\mathrm{d}r
				+\int_{1}^{+\infty}\frac{r^2}{3 \mu}(\partial_{rr} \tilde{S}_1)^2\mathrm{d}r
				-\int_{1}^{+\infty}
				\frac{\tau\epsilon}{3 \mu}r^2\rho [\frac{1}{2}(\partial_{rr}\tilde{S}_1)^2]_r\mathrm{d}r\\
				&-\int_{1}^{+\infty}\bigg(\frac{2r^2}{3}\partial_{rrr} v\partial_{rr} \tilde{S}_1- \frac{2r}{3}\partial_{rr} v\partial_{rr}\tilde{S}_1
				+\frac{4}{3} \partial_{r}v\partial_{rr} \tilde{S}_1-\frac{4}{3r}v\partial_{rr} \tilde{S}_1
				\bigg)\mathrm{d}r
				=\int_{1}^{+\infty}\frac{r^2}{3\mu}l_3\partial_{rr} \tilde{S}_1
			\end{aligned}
		\end{equation}
		where
		\begin{equation}\nonumber
			\begin{aligned}
				&-\int_{1}^{+\infty}
				\frac{\tau\epsilon}{3 \mu}r^2\rho[\frac{1}{2}(\partial_{rr}\tilde{S}_1)^2]_r\mathrm{d}r
				\geq \frac{\tau\epsilon}{6 \mu}\rho(t,1)(\partial_{rr}\tilde{S}_1)^2(t,1)
				- C E^{\frac{1}{2}}(t) \mathcal{D}(t),\\
				&\int_{1}^{+\infty}\frac{r^2}{3\mu}l_3\partial_{rr} \tilde{S}_1\mathrm{d}r
				\leq C E^{\frac{1}{2}}(t) \mathcal{D}(t).
			\end{aligned}
		\end{equation}
		
		Similarly, multiplying the equation $(\ref{60})_4$ by  $\dfrac{r^2}{\lambda}\partial_{rr} \tilde{S}_2$, we get
		
		\begin{equation}
			\begin{aligned}\label{63}
				&\frac{\mathrm{d}}{\mathrm{d} t}\int_{1}^{+\infty}\left[\frac{\tau }{\lambda} \frac{r^2 \rho}{2}(\partial_{rr} \tilde{S}_2)^2\right]\mathrm{d}r
				+\int_{1}^{+\infty}\frac{r^2}{\lambda}(\partial_{rr} \tilde{S}_2)^2\mathrm{d}r
				-\int_{1}^{+\infty}
				\frac{\tau\epsilon}{\lambda}r^2\rho [\frac{1}{2}(\partial_{rr}\tilde{S}_2)^2]_r\mathrm{d}r\\
				&-\int_{1}^{+\infty}\bigg(r^2\partial_{rrr}v\partial_{rr} \tilde{S}_2 +2r\partial_{rr}v\partial_{rr} \tilde{S}_2
				-4\partial_{r}v\partial_{rr} \tilde{S}_2+\frac{4}{r}v\partial_{rr} \tilde{S}_2\bigg)\mathrm{d}r =\int_{1}^{+\infty}\frac{r^2}{\lambda}l_4\partial_{rr}\tilde{S}_2
			\end{aligned}
		\end{equation}
		where
		\begin{equation}\nonumber
			\begin{aligned}
				&-\int_{1}^{+\infty}
				\frac{\tau \epsilon}{\lambda}r^2\rho[\frac{1}{2}(\partial_{rr}\tilde{S}_2)^2]_r\mathrm{d}r
				\geq \frac{\tau\epsilon}{2\lambda}\rho(t,1)(\partial_{rr}\tilde{S}_2)^2(t,1)
				- C E^{\frac{1}{2}}(t) \mathcal{D}(t),\\
				&\int_{1}^{+\infty}\frac{r^2}{\lambda}l_4\partial_{rr} \tilde{S}_2\mathrm{d}r
				\leq C E^{\frac{1}{2}}(t) \mathcal{D}(t).
			\end{aligned}
		\end{equation}
		Thus, by adding (\ref{62}) and (\ref{63}), we derive
		\begin{equation}\label{3.80}
			\begin{aligned}
				&\frac{\mathrm{d}}{\mathrm{d} t}\int_{1}^{+\infty}\left[\frac{\tau}{3 \mu} \frac{r^2 \rho}{2}(\partial_{rr} \tilde{S}_1)^2+\frac{\tau }{\lambda} \frac{r^2 \rho}{2}(\partial_{rr} \tilde{S}_2)^2\right]\mathrm{d}r
				+\int_{1}^{+\infty}\left(\frac{r^2}{3 \mu}(\partial_{rr} \tilde{S}_1)^2+\frac{r^2}{\lambda}(\partial_{rr} \tilde{S}_2)^2\right)\mathrm{d}r
				\\
				&-\int_{1}^{+\infty}\bigg(\frac{2r^2}{3}\partial_{rrr} v\partial_{rr} \tilde{S}_1- \frac{2r}{3}\partial_{rr} v\partial_{rr}\tilde{S}_1
				+\frac{4}{3} \partial_{r}v\partial_{rr} \tilde{S}_1-\frac{4}{3r}v\partial_{rr} \tilde{S}_1
				\bigg)\mathrm{d}r\\
				&-\int_{1}^{+\infty}\bigg(r^2\partial_{rrr}v\partial_{rr} \tilde{S}_2 +2r\partial_{rr}v\partial_{rr} \tilde{S}_2
				-4\partial_{r}v\partial_{rr} \tilde{S}_2+\frac{4}{r}v\partial_{rr} \tilde{S}_2\bigg)\mathrm{d}r\\
				&+\frac{\tau\epsilon}{6 \mu}\rho(t,1)(\partial_{rr}\tilde{S}_1)^2(t,1)
				+\frac{\tau\epsilon}{2\lambda}\rho(t,1)(\partial_{rr}\tilde{S}_2)^2(t,1)
				\leq C E^{\frac{1}{2}}(t) \mathcal{D}(t).
			\end{aligned}
		\end{equation}
		
		To address the third and fourth terms on the left-hand side of \eqref{3.80}, we exploit the structure of system (\ref{3.4}) and couple them with equation (\ref{61}) via the following operations:\\ 
		multiply equation \((\ref{3.4})_3\) by \((\dfrac{12}{\mu r^2}\tilde{S}_1 - \dfrac{2}{\mu}\partial_{rr}\tilde{S}_1 - \dfrac{4}{\mu r}\partial_{r}\tilde{S}_1)\), 
		\((\ref{22})_3\) by \((\dfrac{4}{3\mu}\partial_{r}\tilde{S}_1 + \dfrac{2r}{3\mu}\partial_{rr}\tilde{S}_1 - \dfrac{4}{\mu r}\tilde{S}_1)\),
		\((\ref{60})_3\) by \((\dfrac{2r}{3\mu}\partial_{r}\tilde{S}_1 - \dfrac{2}{\mu}\tilde{S}_1)\),
		\((\ref{22})_4\) by \((\dfrac{4}{\lambda}\partial_{r}\tilde{S}_2 + \dfrac{2r}{\lambda}\partial_{rr}\tilde{S}_2)\),
		\((\ref{60})_4\) by \(\dfrac{2r}{\lambda}\partial_{r}\tilde{S}_2\), respectively.  
		
		Upon adding the resulting equations to \((\ref{3.80})\), we obtain 
			\begin{align}\label{64}
				&\frac{\mathrm{d}}{\mathrm{d} t}\int_{1}^{+\infty}\left[\frac{\tau}{3 \mu} \frac{r^2 \rho}{2}(\partial_{rr} \tilde{S}_1)^2+\frac{\tau }{\lambda} \frac{r^2 \rho}{2}(\partial_{rr} \tilde{S}_2)^2\right]\mathrm{d}r
				+\int_{1}^{+\infty}(\frac{r^2}{3 \mu}(\partial_{rr} \tilde{S}_1)^2+\frac{r^2}{\lambda}(\partial_{rr} \tilde{S}_2)^2)\mathrm{d}r
				\nonumber\\
				&+\frac{\tau\epsilon}{6 \mu}\rho(t,1)(\partial_{rr}\tilde{S}_1)^2(t,1) +\frac{\tau\epsilon}{2\lambda}\rho(t,1)(\partial_{rr}\tilde{S}_2)^2(t,1) \nonumber\\
				&-\int_{1}^{+\infty}\bigg(\frac{2r^2}{3}\partial_{rrr} v\partial_{rr} \tilde{S}_1- \frac{2r}{3}\partial_{rr} v\partial_{rr}\tilde{S}_1
				+\frac{4}{3} \partial_{r}v\partial_{rr} \tilde{S}_1-\frac{4}{3r}v\partial_{rr} \tilde{S}_1
				\bigg)\mathrm{d}r \nonumber\\
				&-\int_{1}^{+\infty}\bigg(r^2\partial_{rrr}v\partial_{rr} \tilde{S}_2 +2r\partial_{rr}v\partial_{rr} \tilde{S}_2
				-4\partial_{r}v\partial_{rr} \tilde{S}_2+\frac{4}{r}v\partial_{rr} \tilde{S}_2\bigg)\mathrm{d}r \nonumber
				\\
				&+\int_{1}^{+\infty}\bigg[\left(\tau \rho\partial_{t}\tilde{S}_1+\tau \rho v\partial_r\tilde{S}_1  +\tilde{S}_1-2\mu(\partial_r v- \dfrac{v}{r})-\tau\epsilon\rho\partial_{r}\tilde{S}_1\right)(\dfrac{12}{\mu r^2}\tilde{S}_1-\dfrac{2}{\mu}\partial_{rr}\tilde{S}_1-\dfrac{4}{\mu r}\partial_{r}\tilde{S}_1) \nonumber\\
				&\quad
				+\bigg(\tau \rho\partial_{tr}\tilde{S}_1+\tau \rho v \partial_{rr}\tilde{S}_1 
				-2\mu\big(\partial_{rr} v- \dfrac{1}{r}\partial_r v
				+\dfrac{1}{r^2} v \big)+\partial_r \tilde{S}_1 \nonumber \\
				&\qquad\qquad\qquad\qquad\qquad\qquad\qquad\qquad\qquad\qquad\qquad\quad-\tau\epsilon\rho\partial_{rr}\tilde{S}_1
				\bigg)(\dfrac{4}{3\mu}\partial_{r}\tilde{S}_1+\dfrac{2r}{3\mu}\partial_{rr}\tilde{S}_1 -\dfrac{4}{\mu r}\tilde{S}_1) \nonumber\\
				&\quad
				+\bigg(\tau \rho\partial_{trr}\tilde{S}_1+\tau \rho v \partial_{rrr}\tilde{S}_1 
				-2\mu\big(\partial_{rrr} v- \dfrac{1}{r}\partial_{rr} v
				+\dfrac{2}{r^2} \partial_{r}v-\dfrac{2}{r^3}v \big)+\partial_{rr} \tilde{S}_1 \nonumber \\ &\qquad\qquad\qquad\qquad\qquad\qquad\qquad\qquad\qquad\qquad\qquad\quad-\tau\epsilon\rho\partial_{rrr}\tilde{S}_1\bigg)(\dfrac{2r}{3\mu}\partial_{r}\tilde{S}_1 -\dfrac{2}{\mu}\tilde{S}_1) \nonumber\\
				&\quad
				+\left(\tau \rho\partial_{tr}\tilde{S}_2+\tau  \rho v \partial_{rr}\tilde{S}_2
				-\lambda \big(\partial_{rr} v+\dfrac{2}{r}\partial_r v-\dfrac{2}{r^2} v\big)+\partial_r \tilde{S}_2 -\tau\epsilon\rho\partial_{rr}\tilde{S}_2\right)
				(\dfrac{4}{\lambda}\partial_{r}\tilde{S}_2+\dfrac{2r}{\lambda}\partial_{rr}\tilde{S}_2) \nonumber \\
				&\quad
				+\left(\tau \rho\partial_{trr}\tilde{S}_2+\tau  \rho v \partial_{rrr}\tilde{S}_2
				-\lambda \big(\partial_{rrr} v+\dfrac{2}{r}\partial_{rr} v-\dfrac{4}{r^2} \partial_{r}v+\dfrac{4}{r^3}v\big)+\partial_{rr} \tilde{S}_2-\tau\epsilon\rho\partial_{rrr}\tilde{S}_2
				\right)\dfrac{2r}{\lambda}\partial_{r}\tilde{S}_2\bigg]\mathrm{d}r \nonumber \\
				&
				~\leq ~C E^{\frac{1}{2}}(t) \mathcal{D}(t).
			\end{align}
		We rewrite the last three terms on the left-hand side of (\ref{64}) as $\int_{1}^{+\infty}(H_1+H_2+H_3+H_4)\mathrm{d}r$, where
		\begin{equation}\nonumber
			\begin{aligned}
				H_1=&	
				\left(\tau \rho\partial_{t}\tilde{S}_1+\tau \rho v\partial_r\tilde{S}_1  \right)(\dfrac{12}{\mu r^2}\tilde{S}_1-\dfrac{2}{\mu}\partial_{rr}\tilde{S}_1-\dfrac{4}{\mu r}\partial_{r}\tilde{S}_1)\\
				&
				+\left(\tau \rho\partial_{tr}\tilde{S}_1+\tau \rho v \partial_{rr}\tilde{S}_1 
				\right)(\dfrac{4}{3\mu}\partial_{r}\tilde{S}_1+\dfrac{2r}{3\mu}\partial_{rr}\tilde{S}_1 -\dfrac{4}{\mu r}\tilde{S}_1)\\
				&+\left(\tau \rho\partial_{trr}\tilde{S}_1+\tau \rho v \partial_{rrr}\tilde{S}_1 
				\right)(\dfrac{2r}{3\mu}\partial_{r}\tilde{S}_1 -\dfrac{2}{\mu}\tilde{S}_1)\\
				&
				+\left(\tau \rho\partial_{tr}\tilde{S}_2+\tau  \rho v \partial_{rr}\tilde{S}_2
				\right)
				(\dfrac{4}{\lambda}\partial_{r}\tilde{S}_2+\dfrac{2r}{\lambda}\partial_{rr}\tilde{S}_2)
				+\left(\tau \rho\partial_{trr}\tilde{S}_2+\tau  \rho v \partial_{rrr}\tilde{S}_2
				\right)\dfrac{2r}{\lambda}\partial_{r}\tilde{S}_2,
			\end{aligned}
		\end{equation}
		\begin{equation}\nonumber
			\begin{aligned}
				H_2
				=&\tilde{S}_1(\dfrac{12}{\mu r^2}\tilde{S}_1-\dfrac{2}{\mu}\partial_{rr}\tilde{S}_1-\dfrac{4}{\mu r}\partial_{r}\tilde{S}_1)
				+\partial_{r}\tilde{S}_1(\dfrac{4}{3\mu}\partial_{r}\tilde{S}_1+\dfrac{2r}{3\mu}\partial_{rr}\tilde{S}_1 -\dfrac{4}{\mu r}\tilde{S}_1)\\
				&
				+\partial_{rr}\tilde{S}_1(\dfrac{2r}{3\mu}\partial_{r}\tilde{S}_1 -\dfrac{2}{\mu}\tilde{S}_1)
				+\partial_{r}\tilde{S}_2
				(\dfrac{4}{\lambda}\partial_{r}\tilde{S}_2+\dfrac{2r}{\lambda}\partial_{rr}\tilde{S}_2)
				+\dfrac{2r}{\lambda}\partial_{rr}\tilde{S}_2\partial_{r}\tilde{S}_2,
			\end{aligned}
		\end{equation}
		\begin{equation}\nonumber
			\begin{aligned}
				H_3
				=
				&-2\mu(\partial_r v- \dfrac{v}{r})(\dfrac{12}{\mu r^2}\tilde{S}_1-\dfrac{2}{\mu}\partial_{rr}\tilde{S}_1-\dfrac{4}{\mu r}\partial_{r}\tilde{S}_1)\\
				&-2\mu\left(\partial_{rr} v- \dfrac{1}{r}\partial_r v
				+\dfrac{1}{r^2} v \right)(\dfrac{4}{3\mu}\partial_{r}\tilde{S}_1+\dfrac{2r}{3\mu}\partial_{rr}\tilde{S}_1 -\dfrac{4}{\mu r}\tilde{S}_1)\\
				&
				-2\mu\left(\partial_{rrr} v- \dfrac{1}{r}\partial_{rr} v
				+\dfrac{2}{r^2} \partial_{r}v-\dfrac{2}{r^3}v \right)(\dfrac{2r}{3\mu}\partial_{r}\tilde{S}_1 -\dfrac{2}{\mu}\tilde{S}_1)\\
				&-\lambda \left(\partial_{rr} v+\dfrac{2}{r}\partial_r v-\dfrac{2}{r^2} v\right)
				(\dfrac{4}{\lambda}\partial_{r}\tilde{S}_2+\dfrac{2r}{\lambda}\partial_{rr}\tilde{S}_2)
				-\lambda \left(\partial_{rrr} v+\dfrac{2}{r}\partial_{rr} v-\dfrac{4}{r^2} \partial_{r}v+\dfrac{4}{r^3}v\right)\dfrac{2r}{\lambda}\partial_{r}\tilde{S}_2\\
				&-\bigg(\frac{2r^2}{3}\partial_{rrr} v\partial_{rr} \tilde{S}_1- \frac{2r}{3}\partial_{rr} v\partial_{rr}\tilde{S}_1
				+\frac{4}{3} \partial_{r}v\partial_{rr} \tilde{S}_1-\frac{4}{3r}v\partial_{rr} \tilde{S}_1
				\bigg)\\
				&-\bigg(r^2\partial_{rrr}v\partial_{rr} \tilde{S}_2 +2r\partial_{rr}v\partial_{rr} \tilde{S}_2
				-4\partial_{r}v\partial_{rr} \tilde{S}_2+\frac{4}{r}v\partial_{rr} \tilde{S}_2\bigg),
			\end{aligned}
		\end{equation}
		and
		\begin{equation}\nonumber
			\begin{aligned}
				H_4=
				&-\tau\epsilon\rho\partial_{r}\tilde{S}_1(\dfrac{12}{\mu r^2}\tilde{S}_1-\dfrac{2}{\mu}\partial_{rr}\tilde{S}_1-\dfrac{4}{\mu r}\partial_{r}\tilde{S}_1)
				-\tau\epsilon\rho\partial_{rr}\tilde{S}_1
				(\dfrac{4}{3\mu}\partial_{r}\tilde{S}_1+\dfrac{2r}{3\mu}\partial_{rr}\tilde{S}_1 -\dfrac{4}{\mu r}\tilde{S}_1)\\
				&-{\tau\epsilon\rho\partial_{rrr}\tilde{S}_1}(\dfrac{2r}{3\mu}\partial_{r}\tilde{S}_1 -\dfrac{2}{\mu}\tilde{S}_1)
				-\tau\epsilon\rho\partial_{rr}\tilde{S}_2
				(\dfrac{4}{\lambda}\partial_{r}\tilde{S}_2+\dfrac{2r}{\lambda}\partial_{rr}\tilde{S}_2)
				-\tau\epsilon\rho\partial_{rrr}\tilde{S}_2
				\dfrac{2r}{\lambda}\partial_{r}\tilde{S}_2.
			\end{aligned}
		\end{equation}
		Next, we handle each \(\int_{1}^{+\infty}H_i\mathrm{d}r\) individually, \(i=1,2,3,4\).
		\begin{equation}\nonumber
			\begin{aligned}
				&\int_{1}^{+\infty}H_1\mathrm{d}r\\
				\geq&
				\int_{1}^{+\infty}\tau\rho\bigg(
				\frac{12}{\mu r^2}\partial_{t}\tilde{S}_1\tilde{S}_1
				-\frac{2}{\mu}\partial_{t}\tilde{S}_1\partial_{rr}\tilde{S}_1
				-\frac{4}{\mu r} \partial_{t}\tilde{S}_1\partial_{r}\tilde{S}_1
				+\frac{4}{3\mu} \partial_{tr}\tilde{S}_1\partial_{r}\tilde{S}_1
				+\frac{2r}{3\mu} \partial_{tr}\tilde{S}_1\partial_{rr}\tilde{S}_1
				\\
				&\qquad\qquad\quad-\frac{4}{\mu r} \partial_{tr}\tilde{S}_1\tilde{S}_1
				+\frac{2r}{3\mu} \partial_{trr}\tilde{S}_1\partial_{r}\tilde{S}_1
				-\frac{2}{\mu} \partial_{trr}\tilde{S}_1\tilde{S}_1
				+\frac{2r}{3\mu} v\partial_{rrr}\tilde{S}_1\partial_{r}\tilde{S}_1
				-\frac{2}{\mu} v\partial_{rrr}\tilde{S}_1\tilde{S}_1
				\bigg)\mathrm{d}r\\
				&+\int_{1}^{+\infty}\tau\rho\bigg(
				\frac{4}{\lambda}\partial_{tr}\tilde{S}_2\partial_{r}\tilde{S}_2
				+\frac{2r}{\lambda}\partial_{tr}\tilde{S}_2\partial_{rr}\tilde{S}_2
				+\frac{2r}{\lambda}\partial_{trr}\tilde{S}_2\partial_{r}\tilde{S}_2
				+\frac{2r}{\lambda}v\partial_{rrr}\tilde{S}_2\partial_{r}\tilde{S}_2
				\bigg)\mathrm{d}r
				-C E^{\frac{1}{2}}(t) \mathcal{D}(t)\\
				\geq&
				\int_{1}^{+\infty}\tau\rho\bigg(
				\frac{6}{\mu r^2}(\tilde{S}_1^2)_t
				-\frac{2}{\mu}(\tilde{S}_1\partial_{rr}\tilde{S}_1)_t
				-\frac{4}{\mu r}(\partial_{r}\tilde{S}_1\tilde{S}_1)_t
				+\frac{2}{3\mu} [(\partial_{r}\tilde{S}_1)^2]_t
				+\frac{2r}{3\mu} (\partial_{r}\tilde{S}_1\partial_{rr}\tilde{S}_1)_t
				\\
				&\qquad\qquad-(\frac{2r}{3\mu} v\partial_{r}\tilde{S}_1)_r\partial_{rr}\tilde{S}_1+(\frac{2}{\mu} v\tilde{S}_1)_r\partial_{rr}\tilde{S}_1
				\bigg)\mathrm{d}r\\
				&+\int_{1}^{+\infty}\tau\rho\bigg(
				\frac{2}{\lambda}[(\partial_{r} \tilde{S}_2)^2]_t
				+\frac{2r}{\lambda}(\partial_{r} \tilde{S}_2\partial_{rr} \tilde{S}_2)_t
				-(\frac{2r}{\lambda}v\partial_{r} \tilde{S}_2)_r\partial_{rr} \tilde{S}_2
				\bigg)\mathrm{d}r
				-C E^{\frac{1}{2}}(t) \mathcal{D}(t)\\
				\geq&
				\frac{\mathrm{d}}{\mathrm{d} t}\int_{1}^{+\infty}\tau\rho\bigg(
				\frac{2}{3\mu}(\partial_{r}\tilde{S}_1)^2+\frac{6}{\mu r^2}(\tilde{S}_1)^2 +\frac{2r}{3\mu}\partial_{rr}\tilde{S}_1 \partial_{r}\tilde{S}_1-\frac{2}{\mu}\partial_{rr}\tilde{S}_1\tilde{S}_1-\frac{4}{\mu r}\partial_{r}\tilde{S}_1 \tilde{S}_1
				+\frac{2}{\lambda}(\partial_{r} \tilde{S}_2)^2\\
				&\qquad\qquad\qquad\qquad\qquad\qquad\qquad\qquad\qquad\qquad\qquad\qquad\qquad\quad+\frac{2r}{\lambda}\partial_{rr} \tilde{S}_2\partial_{r} \tilde{S}_2\bigg)\mathrm{d}r
				-C E^{\frac{1}{2}}(t) \mathcal{D}(t),
			\end{aligned}
		\end{equation}
		and
			\begin{align*}
				\int_{1}^{+\infty}H_2\mathrm{d}r
				=\int_{1}^{+\infty}\bigg(\frac{4}{3\mu}(\partial_{r}\tilde{S}_1)^2+\frac{12}{\mu r^2}(\tilde{S}_1)^2 +\frac{4r}{3\mu}\partial_{rr}\tilde{S}_1 \partial_{r}\tilde{S}_1-\frac{4}{\mu}\partial_{rr}\tilde{S}_1\tilde{S}_1-\frac{8}{\mu r}\partial_{r}\tilde{S}_1 \tilde{S}_1
				+\frac{4}{\lambda}(\partial_{r}  \tilde{S}_2)^2\\
				+\frac{4r}{\lambda}\partial_{rr} \tilde{S}_2\partial_{r} \tilde{S}_2\bigg)\mathrm{d}r.
			\end{align*}
		
		To address the term $\int_{1}^{+\infty}H_3\mathrm{d}r$, we focus on $H_3+G_3$,
		
		\begin{equation*}
			\begin{aligned}
				&\int_{1}^{+\infty}(H_3+G_3)\mathrm{d}r\\
				=&\int_{1}^{+\infty}\bigg[-\frac{2}{3}(r^2\partial_{rrr}v\partial_{rr}\tilde{S}_1+r^2\partial_{rr}v\partial_{rrr}\tilde{S}_1+2r\partial_{rr}v\partial_{rr}\tilde{S}_1)\\
				&\quad\quad\quad-(\frac{4r}{3}\partial_{rr}v\partial_{rr}\tilde{S}_1+\frac{4r}{3}\partial_{r}v\partial_{rrr}\tilde{S}_1+\frac{4}{3}\partial_{r}v\partial_{rr}\tilde{S}_1)+(\frac{4}{3}\partial_{r}v\partial_{rr}\tilde{S}_1+\frac{4}{3}v\partial_{rrr}\tilde{S}_1)\\
				&\quad\quad\quad-
				(\frac{4r}{3}\partial_{rrr}v\partial_{r}\tilde{S}_1+\frac{4r}{3}\partial_{rr}v\partial_{rr}\tilde{S}_1+\frac{4}{3}\partial_{rr}v\partial_{r}\tilde{S}_1)-\frac{8}{3}(\partial_{rr}v\partial_{r}\tilde{S}_1+\partial_{r}v\partial_{rr}\tilde{S}_1)\\
				&\quad\quad\quad+ (\frac{8}{3r}v\partial_{rr}\tilde{S}_1+\frac{8}{3r}\partial_{r}v\partial_{r}\tilde{S}_1-\frac{8}{3r^2}v\partial_{r}\tilde{S}_1)
				+4(\partial_{rrr}v\tilde{S}_1+\partial_{rr}v\partial_{r}\tilde{S}_1)\\
				&\quad\quad\quad+(\frac{8}{r}\partial_{rr}v\tilde{S}_1+\frac{8}{r}\partial_{r}v
				\partial_{r}\tilde{S}_1-\frac{8}{r^2}\partial_{r}v\tilde{S}_1)
				-( \frac{8}{r^2}\partial_{r}v\tilde{S}_1+\frac{8}{r^2}v\partial_{r}\tilde{S}_1-\frac{16}{r^3}v\tilde{S}_1)\\
				&\quad\quad\quad-(r^2\partial_{rrr}v\partial_{rr}\tilde{S}_2+r^2\partial_{rr}v\partial_{rrr}\tilde{S}_2+2r\partial_{rr}v\partial_{rr}\tilde{S}_2)\\
				&\quad\quad\quad-(2r\partial_{rrr}v\partial_{r}\tilde{S}_2+2r\partial_{rr}v\partial_{rr}\tilde{S}_2+2\partial_{rr}v\partial_{r}\tilde{S}_2)\\
				&\quad\quad\quad-(2r\partial_{rr}v\partial_{rr}\tilde{S}_2+2r\partial_{r}v\partial_{rrr}\tilde{S}_2+2\partial_{r}v\partial_{rr}\tilde{S}_2)
				-(4\partial_{rr}v\partial_{r}\tilde{S}_2+4\partial_{r}v\partial_{rr}\tilde{S}_2)\\
				&\quad\quad\quad+(2\partial_{r}v\partial_{rr}\tilde{S}_2+2v\partial_{rrr}\tilde{S}_2)+(\frac{4}{r}\partial_{r}v\partial_{r}\tilde{S}_2+\frac{4}{r}v\partial_{rr}\tilde{S}_2-\frac{4}{r^2}v\partial_{r}\tilde{S}_2)\bigg]\mathrm{d}r\\
				:=&-\int_{1}^{+\infty}(I_2+I_3)\mathrm{d}r
			\end{aligned}
		\end{equation*}
		where
		
		\begin{equation}\nonumber
			\begin{aligned}
				I_2&=\bigg(\frac{2}{3}\partial_{r}(r^2\partial_{rr}v\partial_{rr}\tilde{S}_1)
				+\frac{2}{3}\partial_{r}(2r\partial_{r}v\partial_{rr}\tilde{S}_1)
				-\frac{4}{3}\partial_{r}(v\partial_{rr}\tilde{S}_1)
				+\frac{2}{3}\partial_{r}(2r\partial_{rr}v\partial_{r}\tilde{S}_1)
				+\frac{2}{3}\partial_{r}(4\partial_{r}v\partial_{r}\tilde{S}_1)\\
				&\quad\quad-\frac{4}{3}\partial_{r} (\frac{2v}{r}\partial_{r}\tilde{S}_1)
				-\frac{2}{3}\partial_{r}(6\partial_{rr}v\tilde{S}_1)
				-\frac{2}{3}\partial_{r}(\frac{12}{r}\partial_{r}v\tilde{S}_1)
				+\frac{4}{3}\partial_{r}( \frac{6v}{r^2}\tilde{S}_1)\bigg),\\
				I_3&=\bigg(\partial_{r}(r^2\partial_{rr}v\partial_{rr}\tilde{S}_2)+\partial_{r}(2r\partial_{rr}v\partial_{r}\tilde{S}_2)++\partial_{r}(2r\partial_{r}v\partial_{rr}\tilde{S}_2)+\partial_{r}(4\partial_{r}v\partial_{r}\tilde{S}_2)-2\partial_{r}(v\partial_{rr}\tilde{S}_2)\\
				&\quad\quad-2\partial_{r}(\frac{2v}{r}\partial_{r}\tilde{S}_2)\bigg).
			\end{aligned}
		\end{equation}
		
		Furthermore, we get
			\begin{align*}
				&\int_{1}^{+\infty}H_4\mathrm{~d}r\\
				=&\int_{1}^{+\infty}
				\bigg(-\tau\epsilon\rho\frac{12}{\mu r^2}\tilde{S}_1\partial_{r}\tilde{S}_1
				+\tau\epsilon\rho\frac{2}{3\mu}\partial_{r}\tilde{S}_1\partial_{rr}\tilde{S}_1
				+\tau\epsilon\rho\frac{4}{\mu r}(\partial_{r}\tilde{S}_1)^2
				-\tau\epsilon\rho\frac{2r}{3\mu}(\partial_{rr}\tilde{S}_1)^2
				+\tau\epsilon\rho\frac{4}{\mu r}\tilde{S}_1\partial_{rr}\tilde{S}_1\\
				&-\tau\epsilon\rho\frac{2r}{3\mu}\partial_{r}\tilde{S}_1\partial_{rrr}\tilde{S}_1
				+\tau\epsilon\rho\frac{2}{\mu}\tilde{S}_1\partial_{rrr}\tilde{S}_1
				-\tau\epsilon\rho\frac{4}{\lambda}\partial_{r}\tilde{S}_2\partial_{rr}\tilde{S}_2
				-\tau\epsilon\rho\frac{2r}{\lambda}(\partial_{rr}\tilde{S}_2)^2
				-\tau\epsilon\rho\frac{2r}{\lambda}\partial_{r}\tilde{S}_2\partial_{rrr}\tilde{S}_2\bigg)\mathrm{d}r\\
				\geq& C\epsilon\int_{1}^{+\infty}
				-(\tilde{S}_1^2+(\partial_{r}\tilde{S}_1)^2
				+(\partial_{rr}\tilde{S}_1)^2+(\partial_{r}\tilde{S}_2)^2
				+(\partial_{rr}\tilde{S}_2)^2)\mathrm{d}r\mathrm{d}t
				+\frac{2\tau\epsilon}{3\mu}(\rho\partial_{r}\tilde{S}_1\partial_{rr}\tilde{S}_1)(t,1)\\
				&
				-\frac{2\tau\epsilon}{\mu}(\rho\tilde{S}_1\partial_{rr}\tilde{S}_1)(t,1)
				+\frac{2\tau\epsilon}{\lambda}(\rho\partial_{r}\tilde{S}_2\partial_{rr}\tilde{S}_2)(t,1)
				-C\left(E(0)+ E^{\frac{1}{2}}(t) \int_{0}^{t}\mathcal{D}(s)\mathrm{d}s\right) \displaybreak \\
				\geq&-C\epsilon\int_{1}^{+\infty}\left((\partial_{rr}\tilde{S}_1)^2+(\partial_{rr}\tilde{S}_2)^2\right)\mathrm{d}r
				-\frac{4\tau\epsilon}{3\mu}\eta\rho(t,1)(\partial_{rr}\tilde{S}_1)^2(t,1)
				-\frac{2\tau\epsilon}{\lambda}\eta\rho(t,1)(\partial_{rr}\tilde{S}_2)^2(t,1)\\
				&-\frac{2\tau\epsilon}{\mu} C(\eta)((\tilde{S}_1)^2+(\partial_{r}\tilde{S}_1)^2)(t,1)
				-\frac{2\tau\epsilon}{\lambda} C(\eta)((\partial_{r}\tilde{S}_2)^2)(t,1)
				-C\left(E(0)+ E^{\frac{1}{2}}(t)  \int_{0}^{t}\mathcal{D}(s)\mathrm{d}s
				\right).
			\end{align*}
		Here and after, $\eta$ is any positive constant to be chosen later.	
		
		By adding (\ref{64}) to (\ref{61}) and integrating over $(0,t)$, we derive that
		\begin{equation}\label{65}
			\begin{aligned}
				&\int_{1}^{+\infty}\bigg(\frac{r^2 P^{\prime}(\rho)}{2\rho} \left(\partial_{rr} \rho\right)^2+\frac{ P^{\prime}(\rho)}{\rho}2(\partial_{r}\rho)^2+\frac{P^{\prime}(\rho)}{\rho}2r\partial_{r}\rho\partial_{rr}\rho\\%\ρ相关
				&\quad\quad
				+\frac{r^2\rho}{2} (\partial_{rr}v)^2+2\rho(\partial_{r}v)^2+\frac{2\rho}{r^2}v^2+2r\rho\partial_{rr}v\partial_{r}v-2\rho\partial_{rr}v v-\frac{4\rho v}{r}\partial_{r}v\\%v相关
				&\quad\quad
				+\tau\rho(\frac{r^2}{6\mu}(\partial_{rr}\tilde{S}_1)^2+\frac{2}{3\mu}(\partial_{r}\tilde{S}_1)^2+\frac{6}{\mu r^2}(\tilde{S}_1)^2 +\frac{2r}{3\mu}\partial_{rr}\tilde{S}_1 \partial_{r}\tilde{S}_1-\frac{2}{\mu}\partial_{rr}\tilde{S}_1\tilde{S}_1-\frac{4}{\mu r}\partial_{r}\tilde{S}_1 \tilde{S}_1\\%S_1相关
				&\quad\quad
				+\frac{r^2}{2\lambda}(\partial_{rr} \tilde{S}_2)^2+\frac{2}{\lambda}(\partial_{r} \tilde{S}_2)^2+\frac{2r}{\lambda}\partial_{rr} \tilde{S}_2\partial_{r} \tilde{S}_2)\bigg)\mathrm{d}r
				\\%S_2相关
				&+\int_{0}^{t}\int_{1}^{+\infty}\bigg(\frac{r^2}{3\mu}(\partial_{rr}\tilde{S}_1)^2+\frac{4}{3\mu}(\partial_{r}\tilde{S}_1)^2+\frac{12}{\mu r^2}(\tilde{S}_1)^2 +\frac{4r}{3\mu}\partial_{rr}\tilde{S}_1 \partial_{r}\tilde{S}_1-\frac{4}{\mu}\partial_{rr}\tilde{S}_1\tilde{S}_1-\frac{8}{\mu r}\partial_{r}\tilde{S}_1 \tilde{S}_1\\
				&\quad\quad
				+\frac{r^2}{\lambda}(\partial_{rr}\tilde{S}_2)^2+\frac{4}{\lambda}(\partial_{r}  \tilde{S}_2)^2+\frac{4r}{\lambda}\partial_{rr} \tilde{S}_2\partial_{r} \tilde{S}_2\bigg)\mathrm{d}r\mathrm{d}t\\
				&+\int_{0}^{t}\left(\frac{\tau\epsilon}{3\mu}\rho(t,1)(\partial_{rr}\tilde{S}_1)^2(t,1)
				+\frac{3\tau\epsilon}{4\lambda}\rho(t,1)(\partial_{rr}\tilde{S}_2)^2(t,1)\right)\mathrm{d}t\\ 
				&\leq \int_{0}^{t}\int_{1}^{+\infty} \left(I_1+I_2+I_3\right)\mathrm{d}r\mathrm{d}t+C \left(E(0)+E^{\frac{3}{2}}(t)+E^{\frac{1}{2}}(t) \int_{0}^{t}\mathcal{D}(s)\mathrm{d}s
				\right).
			\end{aligned}
		\end{equation}
		
		Next, we address the first term on the right-hand side of equation (\ref{65}). By applying the momentum equation and equations $(\ref{22})_1$ and $(\ref{22})_4$, we have
		
			\begin{align*}
				&\int_{1}^{+\infty}(I_1+I_2+I_3)\mathrm{d}r\\
				=&\bigg(\frac{2}{3}r^2\partial_{rr}v\partial_{rr}\tilde{S}_1
				+\frac{4}{3}r\partial_{r}v\partial_{rr}\tilde{S}_1
				-\frac{4}{3}v\partial_{rr}\tilde{S}_1
				+\frac{4}{3}r\partial_{rr}v\partial_{r}\tilde{S}_1
				+\frac{8}{3}\partial_{r}v\partial_{r}\tilde{S}_1
				-\frac{8}{3}\frac{v}{r}\partial_{r}\tilde{S}_1\\
				&\qquad
				-4\partial_{rr}v\tilde{S}_1-\frac{8}{r}\partial_{r}v\tilde{S}_1
				+\frac{8}{r^2}v\tilde{S}_1\bigg)\bigg|_{1}^{+\infty}\\
				&+
				\bigg(r^2\partial_{rr}v\partial_{rr}\tilde{S}_2+2r\partial_{rr}v\partial_{r}\tilde{S}_2+2r\partial_{r}v\partial_{rr}\tilde{S}_2+4\partial_{r}v\partial_{r}\tilde{S}_2-2v\partial_{rr}\tilde{S}_2-\frac{4v}{r}\partial_{r}\tilde{S}_2\bigg)\bigg|_{1}^{+\infty}\\
				&-\bigg(r^2\partial_{rr}P\partial_{rr}v +2r\partial_{r}P\partial_{rr}v+2r\partial_{rr}P\partial_{r}v+
				4\partial_{r}v\partial_{r}P
				-2v\partial_{rr}P-\frac{4}{r}v\partial_{r}P\bigg)\bigg|_{1}^{+\infty}\\
				=&r^2\partial_{rr}v\left(\dfrac{2}{3}\partial_{rr} \tilde{S}_1+\dfrac{2}{r} \partial_r \tilde{S}_1-\dfrac{2}{r^2}\tilde{S}_1+\partial_{rr} \tilde{S}_2-\partial_{rr}P\right)\bigg|_{1}^{+\infty}\\
				&+2r\partial_{r}v\left(\dfrac{2}{3}\partial_{rr} \tilde{S}_1+\dfrac{2}{r} \partial_r \tilde{S}_1-\dfrac{2}{r^2}\tilde{S}_1+\partial_{rr} \tilde{S}_2-\partial_{rr}P\right)\bigg|_{1}^{+\infty}\\
				&-r\partial_{rr}v\left(\dfrac{2}{3}\partial_r \tilde{S}_1+\dfrac{2}{r} \tilde{S}_1+\partial_r \tilde{S}_2-\partial_{r}P\right)\bigg|_{1}^{+\infty}
				-2\partial_{r}v\left(\dfrac{2}{3}\partial_r \tilde{S}_1+\dfrac{2}{r} \tilde{S}_1+\partial_r \tilde{S}_2-\partial_{r}P\right)\bigg|_{1}^{+\infty}\\
				&-2v\left(\dfrac{2}{3}\partial_{rr} \tilde{S}_1+\dfrac{2}{r} \partial_r \tilde{S}_1-\dfrac{2}{r^2}\tilde{S}_1+\partial_{rr} \tilde{S}_2-\partial_{rr}P\right)\bigg|_{1}^{+\infty}
				+\frac{2}{r}v\left(\dfrac{2}{3}\partial_r \tilde{S}_1+\dfrac{2}{r} \tilde{S}_1+\partial_r \tilde{S}_2-\partial_{r}P\right)\bigg|_{1}^{+\infty}\\
				&+2r\frac{\partial_{r}P}{\rho}(\partial_{tr}\rho+2\partial_{r}\rho\partial_{r}v)\bigg|_{1}^{+\infty}
				+2r\partial_{r}\tilde{S}_2\left(\frac{\tau\rho}{\lambda}\left(\partial_{tr} \tilde{S}_2+(v-\epsilon)\partial_{rr} \tilde{S}_2\right)+\partial_{r} \tilde{S}_2\right)\bigg|_{1}^{+\infty}.
			\end{align*}

		Furthermore, using momentum equation, $(\ref{22})_4$ and the boundary condition (\ref{1.6}), we have 
		\begin{equation*}
			\begin{aligned}
				&\int_{0}^{t}\int_{1}^{+\infty}(I_1+I_2+I_3)\mathrm{d}r\mathrm{d}t\\
				=& \int_{0}^{t}\bigg[\left((r^2\partial_{rr}v+2r\partial_{r}v-2v)(\rho\partial_{tr}v+\rho v\partial_{rr}v) +(-r\partial_{rr}v-2\partial_{r}v+\frac{2v}{r})
				(\rho\partial_{t}v+\rho v \partial_{r}v)\right)\bigg|_{1}^{+\infty}\\
				&\qquad+2r\frac{\partial_{r}P}{\rho}(\partial_{tr}\rho+2\partial_{r}\rho\partial_{r}v)\bigg|_{1}^{+\infty}
				+2r\partial_{r}\tilde{S}_2\left(\frac{\tau\rho}{\lambda}\left(\partial_{tr} \tilde{S}_2+(v-\epsilon)\partial_{rr} \tilde{S}_2\right)+\partial_{r} \tilde{S}_2\right)\bigg|_{1}^{+\infty}\bigg]\mathrm{d}t\\
				=&\int_{0}^{t}
				\bigg[-(r^2\partial_{rr}v+2r\partial_{r}v)\rho\partial_{tr}v(t,1)
				-2r\frac{\partial_{r}P}{\rho}(\partial_{tr}\rho+2\partial_{r}\rho\partial_{r}v)(t,1)\\
				&\qquad-2r\partial_{r}\tilde{S}_2\left(\frac{\tau\rho}{\lambda}\left(\partial_{tr} \tilde{S}_2+(v-\epsilon)\partial_{rr} \tilde{S}_2\right)+\frac{1}{\lambda}\partial_{r} \tilde{S}_2\right)(t,1)\bigg]
				\mathrm{d}t
			\end{aligned}
		\end{equation*}
	Next, we estimate each term in the above equation.
		\begin{equation}\nonumber
			\begin{aligned}
				&\int_{0}^{t}-(r^2\partial_{rr}v+2r\partial_{r}v)\rho\partial_{tr}v(t,1)\mathrm{d}t
				=-\int_{0}^{t}(\rho\partial_{rr}v\partial_{tr}v)(t,1)\mathrm{d}t- (\rho(\partial_{r}v)^2)(t,1)\bigg|_0^t\\
				&
				\leq C\int_{0}^{t}|\partial_{rr}v\partial_{tr}v(t,1)|\mathrm{d}t+E(0)
				\leq C\int_{0}^{t}((\partial_{rr}v)^2(t,1)+(\partial_{tr}v)^2(t,1))\mathrm{d}t+E(0),
			\end{aligned}
		\end{equation}
		$$
		\begin{aligned}
			&\int_{0}^{t}-2r\frac{\partial_{r}P}{\rho}(\partial_{tr}\rho+2\partial_{r}\rho\partial_{r}v)(t,1)\mathrm{d}t\\
			=& -\frac{P^{\prime}(\rho)}{\rho}(\partial_{r}\rho)^2(t,1)\bigg|_0^t
			+\int_{0}^{t}[\frac{P^{\prime}(\rho)}{\rho}]_t(\partial_{r}\rho)^2(t,1)\mathrm{d}t-\int_{0}^{t}\frac{4\partial_{r}P}{\rho}\partial_{r}\rho\partial_{r}v(t,1)\mathrm{d}t\\
			\leq&-\frac{P^{\prime}(\rho)}{\rho}(\partial_{r}\rho)^2(t,1)
			+\int_{0}^{t}C\|\partial_{t}\rho(\partial_{r}\rho)^2\|_{L^{\infty}}\mathrm{d}t
			-\int_{0}^{t}C\|\partial_{r}v(\partial_{r}\rho)^2\|_{L^{\infty}}\mathrm{d}t\\
			\leq&\int_{0}^{t}C\|\partial_{t}\rho\|_{H^{1}}\|\partial_{r}\rho\|_{H^{1}}^2\mathrm{d}t
			-\int_{0}^{t}C\|\partial_{r}v\|_{H^{1}}\|\partial_{r}\rho\|_{H^{1}}^2\mathrm{d}t
			\leq C \left(E(0)+E^{\frac{1}{2}}(t) \int_{0}^{t}\mathcal{D}(s)\mathrm{d}s\right),
		\end{aligned}
		$$
		and
		$$
		\begin{aligned}
			&\int_{0}^{t}-2r\partial_{r}\tilde{S}_2\left(\frac{\tau\rho}{\lambda}\left(\partial_{tr} \tilde{S}_2+(v-\epsilon)\partial_{rr} \tilde{S}_2\right)+\frac{1}{\lambda}\partial_{r} \tilde{S}_2\right)(t,1)\mathrm{d}t\\
			=&-\int_{0}^{t}\frac{\tau}{\lambda}\rho(t,1)(\partial_{r}\tilde{S}_2)^2_t(t,1)\mathrm{d}t
			+\int_{0}^{t}\frac{2\epsilon\tau }{\lambda}\rho(t,1)\partial_{r}\tilde{S}_2(t,1)\partial_{rr}\tilde{S}_2(t,1)\mathrm{d}t
			-\int_{0}^{t}\frac{2}{\lambda}(\partial_{r}\tilde{S}_2)^2(t,1)\mathrm{d}t\\
			\leq&-\frac{\tau}{\lambda}\rho(t,1)(\partial_{r}\tilde{S}_2)^2(t,1)\bigg|_0^t
			+\int_{0}^{t}\frac{2\tau}{\lambda}\partial_{t}\rho(\partial_{r}\tilde{S}_2)^2(t,1)\mathrm{d}t
			+\int_{0}^{t}\frac{2\epsilon\tau }{\lambda}\rho(\eta(\partial_{rr}\tilde{S}_2)^2
			+C(\eta)(\partial_{r}\tilde{S}_2)^2)(t,1)\mathrm{d}t\\
			\leq& \int_{0}^{t}\frac{2\epsilon\tau }{\lambda}\eta\rho(t,1)(\partial_{rr}\tilde{S}_2)^2(t,1)
			+C \left(E(0)+E^{\frac{1}{2}}(t) \int_{0}^{t}\mathcal{D}(s)\mathrm{d}s\right).
		\end{aligned}
		$$
		Here and after, $\eta$ is any positive constant to be chosen later.
		
		Combining the above results and using \eqref{hu3.4}, we can obtain
		\begin{equation}\label{3.85}
			\begin{aligned}
				&\int_{1}^{+\infty}\bigg(\frac{r^2 P^{\prime}(\rho)}{2\rho} \left(\partial_{rr} \rho\right)^2+\frac{2 P^{\prime}(\rho)}{\rho}(\partial_{r}\rho)^2+\frac{2rP^{\prime}(\rho)}{\rho}\partial_{r}\rho\partial_{rr}\rho\\%\ρ相关
				&\qquad\qquad+\frac{r^2\rho}{2} (\partial_{rr}v)^2+2\rho(\partial_{r}v)^2+\frac{2\rho}{r^2}v^2+2r\rho\partial_{rr}v\partial_{r}v-2\rho\partial_{rr}v v-\frac{4\rho v}{r}\partial_{r}v\\%v相关
				&\qquad\qquad+\tau\rho(\frac{r^2}{6\mu}(\partial_{rr}\tilde{S}_1)^2+\frac{2}{3\mu}(\partial_{r}\tilde{S}_1)^2+\frac{6}{\mu r^2}\tilde{S}_1^2 +\frac{2r}{3\mu}\partial_{rr}\tilde{S}_1 \partial_{r}\tilde{S}_1-\frac{2}{\mu}\partial_{rr}\tilde{S}_1\tilde{S}_1\\%S_1相关
				&\qquad\qquad-\frac{4}{\mu r}\partial_{r}\tilde{S}_1 \tilde{S}_1 +\frac{r^2}{2\lambda}(\partial_{rr} \tilde{S}_2)^2+\frac{2}{\lambda}(\partial_{r} \tilde{S}_2)^2+\frac{2r}{\lambda}\partial_{rr} \tilde{S}_2\partial_{r} \tilde{S}_2)\bigg)\mathrm{d}r\\%S_2相关
				&+\int_0^t\int_{1}^{+\infty}\bigg(\frac{r^2}{3\mu}(\partial_{rr}\tilde{S}_1)^2+\frac{4}{3\mu}(\partial_{r}\tilde{S}_1)^2+\frac{12}{\mu r^2}(\tilde{S}_1)^2 +\frac{4r}{3\mu}\partial_{rr}\tilde{S}_1 \partial_{r}\tilde{S}_1-\frac{4}{\mu}\partial_{rr}\tilde{S}_1\tilde{S}_1\\
				&\qquad\qquad\qquad-\frac{8}{\mu r}\partial_{r}\tilde{S}_1 \tilde{S}_1
				+\frac{r^2}{\lambda}(\partial_{rr}\tilde{S}_2)^2+\frac{4}{\lambda}(\partial_{r}  \tilde{S}_2)^2 +\frac{4r}{\lambda}\partial_{rr} \tilde{S}_2\partial_{r} \tilde{S}_2\bigg)\mathrm{d}r\mathrm{d}t\\
				&+\int_{0}^{t}\left(\frac{3\tau\epsilon}{16\mu}(\partial_{rr}\tilde{S}_1)^2(t,1)
				+\frac{9\tau\epsilon}{32\lambda}(\partial_{rr}\tilde{S}_2)^2(t,1)\right)\mathrm{d}t\\
				&\leq C \left(E(0)+E^{\frac{3}{2}}(t)+E^{\frac{1}{2}}(t) \int_{0}^{t}\mathcal{D}(s)\mathrm{d}s
				+\int_{0}^{t}((\partial_{rr}v)^2(t,1)+(\partial_{tr}v)^2(t,1))\mathrm{d}t\right).
				\\ 
			\end{aligned}
		\end{equation}
		Further, we observe that 
		\begin{equation}\nonumber
			\begin{aligned}
				&\int_{1}^{+\infty}(\frac{r^2\rho}{2} (\partial_{rr}v)^2+2\rho(\partial_{r}v)^2+\frac{2\rho}{r^2}v^2+2r\rho\partial_{rr}v\partial_{r}v-2\rho\partial_{rr}v v-\frac{4\rho v}{r}\partial_{r}v)\mathrm{d}r\\
				=&\int_{1}^{+\infty}(\frac{r^2\rho}{2} (\partial_{rr}v)^2+4\rho(\partial_{r}v)^2
				+2r\rho\partial_{rr}v\partial_{r}v
				+2\partial_{r}\rho\partial_{r}v v+\frac{2}{r}\partial_{r}\rho v^2)\mathrm{d}r\\
				\geq&\int_{1}^{+\infty}(\frac{r^2\rho}{2} (\partial_{rr}v)^2+4\rho(\partial_{r}v)^2
				+2r\rho\partial_{rr}v\partial_{r}v)\mathrm{d}r-CE^{\frac{1}{2}}(t) \mathcal{D}(t),
			\end{aligned}
		\end{equation}
		where we used the boundary condition (\ref{1.6}). 
		Moreover, we know
		\begin{equation}\nonumber
			\begin{aligned}
				&\int_{1}^{+\infty}(\frac{r^2\rho}{2} (\partial_{rr}v)^2+4\rho(\partial_{r}v)^2
				+2r\rho\partial_{rr}v\partial_{r}v)\mathrm{d}r\\
				\geq&\int_{1}^{+\infty}(\frac{r^2\rho}{2} (\partial_{rr}v)^2+4\rho(\partial_{r}v)^2
				-\frac{3r^2\rho}{8}(\partial_{rr}v)^2-\frac{8}{3}\rho(\partial_{r}v)^2)\mathrm{d}r
				=\int_{1}^{+\infty}(\frac{r^2\rho}{8} (\partial_{rr}v)^2+\frac{4}{3}\rho(\partial_{r}v)^2)\mathrm{d}r.
			\end{aligned}
		\end{equation}
		Similarly, we have 
		\begin{equation}\nonumber
			\begin{aligned}
				&\int_{1}^{+\infty}(\frac{r^2}{6\mu}(\partial_{rr}\tilde{S}_1)^2+\frac{2}{3\mu}(\partial_{r}\tilde{S}_1)^2+\frac{6}{\mu r^2}(\tilde{S}_1)^2 +\frac{2r}{3\mu}\partial_{rr}\tilde{S}_1 \partial_{r}\tilde{S}_1-\frac{2}{\mu}\partial_{rr}\tilde{S}_1\tilde{S}_1-\frac{4}{\mu r}\partial_{r}\tilde{S}_1 \tilde{S}_1)\mathrm{d}r\\
				\geq&\int_{1}^{+\infty}(\frac{r^2}{24\mu}(\partial_{rr}\tilde{S}_1)^2+\frac{16}{9\mu}(\partial_{r}\tilde{S}_1)^2+\frac{4}{\mu r^2}(\tilde{S}_1)^2 
				)\mathrm{d}r-CE^{\frac{1}{2}}(t) \mathcal{D}(t).
			\end{aligned}
		\end{equation}
		
		Thus, equation (\ref{3.85}) can be simplified to 
		\begin{equation}\label{3.87}
			\begin{aligned}
				&\int_{1}^{+\infty}\bigg(\frac{ P^{\prime}(\rho)}{2\rho}\left(r\partial_{rr} \rho+r\partial_{r} \rho\right)^2
				+\frac{r^2\rho}{8} (\partial_{rr}v)^2+\frac{4}{3}\rho(\partial_{r}v)^2
				\\
				&\qquad\qquad+\tau\rho\left(\frac{r^2}{24\mu}(\partial_{rr}\tilde{S}_1)^2+\frac{16}{9\mu}(\partial_{r}\tilde{S}_1)^2+\frac{4}{\mu r^2}(\tilde{S}_1)^2\right)+\frac{\tau\rho}{2\lambda}\left(r\partial_{rr} \tilde{S}_2+2\partial_{r} \tilde{S}_2\right)^2\bigg)\mathrm{d}r\\%S_2相关
				&+\int_0^t\int_{1}^{+\infty}\bigg(\frac{r^2}{12\mu}(\partial_{rr}\tilde{S}_1)^2+\frac{32}{9\mu}(\partial_{r}\tilde{S}_1)^2+\frac{8}{\mu r^2}(\tilde{S}_1)^2
				+\frac{1}{\lambda}\left(r\partial_{rr} \tilde{S}_2+2\partial_{r} \tilde{S}_2\right)^2\bigg)\mathrm{d}r\mathrm{d}t\\
				&
				+\int_{0}^{t}\left(\frac{\tau\epsilon}{4\mu}(\partial_{rr}\tilde{S}_1)^2(t,1)
				+\frac{3\tau\epsilon}{8\lambda}(\partial_{rr}\tilde{S}_2)^2(t,1)\right)\mathrm{d}t\\
				&\leq C \left(E(0)+E^{\frac{3}{2}}(t)+E^{\frac{1}{2}}(t) \int_{0}^{t}\mathcal{D}(s)\mathrm{d}s
				+\int_{0}^{t}((\partial_{rr}v)^2(t,1)+(\partial_{tr}v)^2(t,1))\mathrm{d}t\right).
				\\ 
			\end{aligned}
		\end{equation}
		Then, we can get the desired result.
	\end{proof}    
	
	To address the boundary term on the right-hand side of \eqref{3.87}, we now establish the following key lemmas.
	
	First, we give estimates of $\partial_{tr} v$ on the boundary.
	\begin{lemma}\label{lem7_tx}
		There exists a constant $C$ such that for any $0\leq t \leq T$
		\begin{equation}\label{bdy_tx}
			\begin{aligned}
				\int_{0}^{t}(\partial_{tr} v)^2(t,1)\mathrm{d}t
				\leq C\left(E(0)+E^{\frac{3}{2}}(t)+E^{\frac{1}{2}}(t) \int_{0}^{t}\mathcal{D}(s)\mathrm{d}s\right).
			\end{aligned}
		\end{equation}
		
	\end{lemma}
	\begin{proof}
		Multiplying  $(\ref{47})_3$ by $\dfrac{r^2}{3\mu}\partial_{tr} v$  and integrating the result, we get
		\begin{align}\label{76}
				&\int_{0}^{t}\int_{1}^{+\infty}\frac{r^2\tau}{3\mu}\rho\partial_{tr} v\partial_{ttr} \tilde{S}_1\mathrm{d}r\mathrm{d}t
				+\int_{0}^{t}\int_{1}^{+\infty}\frac{r^2\tau}{3\mu}\rho (v-\epsilon)\partial_{tr} v\partial_{trr} \tilde{S}_1\mathrm{d}r\mathrm{d}t
				\nonumber \\
				&+\int_{0}^{t}\int_{1}^{+\infty}\frac{r^2}{3\mu}\partial_{tr} v\partial_{tr} \tilde{S}_1\mathrm{d}r\mathrm{d}t-\int_{0}^{t}\int_{1}^{+\infty}\frac{2}{3}\left(r^2\partial_{trr}v \partial_{tr}v-r(\partial_{tr}v)^2 +\partial_{t}v\partial_{tr}v\right)\mathrm{d}r\mathrm{d}t \nonumber \\
				\leq& C\left(E(0)+ E^{\frac{1}{2}}(t) \int_{0}^{t}\mathcal{D}(s)\mathrm{d}s\right)
				.
		\end{align}	
		We estimate each term of the above equation, respectively. Firstly, using the momentum equation $(\ref{15})_2$, we have
		\begin{align}\label{77}
			%\begin{aligned}
				&\int_{0}^{t}\int_{1}^{+\infty}\dfrac{r^2\tau}{3\mu}\rho\partial_{tr} v\partial_{ttr} \tilde{S}_1\mathrm{d}r\mathrm{d}t \nonumber \\
				=&
				\int_{1}^{+\infty}\dfrac{r^2\tau}{3\mu}\rho\partial_{tr} v \partial_{tr} \tilde{S}_1\big|_0^{t}\mathrm{d}r-
				\int_{0}^{t}\int_{1}^{+\infty}\dfrac{r^2\tau}{3\mu}\rho\partial_{ttr} v \partial_{tr} \tilde{S}_1\mathrm{d}r\mathrm{d}t-
				\int_{0}^{t}\int_{1}^{+\infty}\dfrac{r^2\tau}{3\mu}\partial_{t}\rho
				\partial_{tr} v \partial_{tr} \tilde{S}_1\mathrm{d}r\mathrm{d}t \nonumber \\
				\geq& \int_{1}^{+\infty}\dfrac{r^2\tau}{3\mu}\rho\partial_{tr} v\partial_{tr} \tilde{S}_1\mathrm{d}r-
				\int_{0}^{t}\int_{1}^{+\infty}\dfrac{r^2\tau}{2\mu}\rho\partial_{ttr} v(\rho\partial_{tt} v+\rho v\partial_{tr} v+\partial_{tr} P-\dfrac{2}{r} \partial_t \tilde{S}_1
				-\partial_{tr} \tilde{S}_2)\mathrm{d}r\mathrm{d}t \nonumber \\
				&-C\left(E(0)+ E^{\frac{1}{2}}(t) \int_{0}^{t}\mathcal{D}(s)\mathrm{d}s\right)
			\end{align}
		%\end{equation}
		where 
		\begin{equation}\nonumber
			\begin{aligned}
				\int_{1}^{+\infty}\dfrac{r^2\tau}{3\mu}\rho\partial_{tr} v\partial_{tr} \tilde{S}_1\mathrm{d}r\geq&-C\int_{1}^{+\infty}r^2\rho ((\partial_{tr} v)^2+\tau(\partial_{tr} \tilde{S}_1)^2)\mathrm{d}r\\
				\geq&-C\left(E(0)+ E^{\frac{3}{2}}(t)+ E^{\frac{1}{2}}(t) \int_{0}^{t}\mathcal{D}(s)\mathrm{d}s\right),
			\end{aligned}
		\end{equation}
		\begin{equation}\nonumber
			\begin{aligned}
				&-\int_{0}^{t}\int_{1}^{+\infty}\dfrac{r^2\tau}{2\mu}\rho^2\partial_{ttr} v\partial_{tt} v\mathrm{d}r\mathrm{d}t\\
				=&\int_{0}^{t}\int_{1}^{+\infty}\dfrac{\tau}{4\mu}(2r\rho^2+2r^2\rho
				\partial_{r}\rho)(\partial_{tt}v)^2\mathrm{d}r\mathrm{d}t
				\geq-C\left(E(0)+ E^{\frac{1}{2}}(t) \int_{0}^{t}\mathcal{D}(s)\mathrm{d}s\right),
			\end{aligned}
		\end{equation}
		\begin{align*}
		&-\int_{0}^{t}\int_{1}^{+\infty}\dfrac{r^2\tau}{2\mu}\rho^2v\partial_{ttr} v\partial_{tr} v\mathrm{d}r\mathrm{d}t\\
		=&-\int_{1}^{+\infty}\dfrac{\tau}{4\mu}r^2\rho^2v(\partial_{tr} v)^2\big|_0^t\mathrm{d}r 
		+\int_{0}^{t}\int_{1}^{+\infty}\dfrac{r^2\tau}{4\mu}\partial_{t}(\rho^2v)
		(\partial_{tr} v)^2\mathrm{d}r\mathrm{d}t\\
		\geq&-C\left(E(0)+ E^{\frac{3}{2}}(t)+ E^{\frac{1}{2}}(t) \int_{0}^{t}\mathcal{D}(s)\mathrm{d}s\right),
		\end{align*}
		and 
			\begin{align*}
				&-\int_{0}^{t}\int_{1}^{+\infty}\dfrac{\tau r^2}{2\mu}\rho\partial_{ttr} v\partial_{tr} P\mathrm{d}r\mathrm{d}t\\
				=&-\int_{1}^{+\infty}\dfrac{\tau r^2} {2\mu} \rho\partial_{tr} v\partial_{tr} P\big|_0^t\mathrm{d}r 
				+\int_{0}^{t}\int_{1}^{+\infty}\dfrac{\tau r^2}{2\mu}\rho\partial_{tr} v \partial_{ttr} P \mathrm{d}r \mathrm{d}t
				+\int_{0}^{t}\int_{1}^{+\infty}\dfrac{\tau r^2}{2\mu}\partial_{tr} v \partial_{t}\rho\partial_{tr} P \mathrm{d}r\mathrm{d}t\\
				\geq&-C\int_{1}^{+\infty} r^2((\partial_{tr} v)^2+(\partial_{tr} \rho)^2)\mathrm{d}r +\int_{0}^{t}\int_{1}^{+\infty}\dfrac{\tau r^2}{2\mu} \rho P^{\prime}(\rho) \partial_{ttr}\rho \partial_{tr} v\mathrm{d}r\mathrm{d}t\\&
				-C\left(E(0)+ E^{\frac{1}{2}}(t) \int_{0}^{t}\mathcal{D}(s)\mathrm{d}s\right)\\
				\geq&-\int_{0}^{t}\int_{1}^{+\infty}\dfrac{\tau r^2}{2\mu} \rho P^{\prime}(\rho)
				\partial_{tr}v(\partial_{trr}(\rho v)+\dfrac{2}{r} \partial_{tr}(\rho v)-\dfrac{2}{r^2} \partial_{t}(\rho v))\mathrm{d}r\mathrm{d}t\\
				&-C\left(E(0)+ E^{\frac{3}{2}}(t)+ E^{\frac{1}{2}}(t) \int_{0}^{t}\mathcal{D}(s)\mathrm{d}s\right)\\
				\geq&
				-\int_{0}^{t}\int_{1}^{+\infty}\dfrac{\tau r^2}{2\mu}\rho P^{\prime}(\rho)
				v\partial_{tr}v\partial_{trr}\rho\mathrm{d}r\mathrm{d}t
				-C\left(E(0)+ E^{\frac{3}{2}}(t)+ E^{\frac{1}{2}}(t) \int_{0}^{t}\mathcal{D}(s)\mathrm{d}s\right).
			\end{align*}
		
		Note that there is still a three-order term $\int_{0}^{t}\int_{1}^{+\infty}\frac{\tau r^2}{2\mu}\rho P^{\prime}(\rho) v\partial_{tr}v \partial_{trr}\rho \mathrm{d}r\mathrm{d}t$ in the above inequality, which can be canceled out by the convective term, as shown in \eqref{103}. 
		
	Next, we handle the remaining two terms in inequality \eqref{77}. Applying Lemma \ref{lem2} and Lemmas \ref{lem3}-\ref{lem5}, we have
		
		\begin{equation}\nonumber
			\begin{aligned}
				&\int_{0}^{t}\int_{1}^{+\infty}\dfrac{\tau r}{\mu}\rho \partial_{ttr}v \partial_{t}\tilde{S}_1\mathrm{d}r\mathrm{d}t\\
				=&
				\int_{1}^{+\infty}\dfrac{\tau r}{\mu}\rho \partial_{tr}v \partial_{t}\tilde{S}_1\big|_0^{t}\mathrm{d}r
				-\int_{0}^{t}\int_{1}^{+\infty}\dfrac{\tau r}{\mu}\rho \partial_{tr}v \partial_{tt}\tilde{S}_1\mathrm{d}r\mathrm{d}t
				-\int_{0}^{t}\int_{1}^{+\infty}\dfrac{\tau r}{\mu}\partial_{t}\rho \partial_{tr}v \partial_{t}\tilde{S}_1\mathrm{d}r\mathrm{d}t\\
				\geq&
				-C\int_{1}^{+\infty}r^2\rho((\partial_{tr}v)^2+(\partial_{t}\tilde{S}_1)^2)\mathrm{d}r
				-C\int_{0}^{t}\int_{1}^{+\infty}r^2\rho((\partial_{tr}v)^2+(\partial_{tt}\tilde{S}_1)^2)\mathrm{d}r\mathrm{d}t\\&
				-C\left(E(0)+ E^{\frac{1}{2}}(t) \int_{0}^{t}\mathcal{D}(s)\mathrm{d}s\right)\\
				\geq&-C\left(E(0)+ E^{\frac{3}{2}}(t)+ E^{\frac{1}{2}}(t) \int_{0}^{t}\mathcal{D}(s)\mathrm{d}s\right).
			\end{aligned}
		\end{equation}
		In view of equation $(\ref{47})_4$ and Lemmas \ref{lem3.4}, \ref{lem4} and \ref{lem5}, we have	
		\begin{equation}\nonumber
			\begin{aligned}
				&\int_{0}^{t}\int_{1}^{+\infty}\dfrac{\tau r^2}{2\mu}\rho \partial_{ttr}v \partial_{tr}\tilde{S}_2\mathrm{d}r\mathrm{d}t\\
				=&
				\int_{1}^{+\infty}\dfrac{\tau r^2}{2\mu}\rho \partial_{tr}v \partial_{tr}\tilde{S}_2\big|_0^t\mathrm{d}r
				-\int_{0}^{t}\int_{1}^{+\infty}\dfrac{\tau r^2}{2\mu}\rho \partial_{tr}v \partial_{ttr}\tilde{S}_2\mathrm{d}r\mathrm{d} t
				-\int_{0}^{t}\int_{1}^{+\infty}\dfrac{\tau r^2}{2\mu}\partial_{t}\rho \partial_{tr}v \partial_{t}\tilde{S}_2\mathrm{d}r\mathrm{d} t\\
				\geq &
				-C\int_{1}^{+\infty}r^2\rho ((\partial_{tr}v)^2+ (\partial_{tr}\tilde{S}_2)^2)\mathrm{d}r
				-C(E(0)+E^{\frac{1}{2}}(t) \int_{0}^{t}\mathcal{D}(s)\mathrm{d}s)\\
				&-\int_{0}^{t}\int_{1}^{+\infty}\dfrac{r^2}{2\mu}\partial_{tr}v
				\left(\lambda(\partial_{trr} v+\dfrac{2}{r}\partial_{tr} v
				-\dfrac{2}{r^2} \partial_{t}v)-\partial_{tr} \tilde{S}_2
				-\tau \rho (v-\epsilon)\rho\partial_{trr}\tilde{S}_2\right)\mathrm{d}r\mathrm{d} t
				\\
				\geq &
				\int_{0}^{t}\frac{\lambda}{4\mu} (\partial_{tr}v)^2(t,1)\mathrm{d} t
				-\int_{0}^{t}\int_{1}^{+\infty}\dfrac{\lambda r}{2\mu} (\partial_{tr}v)^2\mathrm{d}r\mathrm{d} t
				+\int_{0}^{t}\int_{1}^{+\infty}\dfrac{\lambda}{\mu}  \partial_{t}v \partial_{tr}v \mathrm{d}r\mathrm{d} t
				\\
				&+\int_{0}^{t}\int_{1}^{+\infty}\dfrac{r^2}{\mu} \partial_{tr}v\partial_{tr} \tilde{S}_2 \mathrm{d}r\mathrm{d} t
				+\int_{0}^{t}\int_{1}^{+\infty}\dfrac{r^2\tau}{2\mu} \rho(v-\epsilon) \partial_{tr}v\partial_{trr} \tilde{S}_2 \mathrm{d}r\mathrm{d} t\\
				&-C\left(E(0)+ E^{\frac{3}{2}}(t)+ E^{\frac{1}{2}}(t) \int_{0}^{t}\mathcal{D}(s)\mathrm{d}s\right)\\
				\geq &
				\int_{0}^{t}\int_{1}^{+\infty}\dfrac{r^2\tau}{2\mu} \rho(v-\epsilon) \partial_{tr}v\partial_{trr} \tilde{S}_2 \mathrm{d}r\mathrm{d}t
				-C\left(E(0)+ E^{\frac{3}{2}}(t)+ E^{\frac{1}{2}}(t) \int_{0}^{t}\mathcal{D}(s)\mathrm{d}s\right)
			\end{aligned}
		\end{equation}
		
		Finally, we have
		\begin{equation}
			\begin{aligned}
				&\int_{0}^{t}\int_{1}^{+\infty}\dfrac{r^2\tau}{3\mu}\rho\partial_{tr} v\partial_{ttr} \tilde{S}_1\mathrm{d}r\mathrm{d}t\\
				\geq&
				\int_{0}^{t}\frac{\lambda}{4\mu} (\partial_{tr}v)^2(t,1)\mathrm{d} t
				+\int_{0}^{t}\int_{1}^{+\infty}\dfrac{r^2\tau}{2\mu} \rho(v-\epsilon) \partial_{tr}v\partial_{trr} \tilde{S}_2 \mathrm{d}r\mathrm{d}t\\
				&-\int_{0}^{t}\int_{1}^{+\infty}\dfrac{\tau r^2}{2\mu}\rho P^{\prime}(\rho)
				v\partial_{tr}v\partial_{trr}\rho\mathrm{d}r\mathrm{d}t
				-C\left(E(0)+ E^{\frac{3}{2}}(t)+ E^{\frac{1}{2}}(t) \int_{0}^{t}\mathcal{D}(s)\mathrm{d}s\right).	
			\end{aligned}
		\end{equation}	
		
		For the second term on the left-hand side of equation (\ref{76}), by applying equation $(\ref{47})_2$, we get
		\begin{align}\label{103}
				&\int_{0}^{t}\int_{1}^{+\infty}\dfrac{r^2\tau}{3\mu}\rho (v-\epsilon)\partial_{tr} v\partial_{trr} \tilde{S}_1\mathrm{d}r\mathrm{d}t \nonumber \\
				=&\int_{0}^{t}\int_{1}^{+\infty}\dfrac{r^2\tau}{2\mu}\rho (v-\epsilon)\partial_{tr} v(\rho\partial_{ttr} v+\rho v\partial_{trr} v+\partial_{trr} P-\dfrac{2}{r} \partial_{tr} \tilde{S}_1+\dfrac{2}{r^2} \partial_{t} \tilde{S}_1
				-\partial_{trr} \tilde{S}_2-g_2)
				\mathrm{~d}r\mathrm{d}t \nonumber\\
				\geq& 
				\int_{0}^{t}\int_{1}^{+\infty}\dfrac{r^2\tau}{2\mu}\rho P^{\prime}(\rho)(v-\epsilon) \partial_{tr}v\partial_{trr}\rho \mathrm{d}r\mathrm{d}t
				-\int_{0}^{t}\int_{1}^{+\infty}\dfrac{r^2\tau}{2\mu}\rho (v-\epsilon) \partial_{tr}v\partial_{trr} \tilde{S}_2 \mathrm{d}r\mathrm{d}t \nonumber \\ 
				&
				-C\left(E(0)+ E^{\frac{3}{2}}(t)+ E^{\frac{1}{2}}(t) \int_{0}^{t}\mathcal{D}(s)\mathrm{d}s\right).
		\end{align}
		
		For the third term on the left-hand side of equation (\ref{76}), we get
		\begin{equation}
			\begin{aligned}
				\int_{0}^{t}\int_{1}^{+\infty}\dfrac{r^2}{3\mu}\partial_{tr} v\partial_{tr} \tilde{S}_1\mathrm{d}r\mathrm{d}t
				&\geq -C\int_{0}^{t}\int_{1}^{+\infty} r^2((\partial_{tr} v)^2
				+(\partial_{tr} \tilde{S}_1)^2) \mathrm{d}r\mathrm{d}t\\
				&\geq -C\left(E(0)+ E^{\frac{3}{2}}(t)+ E^{\frac{1}{2}}(t) \int_{0}^{t}\mathcal{D}(s)\mathrm{d}s\right).
			\end{aligned}
		\end{equation}
		
		For the last term on the left-hand side of equation (\ref{76}), we get
		\begin{equation}
			\begin{aligned}
				&-\int_{0}^{t}\int_{1}^{+\infty}\frac{2}{3}\left(r^2\partial_{trr}v \partial_{tr}v-r(\partial_{tr}v)^2 +\partial_{t}v\partial_{tr}v\right)\mathrm{d}r\mathrm{d}t\\
				&\geq \int_{0}^{t}\frac{1}{3}(\partial_{tr} v)^2(t,1)\mathrm{d}t- C\left(E(0)+ E^{\frac{3}{2}}(t)+ E^{\frac{1}{2}}(t) \int_{0}^{t}\mathcal{D}(s)\mathrm{d}s\right).
			\end{aligned}
		\end{equation}
		
		Combining the above results, we have
		\begin{equation}\nonumber
			\begin{aligned}
				\int_{0}^{t}\frac{1}{3}(\partial_{tr} v)^2(t,1)\mathrm{d}t
				\leq
				\int_{0}^{t}\int_{1}^{+\infty}\dfrac{r^2\tau\epsilon}{2\mu}\rho P^{\prime}(\rho)\partial_{tr}v\partial_{trr}\rho \mathrm{d}r\mathrm{d}t
				+C\left(E(0)+ E^{\frac{3}{2}}(t)+ E^{\frac{1}{2}}(t) \int_{0}^{t}\mathcal{D}(s)\mathrm{d}s\right).
			\end{aligned}
		\end{equation}
		
		%\begin{equation}
		%	\begin{aligned}
			%		&\int_{0}^{t}\int_{1}^{+\infty}\dfrac{r^2\tau\epsilon}{2\mu}\rho P^{\prime}(\rho)\partial_{tr}v\partial_{trr}\rho \mathrm{d}r\mathrm{d}t\\
			%		=&\int_{0}^{t}\int_{1}^{+\infty}-\dfrac{r^2\tau\epsilon}{2\mu} P^{\prime}(\rho)\partial_{trr}\rho 
			%		(\partial_{tt} \rho+v\partial_{tr}\rho+\partial_{t}\rho\partial_{r}v+\partial_{r}\rho\partial_{t}v+\frac{2}{r}\partial_{t}(\rho v))\mathrm{d}r\mathrm{d}t\\
			%		\leq&
			%		\\
			%		& \leq C\left(E(0)+E^{\frac{3}{2}}(t) +E^{\frac{1}{2}}(t) \int_{0}^{t}\mathcal{D}(s)\mathrm{d}s\right)
			%		,
			%	\end{aligned}
		%\end{equation}
		In view of equation $(\ref{47})_1$ and Lemmas \ref{lem3.4}, \ref{lem4} and \ref{lem5}, we derive
		\begin{equation}\nonumber
			\begin{aligned}
				&\int_{0}^{t}\int_{1}^{+\infty}\dfrac{r^2\tau\epsilon}{2\mu}\rho P^{\prime}(\rho)\partial_{tr}v\partial_{trr}\rho \mathrm{d}r\mathrm{d}t\\
				=&-\int_{0}^{t}\dfrac{r^2\tau\epsilon}{2\mu}\rho P^{\prime}(\rho)\partial_{tr}v\partial_{tr}\rho(t,1)\mathrm{d}t-
				\int_{0}^{t}\int_{1}^{+\infty}\dfrac{\tau\epsilon}{2\mu}[r^2\rho P^{\prime}(\rho)\partial_{tr}v]_r\partial_{tr}\rho \mathrm{d}r\mathrm{d}t\\
				\leq&-\int_{0}^{t}\dfrac{\tau\epsilon}{2\mu}\rho P^{\prime}(\rho)\partial_{tr}v\partial_{tr}\rho(t,1)\mathrm{d}t
				+\int_{0}^{t}\int_{1}^{+\infty}\dfrac{\tau\epsilon}{2\mu}r^2P^{\prime}(\rho)\partial_{tr}\rho(\partial_{ttr} \rho+v\partial_{trr}\rho+\frac{2}{r}\rho\partial_{tr}v-\frac{2}{r^2}\rho \partial_{t}v)\mathrm{d}r\mathrm{d}t\\
				&-
				\int_{0}^{t}\int_{1}^{+\infty}\dfrac{\tau\epsilon}{2\mu}2r\rho P^{\prime}(\rho)\partial_{tr}v\partial_{tr}\rho \mathrm{d}r\mathrm{d}t
				+ C\left(E(0)+E^{\frac{1}{2}}(t) \int_{0}^{t}\mathcal{D}(s)\mathrm{d}s\right)
				\\
				\leq&\dfrac{C\tau\epsilon}{2\mu}\int_{0}^{t}\left((\partial_{tr}v)^2+(\partial_{tr}\rho)^2\right)(t,1) \mathrm{d}t
				+\dfrac{C\tau\epsilon}{\mu}\int_{0}^{t}\int_{1}^{+\infty}\left((\partial_{t}v)^2+(\partial_{tr}v)^2+(\partial_{tr}\rho)^2\right)
				\mathrm{d}r\mathrm{d}t
				\\
				&+ C\left(E(0)+E^{\frac{3}{2}}(t) +E^{\frac{1}{2}}(t) \int_{0}^{t}\mathcal{D}(s)\mathrm{d}s\right)\\
				\leq&\dfrac{C\tau\epsilon}{2\mu}\int_{0}^{t}(\partial_{tr}v)^2(t,1) \mathrm{d}t
				+ C\left(E(0)+E^{\frac{3}{2}}(t) +E^{\frac{1}{2}}(t) \int_{0}^{t}\mathcal{D}(s)\mathrm{d}s\right),
			\end{aligned}
		\end{equation}
		where we use equation $(\ref{15})_2$, leading to
		\begin{equation}\nonumber
			\begin{aligned}	
				\dfrac{C\tau\epsilon}{2\mu}\int_{0}^{t}(\partial_{tr}\rho)^2(t,1) \mathrm{d}t
				&\leq \dfrac{C\tau\epsilon}{2\mu}\int_{0}^{t}((\partial_{tr}\tilde{S}_1)^2+(\partial_{tr}\tilde{S}_2)^2+(\partial_{t}\tilde{S}_1)^2)(t,1)\mathrm{d}t\\
				&\leq C\left(E(0)+E^{\frac{3}{2}}(t) +E^{\frac{1}{2}}(t) \int_{0}^{t}\mathcal{D}(s)\mathrm{d}s\right).
			\end{aligned}
		\end{equation}
		Thus, we can obtain
		\begin{equation}\label{89}
			\begin{aligned}
				\int_{0}^{t}(\partial_{tr} v)^2(t,1)\mathrm{d}t
				\leq C\left(E(0)+E^{\frac{3}{2}}(t)+E^{\frac{1}{2}}(t) \int_{0}^{t}\mathcal{D}(s)\mathrm{d}s\right).
			\end{aligned}
		\end{equation}
	This finish the proof of Lemma \ref{lem7_tx}.
	\end{proof}                                                                                                                                                                                                       
	\begin{lemma}\label{lem7_xx}
		There exists a constant $C$ such that for any $0\leq t \leq T$
		\begin{align}\label{bdy}
				&\int_{0}^{t}(\partial_{rr} v)^2(t,1)\mathrm{~d}t \nonumber\\
				\leq& \frac{\tau C}{2\mu}(\eta+\epsilon)\int_{1}^{+\infty}r^2\rho(\partial_{rr} v)^2\mathrm{d}r
				+\eta\int_{0}^{t}\int_{1}^{+\infty}\frac{r^2}{2\mu}\left((\partial_{rr}\tilde{S}_1)^2+(\partial_{rr}\tilde{S}_2)^2\right)\mathrm{d}r\mathrm{d}t \nonumber\\
				&
				+\frac{C\tau\epsilon}{\mu}\int_{0}^{t}\eta\left(
				(\partial_{rr} \tilde{S}_1)^2+(\partial_{rr} \tilde{S}_2)^2\right)(t,1)\mathrm{d}t
				+C(\eta)\epsilon\int_{0}^{t}\int_{1}^{+\infty}r^2((\partial_{rr}\tilde{S}_1)^2+(\partial_{rr}\tilde{S}_2)^2)\mathrm{d}r\mathrm{d}t \nonumber\\
				&
				+C\left(E(0)+ E^{\frac{3}{2}}(t)+ E^{\frac{1}{2}}(t) \int_{0}^{t}\mathcal{D}(s)\mathrm{d}s\right)
			\end{align}
		where $\eta$ is any small constant to be chosen later.
	\end{lemma}
	\begin{proof}
		Multiplying the equation $(\ref{60})_3$ by $\frac{r^2}{3\mu}\partial_{rr} v$
		and integrating over $(0,t)\times[1,+\infty)$, we get
		\begin{align}\label{70}
				&\int_{0}^{t}\int_{1}^{+\infty}\frac{r^2\tau}{3\mu}\rho\partial_{rr} v\partial_{trr} \tilde{S}_1\mathrm{d}r\mathrm{d}t +\int_{0}^{t}\int_{1}^{+\infty}\dfrac{r^2\tau}{3\mu}\rho (v-\epsilon) \partial_{rr} v\partial_{rrr} \tilde{S}_1\mathrm{d}r\mathrm{d}t\nonumber \\
				& +\int_{0}^{t}\int_{1}^{+\infty}\frac{r^2}{3\mu}\partial_{rr} v\partial_{rr} \tilde{S}_1\mathrm{d}r\mathrm{d}t 
				-\int_{0}^{t}\int_{1}^{+\infty}\frac{2}{3}\left(r^2\partial_{rrr}v \partial_{rr}v-r(\partial_{rr} v)^2+2\partial_{r}v\partial_{rr}v-\frac{2}{r}v\partial_{rr} v\right)\mathrm{d}r\mathrm{d}t \nonumber\\
				&\leq C\left(E(0)+ E^{\frac{1}{2}}(t) \int_{0}^{t}\mathcal{D}(s)\mathrm{d}s\right).
		\end{align}	
		We estimate each term of the above equation, respectively. Firstly, by using equation $(\ref{22})_2$, we have
		\begin{align}\label{71}
				&\int_{0}^{t}\int_{1}^{+\infty}\dfrac{r^2\tau}{3\mu}\rho\partial_{rr} v\partial_{trr} \tilde{S}_1\mathrm{d}r\mathrm{d}t \nonumber \\
				=&
				\int_{1}^{+\infty}\dfrac{r^2\tau}{3\mu}\rho\partial_{rr} v\partial_{rr} \tilde{S}_1\big|_0^t\mathrm{d}r-
				\int_{0}^{t}\int_{1}^{+\infty}\dfrac{r^2\tau}{3\mu}\rho\partial_{trr} v\partial_{rr} \tilde{S}_1\mathrm{d}r\mathrm{d}t-
				\int_{0}^{t}\int_{1}^{+\infty}\dfrac{r^2\tau}{3\mu}\partial_{t}\rho\partial_{rr} v\partial_{rr} \tilde{S}_1\mathrm{d}r\mathrm{d}t \nonumber \\
				\geq& \int_{1}^{+\infty}\dfrac{r^2\tau}{3\mu}\rho\partial_{rr} v\partial_{rr} \tilde{S}_1\mathrm{d}r-
				\int_{0}^{t}\int_{1}^{+\infty}\dfrac{r^2\tau}{2\mu}\rho\partial_{trr} v(\rho\partial_{tr} v+\rho v\partial_{rr} v+\partial_{rr} P-\dfrac{2}{r} \partial_r \tilde{S}_1
				+\dfrac{2}{r^2}\tilde{S}_1 \nonumber \\
				&\qquad\qquad\qquad\qquad\qquad\qquad\qquad\qquad-\partial_{rr} \tilde{S}_2)\mathrm{d}r\mathrm{d}t-C\left(E(0)+ E^{\frac{1}{2}}(t) \int_{0}^{t}\mathcal{D}(s)\mathrm{d}s\right),
			\end{align}
		where 
		\begin{equation}\nonumber
			\begin{aligned}
				&-\int_{0}^{t}\int_{1}^{+\infty}\dfrac{r^2\tau}{2\mu}\rho^2\partial_{trr} v\partial_{tr} v\mathrm{d}r\mathrm{d}t
				-\int_{0}^{t}\int_{1}^{+\infty}\dfrac{r^2\tau}{2\mu}\rho^2v\partial_{trr} v\partial_{rr} v\mathrm{d}r\mathrm{d}t\\
				\geq&\int_{0}^{t}\frac{\tau}{4\mu}\rho^2(\partial_{tr}v)^2(t,1)\mathrm{d}t
				+\int_{0}^{t}\int_{1}^{+\infty}\dfrac{\tau}{4\mu}(2r\rho^2+2r^2\rho\partial_{r}\rho)(\partial_{tr}v)^2\mathrm{d}r\mathrm{d}t
				-C\left(E(0)+ E^{\frac{1}{2}}(t) \int_{0}^{t}\mathcal{D}(s)\mathrm{d}s\right)\\
				\geq& -C\left(E(0)+ E^{\frac{3}{2}}(t)+ E^{\frac{1}{2}}(t) \int_{0}^{t}\mathcal{D}(s)\mathrm{d}s\right),
			\end{aligned}
		\end{equation}
		and 
		\begin{equation}\nonumber
			\begin{aligned}
				&-\int_{0}^{t}\int_{1}^{+\infty}\dfrac{\tau r^2}{2\mu}\rho\partial_{trr} v\partial_{rr} P\mathrm{d}r\mathrm{d}t\\
				\geq&\int_{1}^{+\infty}\frac{\tau r^2}{2\mu}\rho\partial_{rr} v\partial_{rr} P\mathrm{d}r +\int_{0}^{t}\int_{1}^{+\infty}\dfrac{\tau r^2}{2\mu}\rho\partial_{rr} v P^{\prime}(\rho)\partial_{trr} \rho \mathrm{d}r\mathrm{d}t
				-C\left(E(0)+ E^{\frac{1}{2}}(t) \int_{0}^{t}\mathcal{D}(s)\mathrm{d}s\right).
			\end{aligned}
		\end{equation}	
		By using equation $(\ref{60})_1$, Lemmas \ref{lem3.4} and \ref{lem5}, we have
	
			\begin{align*}
				&\int_{0}^{t}\int_{1}^{+\infty}\dfrac{\tau r^2}{2\mu}\rho\partial_{rr} v P^{\prime}(\rho)\partial_{trr} \rho \mathrm{d}r\mathrm{d}t\\
				=&-\int_{0}^{t}\int_{1}^{+\infty}\frac{\tau r^2}{2\mu}\rho\partial_{rr} v P^{\prime}(\rho)(\partial_{rrr}(\rho v)+\frac{2}{r} \partial_{rr}(\rho v)-\frac{4}{r^2} \partial_{r}(\rho v)+\frac{4}{r^3}\rho v)\mathrm{d}r\mathrm{d}t\\
				\geq&
				\int_{0}^{t}\dfrac{\tau r^2}{4\mu}P^{\prime}(\rho)\rho^2(\partial_{rr} v)^2(t,1)\mathrm{d}t
				-\int_{0}^{t}\int_{1}^{+\infty}\dfrac{\tau r}{2\mu}
				\rho^2 P^{\prime}(\rho)(\partial_{rr} v)^2\mathrm{d}r\mathrm{d}t\\
				&-\int_{0}^{t}\int_{1}^{+\infty}\dfrac{\tau r^2}{2\mu}
				\rho P^{\prime}(\rho)v\partial_{rr} v\partial_{rrr}\rho\mathrm{d}r\mathrm{d}t
				+\int_{0}^{t}\int_{1}^{+\infty}\dfrac{2\tau}{\mu}
				\rho^2 P^{\prime}(\rho)\partial_{rr} v\partial_{r} v\mathrm{d}t\\
				&
				-\int_{0}^{t}\int_{1}^{+\infty}\dfrac{2\tau}{\mu r}
				\rho^2 P^{\prime}(\rho)v\partial_{rr}v\mathrm{d}r\mathrm{d}t
				-C\left(E(0)+E^{\frac{1}{2}}(t) \int_{0}^{t}\mathcal{D}(s)\mathrm{d}s\right)
				\displaybreak \\
				\geq&
				-\int_{0}^{t}\int_{1}^{+\infty}\dfrac{\tau r^2}{2\mu}
				\rho P^{\prime}(\rho)v\partial_{rr} v\partial_{rrr}\rho\mathrm{d}r\mathrm{d}t
				-C\int_{0}^{t}\int_{1}^{+\infty}r^2((\partial_{rr} v)^2+(\partial_{r} v)^2+v^2)\mathrm{d}r\mathrm{d}t\\
				&\qquad\qquad\qquad\qquad\qquad\qquad\qquad\qquad\qquad
				-C\left(E(0)+ E^{\frac{3}{2}}(t)+ E^{\frac{1}{2}}(t) \int_{0}^{t}\mathcal{D}(s)\mathrm{d}s\right)\\
				\geq&-\int_{0}^{t}\int_{1}^{+\infty}\dfrac{\tau r^2}{2\mu}
				\rho P^{\prime}(\rho)v\partial_{rr} v\partial_{rrr}\rho\mathrm{d}r\mathrm{d}t
				-C\epsilon\int_{0}^{t}\int_{1}^{+\infty}r^2((\partial_{rr}\tilde{S}_1)^2+(\partial_{rr}\tilde{S}_2)^2)\mathrm{d}r\mathrm{d}t\\
				&\qquad\qquad\qquad\qquad\qquad\qquad\qquad\qquad\qquad
				-C\left(E(0)+ E^{\frac{3}{2}}(t)+ E^{\frac{1}{2}}(t) \int_{0}^{t}\mathcal{D}(s)\mathrm{d}s\right)
			\end{align*}
		since $P^{\prime}(\rho)>0$. 
		Next, we continue to handle the remaining terms in (\ref{71}). By applying Lemmas \ref{lem4} and  $\ref{lem5}$, we have	 
		\begin{equation}\nonumber
			\begin{aligned}
				&\int_{0}^{t}\int_{1}^{+\infty}\frac{r\tau}{\mu}\rho\partial_{trr}v
				\partial_{r} \tilde{S}_1\mathrm{d}r\mathrm{d}t\\
				=&\int_{1}^{+\infty} \frac{r\tau}{\mu}\rho\partial_{rr}v
				\partial_{r} \tilde{S}_1\big|_0^t\mathrm{d}r
				-\int_{0}^{t}\int_{1}^{+\infty} \frac{r\tau}{\mu}
				\rho\partial_{tr} \tilde{S}_1\partial_{rr}v\mathrm{d}r\mathrm{d}t
				-\int_{0}^{t}\int_{1}^{+\infty} \frac{r\tau}{\mu}
				\partial_{t}\rho\partial_{r} \tilde{S}_1 \partial_{rr}v\mathrm{d}r\mathrm{d}t\\
				\geq&\int_{1}^{+\infty} \frac{r\tau}{\mu}\rho\partial_{rr}v
				\partial_{r} \tilde{S}_1\mathrm{d}r
				-C\int_{0}^{t}\int_{1}^{+\infty} r^2((\partial_{rr}v)^2+(\partial_{tr} \tilde{S}_1)^2)\mathrm{d}r\mathrm{d}t
				-C\left(E(0)+ E^{\frac{1}{2}}(t) \int_{0}^{t}\mathcal{D}(s)\mathrm{d}s\right)\\
				\geq&\int_{1}^{+\infty} \frac{r\tau}{\mu}\rho\partial_{rr}v
				\partial_{r} \tilde{S}_1\mathrm{d}r
				-C\epsilon\int_{0}^{t}\int_{1}^{+\infty}r^2((\partial_{rr} \tilde{S}_1)^2
				+(\partial_{rr} \tilde{S}_2)^2)\mathrm{d}r\mathrm{d}t\\
				&\qquad\qquad\qquad\qquad\qquad\qquad\qquad\qquad\qquad\qquad\qquad\quad-C\left(E(0)+ E^{\frac{3}{2}}(t)+ E^{\frac{1}{2}}(t) \int_{0}^{t}\mathcal{D}(s)\mathrm{d}s\right).
			\end{aligned}
		\end{equation}
		By applying Lemmas \ref{lem2} and  $\ref{lem5}$, we have	
			\begin{align*}
				&-\int_{0}^{t}\int_{1}^{+\infty}\frac{\tau}{\mu}\rho\partial_{trr}v
				\tilde{S}_1\mathrm{d}r\mathrm{d}t\\
				=&-\int_{1}^{+\infty} \frac{\tau}{\mu}\rho\partial_{rr}v \tilde{S}_1\big|_0^t\mathrm{d}r
				+\int_{0}^{t}\int_{1}^{+\infty} \frac{\tau}{\mu}
				\rho \partial_{t}\tilde{S}_1\partial_{rr}v\mathrm{d}r\mathrm{d}t +\int_{0}^{t}\int_{1}^{+\infty} \frac{\tau}{\mu}
				\partial_{t}\rho \partial_{rr}v\tilde{S}_1\mathrm{d}r\mathrm{d}t\\
				\geq&-\int_{1}^{+\infty} \frac{\tau}{\mu}\rho\partial_{rr}v \tilde{S}_1\mathrm{d}r
				-\int_{0}^{t}\int_{1}^{+\infty} \frac{\tau\rho}{2\mu}((\partial_{rr}v)^2
				+(\partial_{t}\tilde{S}_1)^2) \mathrm{d}r
				+\int_{0}^{t}\int_{1}^{+\infty} \frac{\tau}{\mu}
				\partial_{t}\rho \partial_{rr}v\tilde{S}_1\mathrm{d}r\mathrm{d}t\\
				\geq&-\int_{1}^{+\infty} \frac{\tau}{\mu}\rho\partial_{rr}v \tilde{S}_1\mathrm{d}r
				-C\epsilon\int_{0}^{t}\int_{1}^{+\infty}r^2((\partial_{rr} \tilde{S}_1)^2
				+(\partial_{rr} \tilde{S}_2)^2)\mathrm{d}r\mathrm{d}t\\
				&\qquad\qquad\qquad\qquad\qquad\qquad\qquad\qquad\qquad\qquad\qquad-C\left(E(0)+ E^{\frac{3}{2}}(t)+ E^{\frac{1}{2}}(t) \int_{0}^{t}\mathcal{D}(s)\mathrm{d}s\right).
			\end{align*}
	
		In view of equation $(\ref{60})_4$ and Lemmas \ref{lem3.4} and \ref{lem5}, we have
			\begin{align*}
				&\int_{0}^{t}\int_{1}^{+\infty}\frac{r^2\tau}{2\mu}\rho\partial_{trr}v
				\partial_{rr} \tilde{S}_2\mathrm{d}r\mathrm{d}t\\
				=&\int_{1}^{+\infty} \frac{r^2\tau}{2\mu}\rho\partial_{rr}v
				\partial_{rr} \tilde{S}_2\big|_0^t\mathrm{d}r
				-\int_{0}^{t}\int_{1}^{+\infty} \frac{r^2\tau}{2\mu}\rho\partial_{rr}v
				\partial_{trr} \tilde{S}_2\mathrm{d}r\mathrm{d}t
				-\int_{0}^{t}\int_{1}^{+\infty} \frac{r^2\tau}{2\mu}\partial_{t}\rho\partial_{rr}v
				\partial_{rr} \tilde{S}_2\mathrm{d}r\mathrm{d}t\\
				\geq&	\int_{1}^{+\infty} \frac{r^2\tau}{2\mu}\rho\partial_{rr}v
				\partial_{rr} \tilde{S}_2\mathrm{d}r
				-C\left(E(0)+ E^{\frac{1}{2}}(t) \int_{0}^{t}\mathcal{D}(s)\mathrm{d}s\right)\\
				&-\int_{0}^{t}\int_{1}^{+\infty}\frac{r^2}{2\mu}\partial_{rr}v
				(\lambda(\partial_{rrr} v+\dfrac{2}{r}\partial_{rr} v-\dfrac{4}{r^2} \partial_{r}v+\dfrac{4}{r^3}v)-\tau  \rho (v-\epsilon) \partial_{rrr}\tilde{S}_2 
				-\partial_{rr} \tilde{S}_2)\mathrm{d}r\mathrm{d}t \displaybreak \\
				\geq&\int_{1}^{+\infty}\frac{r^2\tau}{2\mu}\rho\partial_{rr}v\partial_{rr} \tilde{S}_2\mathrm{d}r
				+\int_{0}^{t}\int_{1}^{+\infty}\dfrac{r^2\tau}{2\mu}\rho  (v-\epsilon)  \partial_{rr}v\partial_{rrr} \tilde{S}_2 \mathrm{d}r\mathrm{d}t
				-\eta\int_{0}^{t}\int_{1}^{+\infty}\frac{r^2}{2\mu}(\partial_{rr} \tilde{S}_2)^2\mathrm{d}r\mathrm{d}t\\
				&
				-C(\eta)\epsilon\int_{0}^{t}\int_{1}^{+\infty}r^2((\partial_{rr} \tilde{S}_1)^2
				+(\partial_{rr} \tilde{S}_2)^2)\mathrm{d}r\mathrm{d}t
				-C\left(E(0)+ E^{\frac{3}{2}}(t)+ E^{\frac{1}{2}}(t) \int_{0}^{t}\mathcal{D}(s)\mathrm{d}s\right).
			\end{align*}
		
		For the second term on the left-hand side of (\ref{70}), by applying $(\ref{60})_2$, we get
		\begin{align}\label{88}
				&\int_{0}^{t}\int_{1}^{+\infty}\dfrac{r^2\tau}{3\mu}\rho (v-\epsilon)\partial_{rr} v\partial_{rrr} \tilde{S}_1\mathrm{d}r\mathrm{d}t \nonumber\\
				=&\int_{0}^{t}\int_{1}^{+\infty}\dfrac{r^2\tau}{2\mu}\rho (v-\epsilon)\partial_{rr} v\big(\rho\partial_{trr} v+\rho v\partial_{rrr} v+\partial_{rrr} P-\dfrac{2}{r} \partial_{rr} \tilde{S}_1+\dfrac{4}{r^2} \partial_{r} \tilde{S}_1
				+\dfrac{4}{r^3}\tilde{S}_1-\partial_{rrr} \tilde{S}_2 \nonumber\\
				&\qquad\qquad\qquad\qquad\qquad\qquad\qquad\qquad\qquad\qquad\qquad\qquad\qquad\qquad\qquad\qquad\qquad\qquad-l_2\big)
				\mathrm{d}r\mathrm{d}t \nonumber\\
				\geq&\int_{1}^{+\infty}-\dfrac{\tau\epsilon}{4\mu}r^2\rho^2(\partial_{rr} v)^2\mathrm{d}r
				-C\epsilon\int_{0}^{t}\int_{1}^{+\infty}r^2\left((\partial_{rr} \tilde{S}_1)^2+(\partial_{rr} \tilde{S}_2)^2\right)\mathrm{d}r\mathrm{d}t
				\nonumber\\
				&+\int_{0}^{t}\int_{1}^{+\infty}\dfrac{\tau r^2}{2\mu}\rho v \partial_{rr}v P^{\prime}(\rho) \partial_{rrr}\rho\mathrm{d}r\mathrm{d}t
				-\int_{0}^{t}\int_{1}^{+\infty}\dfrac{r^2\tau}{2\mu}\rho (v-\epsilon) \partial_{rr}v\partial_{rrr} \tilde{S}_2 \mathrm{d}r\mathrm{d}t \nonumber\\
				&-\int_{0}^{t}\int_{1}^{+\infty}\dfrac{r^2\tau\epsilon}{2\mu}\rho P^{\prime}(\rho)\partial_{rr}v\partial_{rrr}\rho \mathrm{d}r\mathrm{d}t
				-C\left(E(0)+ E^{\frac{1}{2}}(t) \int_{0}^{t}\mathcal{D}(s)\mathrm{d}s\right)
			\end{align}
		where we use Lemmas \ref{lem1}, \ref{lem3} and \ref{lem5}, leading to
		\begin{equation}\nonumber
			\begin{aligned}
				&-\int_{0}^{t}\int_{1}^{+\infty}\frac{\tau r}{\mu}\rho(v-\epsilon) \partial_{rr} \tilde{S}_1\partial_{rr}v\mathrm{d}r\mathrm{d}t\\
				\geq&-\frac{\tau\epsilon}{\mu} \int_{0}^{t}\int_{1}^{+\infty}Cr^2 \left((\partial_{rr}v)^2+(\partial_{rr} \tilde{S}_1)^2\right)\mathrm{d}r\mathrm{d}t
				-C\left(E(0)+ E^{\frac{1}{2}}(t) \int_{0}^{t}\mathcal{D}(s)\mathrm{d}s\right)\\
				\geq&-C\epsilon\int_{0}^{t}\int_{1}^{+\infty}r^2\left((\partial_{rr} \tilde{S}_1)^2+(\partial_{rr} \tilde{S}_2)^2\right)\mathrm{d}r\mathrm{d}t
				-C\left(E(0)+ E^{\frac{3}{2}}(t)+ E^{\frac{1}{2}}(t) \int_{0}^{t}\mathcal{D}(s)\mathrm{d}s\right),
			\end{aligned}
		\end{equation}
		and
		\begin{equation}\nonumber
			\begin{aligned}
				&-\int_{0}^{t}\int_{1}^{+\infty}\frac{2\epsilon\tau}{\mu}\rho\partial_{r} \tilde{S}_1\partial_{rr}v\mathrm{d}r\mathrm{d}t
				-\int_{0}^{t}\int_{1}^{+\infty}\dfrac{2\tau\epsilon}{\mu r}\rho \partial_{rr}v\tilde{S}_1\mathrm{d}r\mathrm{d}t\\
				\geq&-C\epsilon\int_{0}^{t}\int_{1}^{+\infty}r^2\left((\partial_{rr} \tilde{S}_1)^2+(\partial_{rr} \tilde{S}_2)^2\right)\mathrm{d}r\mathrm{d}t
				-C\left(E(0)+ E^{\frac{3}{2}}(t)+ E^{\frac{1}{2}}(t) \int_{0}^{t}\mathcal{D}(s)\mathrm{d}s\right).
			\end{aligned}
		\end{equation}
		
		For the third term on the left-hand side of equation (\ref{70}), we have
		\begin{equation}
			\begin{aligned}
				&\int_{0}^{t}\int_{1}^{+\infty}\dfrac{r^2}{3\mu}\partial_{rr} v\partial_{rr} \tilde{S}_1\mathrm{d}r\mathrm{d}t
				\geq \int_{0}^{t}\int_{1}^{+\infty}\left(-C(\eta) \dfrac{r^2}{3\mu}(\partial_{rr} v)^2
				- \eta \dfrac{r^2}{3\mu}(\partial_{rr} \tilde{S}_1)^2\right) \mathrm{d}r\mathrm{d}t\\
				\geq& -C(\eta)\epsilon\int_{0}^{t}\int_{1}^{+\infty}r^2\left((\partial_{rr} \tilde{S}_1)^2+(\partial_{rr} \tilde{S}_2)^2\right)\mathrm{d}r\mathrm{d}t
				-\eta \int_{0}^{t}\int_{1}^{+\infty}\dfrac{r^2}{3\mu}(\partial_{rr} \tilde{S}_1)^2\mathrm{d}r\mathrm{d}t
			\end{aligned}
		\end{equation}
		where $\eta$ is sufficiently small.
		For the last term on the left-hand side of equation (\ref{70}), we have
			\begin{align*}
				&-\int_{0}^{t}\int_{1}^{+\infty}\frac{2}{3}\left(r^2\partial_{rrr}v \partial_{rr}v-r(\partial_{rr} v)^2+2\partial_{r}v\partial_{rr}v-\frac{2}{r}v\partial_{rr} v\right)\mathrm{d}r\mathrm{d}t\\
				\geq&\int_{0}^{t}\frac{1}{3}(\partial_{rr} v)^2(t,1) \mathrm{d}t
				-C\bigg(\epsilon\int_{0}^{t}\int_{1}^{+\infty}r^2\left((\partial_{rr} \tilde{S}_1)^2+(\partial_{rr} \tilde{S}_2)^2\right)\mathrm{d}r\mathrm{d}t\\
				&\qquad\qquad\qquad\qquad\qquad\qquad\quad+
				E(0)+ E^{\frac{3}{2}}(t)+ E^{\frac{1}{2}}(t) \int_{0}^{t}\mathcal{D}(s)\mathrm{d}s\bigg).
			\end{align*}
		
		
		Combining the above result, one can obtain
			\begin{align*}
				&\int_{0}^{t}\frac{1}{3}
				(\partial_{rr}v)^2(t,1)\mathrm{d}t\\
				\leq& \int_{1}^{+\infty}\left(
				\frac{\tau r^2}{2\mu}\rho\partial_{rr} v (\partial_{rr} P
				-\frac{2}{3}\partial_{rr} \tilde{S}_1
				-\frac{2}{r}\partial_{r}\tilde{S}_1+\frac{2}{r^2}\tilde{S}_1-\partial_{rr}\tilde{S}_2 ) 
				+\frac{\epsilon\tau r^2}{4\mu}\rho^2(\partial_{rr} v)^2 \right)\mathrm{d}r\\
				&+\eta\int_{0}^{t}\int_{1}^{+\infty}
				\frac{r^2}{2\mu}\left((\partial_{rr} \tilde{S}_1)^2+(\partial_{rr} \tilde{S}_2)^2\right)
				\mathrm{d}r\mathrm{d}t
				+\int_{0}^{t}\int_{1}^{+\infty}\dfrac{r^2\tau\epsilon}{2\mu}\rho P^{\prime}(\rho)\partial_{rr}v\partial_{rrr}\rho \mathrm{d}r\mathrm{d}t\\
				&+C(\eta)\epsilon\int_{0}^{t}\int_{1}^{+\infty}r^2\left((\partial_{rr} \tilde{S}_1)^2+(\partial_{rr} \tilde{S}_2)^2\right)\mathrm{d}r\mathrm{d}t\\
				\leq&\frac{\tau C}{2\mu}\int_{1}^{+\infty} r^2\rho \left(\eta(\partial_{rr}v)^2+C(\eta)(\partial_{tr}v)^2
				+ \epsilon (\partial_{rr}v)^2 \right)\mathrm{d}r
				+\eta\int_{0}^{t}\int_{1}^{+\infty}
				\frac{r^2}{2\mu}\left((\partial_{rr} \tilde{S}_1)^2+(\partial_{rr} \tilde{S}_2)^2\right)
				\mathrm{d}r\mathrm{d}t\\
				&+\int_{0}^{t}\int_{1}^{+\infty}\dfrac{r^2\tau\epsilon}{2\mu}\rho P^{\prime}(\rho)\partial_{rr}v\partial_{rrr}\rho \mathrm{d}r\mathrm{d}t
				+C(\eta)\epsilon\int_{0}^{t}\int_{1}^{+\infty}r^2\left((\partial_{rr} \tilde{S}_1)^2+(\partial_{rr} \tilde{S}_2)^2\right)\mathrm{d}r\mathrm{d}t\\
				&
				+C\left(E(0)+ E^{\frac{3}{2}}(t)+ E^{\frac{1}{2}}(t) \int_{0}^{t}\mathcal{D}(s)\mathrm{d}s\right).
			\end{align*}
		In view of equation $(\ref{60})_1$ and $(\ref{60})_2$ and Lemmas \ref{lem5}, we obtain
		\begin{equation}\nonumber
			\begin{aligned}
				&\int_{0}^{t}\int_{1}^{+\infty}\dfrac{r^2\tau\epsilon}{2\mu}\rho P^{\prime}(\rho)\partial_{rr}v\partial_{rrr}\rho \mathrm{d}r\mathrm{d}t\\
				=&-\int_{0}^{t}\dfrac{r^2\tau\epsilon}{2\mu}\rho P^{\prime}(\rho)\partial_{rr}v\partial_{rr}\rho(t,1)\mathrm{d}t-
				\int_{0}^{t}\int_{1}^{+\infty}\dfrac{\tau\epsilon}{2\mu}[r^2\rho P^{\prime}(\rho)\partial_{rr}v]_r\partial_{rr}\rho \mathrm{d}r\mathrm{d}t\\
				\leq&-\int_{0}^{t}\dfrac{r^2\tau\epsilon}{2\mu}\rho P^{\prime}(\rho)\partial_{rr}v\partial_{rr}\rho(t,1)\mathrm{d}t
				+\int_{0}^{t}\int_{1}^{+\infty}\dfrac{r^2\tau\epsilon}{2\mu}P^{\prime}(\rho)(\partial_{trr} \rho+v\partial_{rrr}\rho+\frac{2}{r}\rho\partial_{rr}v+\frac{4}{r^3}\rho v)\partial_{rr}\rho\mathrm{d}r\mathrm{d}t\\
				&-
				\int_{0}^{t}\int_{1}^{+\infty}\dfrac{\tau\epsilon}{2\mu}2r\rho P^{\prime}(\rho)\partial_{rr}v\partial_{rr}\rho \mathrm{d}r\mathrm{d}t
				+ CE^{\frac{1}{2}}(t) \int_{0}^{t}\mathcal{D}(s)\mathrm{d}s
				\\
				\leq& \frac{C\tau\epsilon}{2\mu}\int_{0}^{t}\left(C(\eta)(\partial_{rr}v)^2(t,1)+\eta(\partial_{rr}\rho)^2(t,1)\right)\mathrm{d}t
				-\int_{1}^{+\infty}\dfrac{\tau\epsilon}{4\mu}P^{\prime}(\rho)(\partial_{rr}\rho)^2\big|_0^t\mathrm{d}r\\
				&
				+ C\left(E(0)+E^{\frac{3}{2}}(t) +E^{\frac{1}{2}}(t) \int_{0}^{t}\mathcal{D}(s)\mathrm{d}s\right)\\
				\leq& \frac{C\tau\epsilon}{2\mu}\int_{0}^{t}C(\eta)(\partial_{rr}v)^2(t,1)\mathrm{d}t
				+\frac{C\tau\epsilon}{\mu}\int_{0}^{t}\eta\left(
				(\partial_{tr}v)^2+(\partial_{rr} \tilde{S}_1)^2+(\partial_{rr} \tilde{S}_2)^2+(\partial_{r} \tilde{S}_1)^2+(\tilde{S}_1)^2\right)(t,1)\mathrm{d}t\\
				&-\int_{1}^{+\infty}\dfrac{\tau\epsilon}{4\mu}P^{\prime}(\rho)(\partial_{rr}\rho)^2\mathrm{d}r
				+ C\left(E(0)+E^{\frac{3}{2}}(t) +E^{\frac{1}{2}}(t) \int_{0}^{t}\mathcal{D}(s)\mathrm{d}s\right)\\
				\leq&
				\frac{C\tau\epsilon}{2\mu}\int_{0}^{t}C(\eta)(\partial_{rr}v)^2(t,1)\mathrm{d}t
				+\frac{C\tau\epsilon}{\mu}\int_{0}^{t}\eta\left(
				(\partial_{rr} \tilde{S}_1)^2+(\partial_{rr} \tilde{S}_2)^2\right)(t,1)\mathrm{d}t\\
				&
				+ C\left(E(0)+E^{\frac{3}{2}}(t) +E^{\frac{1}{2}}(t) \int_{0}^{t}\mathcal{D}(s)\mathrm{d}s\right)
			\end{aligned}
		\end{equation}
		since $P^{\prime}(\rho)>0$.
	
		Thus, 
		by using Lemmas $\ref{lem1}, \ref{lem2-1}$ and \ref{lem5}, we finally have
		
		\begin{equation}
		\begin{aligned}\label{82}
				&\int_{0}^{t}(\partial_{rr} v)^2(t,1)\mathrm{~d}t \\
				\leq& 
				\frac{\tau C}{2\mu}(\eta+\epsilon)\int_{1}^{+\infty} r^2\rho(\partial_{rr}v)^2\mathrm{d}r +\eta\int_{0}^{t}\int_{1}^{+\infty}\frac{r^2}{2\mu}\left((\partial_{rr}\tilde{S}_1)^2+(\partial_{rr}\tilde{S}_2)^2\right)\mathrm{d}r\mathrm{d}t \\&
				+\frac{C\tau\epsilon}{2\mu}\int_{0}^{t}C(\eta)(\partial_{rr}v)^2(t,1)\mathrm{d}t
				+\frac{C\tau\epsilon}{\mu}\int_{0}^{t}\eta\left(
				(\partial_{rr} \tilde{S}_1)^2+(\partial_{rr} \tilde{S}_2)^2\right)(t,1)\mathrm{d}t \\&
				+C(\eta)\epsilon\int_{0}^{t}\int_{1}^{+\infty}\left((\partial_{rr} \tilde{S}_1)^2+(\partial_{rr} \tilde{S}_2)^2\right)\mathrm{d}r\mathrm{d}t
				+C\left(E(0)+ E^{\frac{3}{2}}(t)+ E^{\frac{1}{2}}(t) \int_{0}^{t}\mathcal{D}(s)\mathrm{d}s\right)
		\end{aligned}
		\end{equation}
	where \(\tau\epsilon\) is sufficiently small. Thus, the desired result can be obtained immediately.
	\end{proof}
	
	Next, according to Lemmas \ref{lem8},\ref{lem7_tx} and \ref{lem7_xx}, we can obtain the following corollary.
	\begin{corollary}\label{cor1}
		There exists a constant C such that for any $0\leq t \leq T$
		\begin{equation}
			\begin{aligned}
				&\int_{1}^{+\infty}\bigg(\left(r\partial_{rr} \rho+r\partial_{r} \rho\right)^2
				+r^2 (\partial_{rr}v)^2+(\partial_{r}v)^2
				+\tau\left(r^2(\partial_{rr}\tilde{S}_1)^2+(\partial_{r}\tilde{S}_1)^2+\frac{1}{ r^2}(\tilde{S}_1)^2\right)\\
				&\qquad\qquad
				+\tau\left(r\partial_{rr} \tilde{S}_2+2\partial_{r} \tilde{S}_2\right)^2\bigg)\mathrm{d}r
				+\int_{0}^{t}\left((\partial_{rr}v)^2(t,1)+(\partial_{tr}v)^2(t,1)
				\right)\mathrm{d}t\\
				&+\int_0^t\int_{1}^{+\infty}\bigg((r^2(\partial_{rr}\tilde{S}_1)^2+(\partial_{r}\tilde{S}_1)^2+\frac{1}{ r^2}(\tilde{S}_1)^2)
				+\left(r\partial_{rr} \tilde{S}_2+2\partial_{r} \tilde{S}_2\right)^2\bigg)\mathrm{d}r\mathrm{d}t\\
				&\leq C \left(E(0)+E^{\frac{3}{2}}(t)+E^{\frac{1}{2}}(t) \int_{0}^{t}\mathcal{D}(s)\mathrm{d}s\right).	
			\end{aligned}
		\end{equation}
	\end{corollary}
	
	\begin{proof}
		Indeed, multiplying  (\ref{bdy}) by $C+1$ and adding it to (\ref{sed-rr}).
		Choose a sufficiently small $\eta$ such that $\frac{(C+1)}{2\mu}\eta<\frac{1}{2}$,
		and then fix $\eta$, we have 
		\begin{equation}
			\begin{aligned}
				&\int_{1}^{+\infty}\bigg(\left(r\partial_{rr} \rho+r\partial_{r} \rho\right)^2
				+\frac{1}{2}r^2 (\partial_{rr}v)^2+(\partial_{r}v)^2
				+\tau\left(r^2(\partial_{rr}\tilde{S}_1)^2+(\partial_{r}\tilde{S}_1)^2+\frac{1}{ r^2}(\tilde{S}_1)^2\right)\\
				&\qquad\qquad
				+\tau\left(r\partial_{rr} \tilde{S}_2+2\partial_{r} \tilde{S}_2\right)^2\bigg)\mathrm{d}r
				+\int_{0}^{t}\left((\partial_{rr}v)^2(t,1)+(\partial_{tr}v)^2(t,1)
				\right)\mathrm{d}t\\
				&+\int_0^t\int_{1}^{+\infty}\bigg((\frac{1}{2}r^2(\partial_{rr}\tilde{S}_1)^2+(\partial_{r}\tilde{S}_1)^2+\frac{1}{ r^2}(\tilde{S}_1)^2)
				+\frac{1}{2}\left(r\partial_{rr} \tilde{S}_2+2\partial_{r} \tilde{S}_2\right)^2\bigg)\mathrm{d}r\mathrm{d}t\\
				\leq &
				\frac{\tau C}{\mu}\epsilon\int_{1}^{+\infty}r^2(\partial_{rr}v)^2\mathrm{d}r
				+C(\eta)\epsilon\int_{0}^{t}\int_{1}^{+\infty}r^2((\partial_{rr}\tilde{S}_1)^2+(\partial_{rr}\tilde{S}_2)^2)\mathrm{d}r\mathrm{d}t\\
				&+C \left(E(0)+E^{\frac{3}{2}}(t)+E^{\frac{1}{2}}(t) \int_{0}^{t}\mathcal{D}(s)\mathrm{d}s\right).	
			\end{aligned}
		\end{equation}
		For a fixed $\eta$, by choosing a sufficiently small $\epsilon$, and applying Lemma \ref{lem7_tx}, we can obtain the desired result.
	\end{proof}
	On the other hand, using $(\ref{22})_2$, Lemma \ref{lem8} and Corollary \ref{cor1}, we can easily get the second order
	dissipation of $\rho$.
	\begin{corollary}
		There exists some constant C such that
		\begin{equation}
			\int_{0}^{t}\int_{1}^{+\infty}r^2(\partial_{rr}\rho)^2\mathrm{d}r\mathrm{d}t
			\leq C \left(E(0)+E^{\frac{3}{2}}(t)+E^{\frac{1}{2}}(t) \int_{0}^{t}\mathcal{D}(s)\mathrm{d}s\right).
		\end{equation}
	\end{corollary}
	Combining Lemmas Lemma \ref{lem3}-\ref{lem8} and Corollary \ref{cor1},  we get
	\begin{lemma}\label{lem10}
		There exists some constant C such that
		\begin{equation}
			\begin{aligned}
				&\left\|rD(\rho, v, \sqrt{\tau} \tilde{S}_1, \sqrt{\tau} \tilde{S}_2)\right\|_{H^{1}}^2
				+\tau^2\left\|r\partial_{t}^2(\rho, v, \sqrt{\tau} \tilde{S}_1, \sqrt{\tau} \tilde{S}_2)\right\|_{L^{2}}^2\\
				&+\int_{0}^{t}\bigg(\left\|r D^2(\rho, v)\right\|_{L^2}^2
				+\left\|rD(\tilde{S}_1, \tilde{S}_2)\right\|_{H^{1}}^2
				+\tau^2\left\|r\partial_{t}^2(\tilde{S}_1,\tilde{S}_2)\right\|_{L^{2}}^2\bigg)\mathrm{~d} t\\&
				\leq C \left(E(0)+E^{\frac{3}{2}}(t)+E^{\frac{1}{2}}(t) \int_{0}^{t}\mathcal{D}(s)\mathrm{d}s\right).
			\end{aligned}
		\end{equation}
	\end{lemma}
	Therefore, combining Lemma \ref{lem3.5} and \ref{lem10}, the Proposition 3.1 follows immediately.
	
	\section{Passing to the Limit and proof of Main Theorems}
	
	[\textbf{Proof of Theorem \ref{thm1}}]
	According to Proposition \ref{prop1}, the local solution $(\rho^\epsilon,v^\epsilon,\tilde{S}_1^{\epsilon},\tilde{S}_2^{\epsilon})$ can
	be extend to $[0,+\infty)$ by usual continuation methods. Thus, for fixed $\epsilon$, we get a global solution
	$(\rho^\epsilon,v^\epsilon,\tilde{S}_1^{\epsilon},\tilde{S}_2^{\epsilon})$ to system (\ref{9})-(\ref{11}) satisfying
	\begin{equation}\label{4.1}
		\begin{aligned}
			&\sup _{0 \leq s \leq t}\bigg( \sum_{k=0}^1\left\|r\partial_{t}^k(\rho-1, v, \sqrt{\tau} \tilde{S}_1, \sqrt{\tau} \tilde{S}_2)\right\|_{H^{2-k}}^2
			+\tau^2\left\|r\partial_{t}^2(\rho-1, v, \sqrt{\tau} \tilde{S}_1, \sqrt{\tau} \tilde{S}_2)\right\|_{L^{2}}^2\bigg)\\
			&+\int_{1}^{+\infty}\bigg(\sum_{\alpha=1}^2\left\|r D^\alpha(\rho, v)\right\|_{L^2}^2
			+\sum_{k=0}^1\left\|r\partial_{t}^k(\tilde{S}_1, \tilde{S}_2)\right\|_{H^{2-k}}^2
			+\tau^2\left\|r\partial_{t}^2(\tilde{S}_1,\tilde{S}_2)\right\|_{L^{2}}^2\bigg)\mathrm{~d} t
			\leq C E(0)
		\end{aligned}
	\end{equation}
	where C is a constant independent of $\tau$ and $\epsilon$. Therefore, for fixed $\tau$, the uniform bounds of $(\rho^\epsilon-1,v^\epsilon, \sqrt{\tau}\tilde{S}_1^{\epsilon}, \sqrt{\tau}\tilde{S}_2^{\epsilon})$ in $L^{\infty}\left([0,\infty), H^2(\Omega)\right)$
	implies that there exists $(\rho^\epsilon-1,v^\epsilon, \sqrt{\tau}\tilde{S}_1^{\epsilon}, \sqrt{\tau}\tilde{S}_2^{\epsilon})
	\in  L^{\infty}\left([0,\infty), H^2(\Omega)\right)$ such that
	
	$$
	(\rho^\epsilon,v^\epsilon, \sqrt{\tau}\tilde{S}_1^{\epsilon}, \sqrt{\tau}\tilde{S}_2^{\epsilon}) \rightharpoonup
	(\rho,v, \sqrt{\tau}\tilde{S}_1, \sqrt{\tau}\tilde{S}_2) \quad \text { weakly - } * \quad \text { in } \quad L^{\infty}\left(\mathbb{R}^{+} ; H^2(\Omega)\right).
	$$
	
	Further, the uniform boundedness of $\partial_{t}^k(\rho-1, v, \sqrt{\tau} \tilde{S}_1, \sqrt{\tau} \tilde{S}_2)$, $k=1,2$ in
	$L^{\infty}\left([0, T], H^{2-k}\right) \cap L^2\left([0, T], H^{2-k}\right)$ for any $T>0$ ensures, via established compactness theorems, that the sequence $(\rho^\epsilon,v^\epsilon, \sqrt{\tau}\tilde{S}_1^{\epsilon}, \sqrt{\tau}\tilde{S}_2^{\epsilon})$ is relative compact in
	$C^0\left([0, T], H^{2-\delta_0}\right) \cap C^1\left([0, T], H^{1-\delta_0}\right)$ for any $\delta_0>0$. 
	As a consequence, as $\delta\rightarrow 0$ and up to subsequences,
	$$
	(\rho^\epsilon,v^\epsilon, \tilde{S}_1^{\epsilon}, \tilde{S}_2^{\epsilon}) \rightarrow(\rho,v, \tilde{S}_1, \tilde{S}_2), \quad \text { strongly } \quad \text { in } \quad C^0\left([0, T], H^{2-\delta_0}\right) \cap C^1\left([0, T], H^{1-\delta_0}\right).
	$$
	
	This strong convergence suffices to pass to the limit in (\ref{9}) and (\ref{4.1}), thereby confirming that 
	$(\rho,v, \tilde{S}_1, \tilde{S}_2)$ is a classical solution to system $(\ref{1.4})-(\ref{1.6})$ satisfying (\ref{1.8}). The uniqueness follows directly from the regularity
	$$
	(\rho, v, \tilde{S}_1, \tilde{S}_2) \in C\left([0, \infty), H^{2-\delta_0}(\Omega)\right) \cap L^{\infty}\left((0, \infty), H^2(\Omega)\right) \cap L^2\left((0, \infty), H^1(\Omega)\right),
	$$
	which completes the proof of Theorem \ref{thm1}.
	
	
	
	[\textbf{Proof of Theorem \ref{thm2}}]
	Let $(\rho^{\tau},v^{\tau},\tilde{S}_1^{\tau},\tilde{S}_2^{\tau})$ be the global solutions obtained in Theorem \ref{thm1} satisfying (\ref{1.8}). 
	
	Then there exist limit functions
	$
	(\rho^{0},v^{0}) \in L^{\infty}\left(\mathbb{R}^{+} ; H^2(\Omega)\right), (\tilde{S}_1^{0},\tilde{S}_2^{0}) \in L^2\left(\mathbb{R}^{+} ; H^2(\Omega)\right),
	$
	such that along a subsequence as $\tau\rightarrow 0$:
	\begin{equation}
		\begin{aligned}
			& \left(\rho^{\tau},v^{\tau}\right) \rightharpoonup^*\left(\rho^{0},v^{0}\right) \quad \text { weakly-* in } L^{\infty}\left([0,\infty), H^2(\Omega)\right), \\
			& \left(\tilde{S}_1^{\tau},\tilde{S}_2^{\tau}\right) \rightharpoonup\left(\tilde{S}_1^{0},\tilde{S}_2^{0}\right) \quad \text { weakly in } L^2\left([0,\infty), H^2(\Omega)\right) .
		\end{aligned}
	\end{equation}
	The uniform boundedness of $\partial_{t}^k(\rho^{\tau},v^{\tau})$ in $L^2\left([0,\infty), H^{2-k}(\Omega)\right)$ for $k=1,2$ ensures, via classical compactness arguments, that the sequence $(\rho^{\tau},v^{\tau})$ is relatively compact in
	$C^0\left([0, T], H^{2-\delta_0}\right) \cap C^1\left([0, T], H^{1-\delta_0}\right)$ for any $\delta_0>0$ and $T>0$. Consequently, along further subsequences as $\tau$$\rightarrow$0, 
	$$
	(\rho^\tau,v^\tau) \rightarrow(\rho^0,v^0), \quad \text { strongly } \quad \text { in } \quad C^0\left([0, T], H^{2-\delta_0}(\Omega)\right) \cap C^1\left([0, T], H^{1-\delta_0}(\Omega)\right).
	$$
	
	From the uniform estimate (\ref{1.8}) and the constitutive relations $(\ref{1.4})_4-(\ref{1.4})_5$, we derive  
	$$
	\sup_{t \geq 0} \|(\tilde{S}_1^\tau, \tilde{S}_2^\tau)\|_{H^1} \leq C E(0).
	$$
	
	Since $\partial_{t}(\tilde{S}_1^{\tau},\tilde{S}_2^{\tau})$ are bounded in $L^2([0,T], H^1(\Omega))$,
	the sequence $(\tilde{S}_1^{\tau},\tilde{S}_2^{\tau})$ is relatively compact in $C^0\left([0, T], H^{1-\delta}(\Omega)\right)$
	for any $T > 0$. Therefore, as $\tau$$\rightarrow$0 and up to subsequences,
	$$
	\left(\tilde{S}_1^{\tau},\tilde{S}_2^{\tau}\right) \rightarrow\left(\tilde{S}_1^{0},\tilde{S}_2^{0}\right) \quad \text { strongly in } C^0\left([0, T], H^{1-\delta}(\Omega)\right).
	$$
	
	The uniform bound on $\sqrt{\tau}(\tilde{S}_1^{\tau},\tilde{S}_2^{\tau})$ implies:
	\[
	\tau(\tilde{S}_1^{\tau},\tilde{S}_2^{\tau}) \to (0, 0) \quad \text{in } L^\infty([0, \infty), H^2(\Omega)),
	\]
	and consequently in the distributional sense:
	\[
	\tau (\partial_t\tilde{S}_1^\tau, \partial_t \tilde{S}_2^\tau) \to (0, 0) \quad \text{in } D'((0, \infty) \times \Omega)
	\quad \text{as} \quad \tau\rightarrow0.
	\]
	Passing to the limit in equations $(\ref{1.4})_3$ and $(\ref{1.4})_4$, we have
	\begin{equation}\label{4.3}
		\tilde{S}_1^0=2\mu(\partial_{r}v^0-\frac{v^0}{r}), \quad \tilde{S}_2^0=\lambda(\partial_{r}v^0+\frac{2v^0}{r}). \quad \text{a.e. in } (0, \infty) \times \Omega.
	\end{equation}
	
	
	Finally, passing to the limit in $(\ref{1.4})_1-(\ref{1.4})_2$ and using (\ref{4.3}), the limit functions \((\rho^0, v^0)\) satisfy the classical compressible Navier-Stokes system. This completes the proof of Theorem \ref{thm2}.
	
	
	
	
	
	%\bibliography{ref}
	\begin{thebibliography}{10}
		
		%1
%		\bibitem{beirao1989long}
%		H. Beirão da Veiga. 
%		\newblock Long time behavior for one-dimensional motion of a general barotropic viscous fluid. 
%		\newblock {\em Archive for Rational Mechanics and Analysis}, 108:141-160, 1989.
		
		\bibitem{chakraborty2015constitutive}
		D. Chakraborty and J.E. Sader.
		\newblock Constitutive models for linear compressible viscoelastic flows of simple liquids at nanometer length scales.
		\newblock {\em Physics of Fluids}, 27(5): 052002, 2015.
		
\bibitem{Chen2007} S.X. Chen, Initial boundary value problems for quasilinear symmetric hyperbolic systems with characteristic boundary, {\it Front. Math. China}{\bf 2}(1) (2007), 87-102. Translated from {\it 
Chinese Annals of Mathematics} Ser. A {\bf 3}(2) (1982), 223-232. 


		\bibitem{Cho-Kim}
		H.J. Choe and H. Kim. 
		\newblock Global existence of the radially symmetric solutions of the Navier–Stokes equations for the isentropic compressible fluids. 
		\newblock {\em Mathematical Methods in the Applied Sciences}, 28(1): 1-28, 2005.
		
		% %2
		% \bibitem{chen1997remarks}
		% G.Q. Chen.
		% \newblock Remarks on spherically symmetric solutions of the compressible euler
		% equations.
		% \newblock {\em Proceedings of the Royal Society of Edinburgh Section A: Mathematics}, 127(2): 243-259, 1997.
		
		% %3
		% \bibitem{chen2015vanishing}
		% G.Q. Chen and M. Perepelitsa.
		% \newblock Vanishing viscosity solutions of the compressible Euler equations
		% with spherical symmetry and large initial data.
		% \newblock {\em Communications in Mathematical Physics}, 338: 771-800, 2015.
		
		% %4
		% \bibitem{chen2018vanishing}
		% G.Q. Chen and Matthew~R.I. Schrecker.
		% \newblock Vanishing viscosity approach to the compressible Euler equations for
		% transonic nozzle and spherically symmetric flows.
		% \newblock {\em Archive for Rational Mechanics and Analysis}, 229: 1239-1279, 2018.
		
		\bibitem{ding2012global}
		S. Ding, H. Wen, L. Yao and C. Zhu. 
		\newblock Global spherically symmetric classical solution to compressible Navier–Stokes equations with large initial data and vacuum. 
		\newblock {\em SIAM Journal on Mathematical Analysis}, 44(2):1257-1278, 2012.
		
		\bibitem{feireisl2004dynamics}
		E. Feireisl. 
		\newblock Dynamics of viscous compressible fluids. Oxford University Press, 2004.
		
		%5
		\bibitem{freistuhler2020galilei}
		H. Freistuhler.
		\newblock A Galilei invariant version of Yong's model.
		\newblock {\em preprint, arXiv:2012.09059}, 2020.
		
		%6
		\bibitem{hoff1992spherically}
		D. Hoff.
		\newblock Spherically symmetric solutions of the Navier-Stokes equations for
		compressible, isothermal flow with large, discontinuous initial data.
		\newblock {\em Indiana University Mathematics Journal}, 41(4): 1225-1302, 1992.
		
		\bibitem{HL24}
		Y. Hu and Y.C. Li.
		\newblock Initial boundary value problem for relaxed isentropic compressible Navier-Stokes equations.
		\newblock {\em preprint}, 2024.
		
		\bibitem{HWAML}
		Y. Hu and N. Wang.
		\newblock Blowup of solutions for compressible Navier-Stokes equations with revised Maxwell’s law. 
		\newblock {\em Applied Mathematics Letters}, 103: 106221, 2020.
		%7
%		\bibitem{hu2016compressible}
%		Y. Hu and R. Racke.
%		\newblock Compressible Navier-Stokes equations with hyperbolic heat
%		conduction.
%		\newblock {\em Journal of Hyperbolic Differential Equations}, 13(02): 233-247, 2016.
		
		%8
		\bibitem{hu2017compressible}
		Y. Hu and R. Racke.
		\newblock Compressible Navier-Stokes equations with revised Maxwell’s law.
		\newblock {\em Journal of Mathematical Fluid Mechanics}, 19: 77-90, 2017.
		
		\bibitem{HZ25}
		Y. Hu and X.N. Zhao.
		\newblock Global well-posedness and relaxation limit for relaxed compressible Navier-Stokes-Fourier equations in bounded domain.
		\newblock {\em preprint, arXiv:2502.00544}, 2025.
		
		%9
		\bibitem{jiang1996global}
		S. Jiang.
		\newblock Global spherically symmetric solutions to the equations of a viscous
		polytropic ideal gas in an exterior domain.
		\newblock {\em Communications in Mathematical Physics}, 178(2): 339-374, 1996.
		
%		%10
%		\bibitem{jiang1998large}
%		S. Jiang.
%		\newblock Large-time behavior of solutions to the equations of a viscous
%		polytropic ideal gas.
%		\newblock {\em Annali di Matematica Pura ed Applicata}, 175(1): 253-275, 1998.
		
		
		\bibitem{jiang2001spherically}
		S. Jiang and P. Zhang.
		\newblock On spherically symmetric solutions of the compressible isentropic Navier--Stokes equations.
		\newblock {\em Communications in Mathematical Physics}, 215: 559-581, 2001.
		
%		\bibitem{kazhikhov1977global}
%		{ A. V. Kazhikhov, V. V. Shelukin, Unique global solution with respect to time of initial boundary value problems for one-dimensional equations of a viscous gas. J. Appl. Math. Mech.
%{bf 41 (2)} (1977) 273-282; traslated from Prikl. Mat. Meh. {\bf 41} (2) (1977) 282-291 (Russian).}
%		
%		\bibitem{kawashima1981global}
%		S. Kawashima and T. Nishia. 
%		\newblock Global solutions to the initial value problem for the equations of one-dimensional motion of viscous polytropic gases. 
%		\newblock {\em Journal of Mathematics of Kyoto University}, 21(4): 825-837, 1981.
		
		%11
		\bibitem{lions1998mathematical}
		P.L. Lions.
		\newblock Mathematical Topics in Fluid Mechanics, Compressible Models, vol. 2. Clarendon. 1998.
		
		% \bibitem{lefloch2007finite}
		% P.G. LeFloch and M. Westdickenberg.
		% \newblock Finite energy solutions to the isentropic euler equations with
		% geometric effects.
		% \newblock {\em Journal de math{\'e}matiques pures et appliqu{\'e}es}, 88(5): 389-429, 2007.
		
%		\bibitem{matsumura1980initial}
%		Matsumura A, Nishida T. 
%		\newblock The initial value problem for the equations of motion of viscous and heat-conductive gases. 
%		\newblock {\em Journal of Mathematics of Kyoto University}, 20(1):67-104, 1980.
		%12
		\bibitem{makino1994initial}
		T. Makino and S. Takeno.
		\newblock Initial boundary value problem for the spherically symmetric motion
		of isentropic gas.
		\newblock {\em Japan Journal of Industrial and Applied Mathematics}, 11: 171-183, 1994.
		
		%13
		\bibitem{maxwell1867iv}
		J.C. Maxwell.
		\newblock On the dynamical theory of gases.
		\newblock {\em Philosophical Transactions of the Royal Society of London},
		157: 49-88, 1967.
		
%		%14
%		\bibitem{nikolaev1983solvability}
%		V.B. Nikolaev.
%		\newblock On the solvability of mixed problem for one-dimensional
%		axisymmetrical viscous gas flow. dinamicheskie zadachi mekhaniki sploshnoj
%		sredy, 63.
%		\newblock {\em Sibirsk. Otd. Acad. Nauk SSSR, Inst. Gidrodinamiki}, 63, 1983.
		
		%15
		\bibitem{pelton2013viscoelastic}
		M. Pelton, D. Chakraborty, E. Malachosky, P. Guyot-Sionnest and J.E. Sader.
		\newblock Viscoelastic flows in simpl liquids generated by vibrating nanostructures.
		\newblock {\em Physical Review Letters}, 111: 244502, 2013.
		
		%16
		\bibitem{peng2021relaxed}
		Y.J. Peng.
		\newblock Relaxed Euler systems and convergence to Navier-Stokes equations.
		\newblock {\em Annales de l'Institut Henri Poincar{\'e} C, Analyse non lin{\'e}aire}, 38(2): 369-401, 2021.
		
		%17
		\bibitem{schochet1986compressible}
		S. Schochet.
		\newblock The compressible Euler equations in a bounded domain: existence of
		solutions and the incompressible limit.
		\newblock {\em Communications in Mathematical Physics}, 104(1): 49-75, 1986.
		
		% %18
		% \bibitem{yang1995functional}
		% T. Yang.
		% \newblock A functional integral approach to shock wave solutions of euler
		% equations with spherical symmetry.
		% \newblock {\em Communications in mathematical physics}, 171: 607-638, 1995.
		
		% %19
		% \bibitem{yang1996functional}
		% T. Yang.
		% \newblock A functional integral approach to shock wave solutions of the euler
		% equations with spherical symmetry (ii).
		% \newblock {\em Journal of differential equations}, 130(1): 162-178, 1996.
		
		%20
		\bibitem{yong2014newtonian}
		W.A. Yong.
		\newblock Newtonian limit of Maxwell fluid flows.
		\newblock {\em Archive for Rational Mechanics and Analysis}, 214: 913-922, 2014.
		
	\end{thebibliography}
	
	
	
	
	
	
	
	
	
	
	
	
	
	
	
	
\end{document} 
